%! TeX program = lualatex
\documentclass[12pt,a4paper,oneside]{report}

\usepackage{amsmath}
\usepackage{mathtools}
\usepackage{amssymb}
\usepackage{amsfonts}
\usepackage{amsthm}
\usepackage{bm}
\usepackage{bbm}
\usepackage{centernot}
\usepackage{enumitem}

\usepackage{fontspec}
\usepackage[a4paper,margin=1.2in]{geometry}
\usepackage[colorlinks=true, allcolors=blue, plainpages=false, pdfpagelabels]{hyperref}
\usepackage{mathrsfs}

\usepackage[backend=biber,style=alphabetic,maxbibnames=99]{biblatex}
\usepackage{stmaryrd}
\SetSymbolFont{stmry}{bold}{U}{stmry}{m}{n}
\usepackage{import}
\usepackage{xifthen}
\usepackage{pdfpages}
\usepackage{booktabs}

\usepackage{graphicx}
\usepackage{tikz}
\usetikzlibrary{calc}
\usetikzlibrary{arrows.meta}
\usetikzlibrary{shapes.geometric}
\usetikzlibrary{positioning}

\usepackage{microtype}

\newcommand{\incfig}[1]{%
    \def\svgwidth{0.5\columnwidth}
    \import{./figures/}{#1.pdf_tex}
}

\makeatletter
\newtheorem*{rep@theorem}{\rep@title}
\newcommand{\newreptheorem}[2]{%
\newenvironment{rep#1}[1]{%
 \def\rep@title{#2 ##1}%
 \begin{rep@theorem}}%
 {\end{rep@theorem}}}
\makeatother

\theoremstyle{plain}
\newtheorem{theorem}{Theorem}[chapter]

\newtheorem{corollary}[theorem]{Corollary}
\newtheorem{claim}[theorem]{Claim}
\newtheorem{lemma}[theorem]{Lemma}
\newtheorem{conjecture}[theorem]{Conjecture}
\newtheorem{problem}[theorem]{Problem}
\newtheorem{proposition}[theorem]{Proposition}
\newtheorem{question}[theorem]{Question}

\newtheorem*{question*}{Question}

\theoremstyle{definition}
\newtheorem{definition}[theorem]{Definition}
\newtheorem{construction}[theorem]{Construction}
\newtheorem{remark}[theorem]{Remark}

% paired solution
\newenvironment{sol}
  {\par\noindent\textbf{Solution}.}
  {\par\bigskip}

\input{preamble/mathdefs}

\addbibresource{preamble/ref.bib}

\newenvironment{comments}
  {\par\medskip\begingroup\color{red!70!black}%
   \noindent\textbf{Comments:}%
   \begin{enumerate}[leftmargin=*,itemsep=0.25em,topsep=0.4em]}
  {\end{enumerate}\endgroup\medskip}

\newcommand{\sharat}[1]{\textcolor{red}{Sharat: #1}} %%% don't remove this macro

\begin{document}

% ============================================================
% FRONT MATTER
% ============================================================

\pagenumbering{Roman}
% ============================================================
% Thesis metadata — fill in your details here
% ============================================================
\newcommand{\thesistitle}{Information Gaps and Algorithms for the Stochastic Traveling Salesperson Problem}
\newcommand{\thesisauthor}{Teng Liu}
\newcommand{\thesisdate}{July 2026}
\newcommand{\thesisadvisor}{Dr. Sharat Ibrahimpur}
\newcommand{\thesissecondadvisor}{Prof.\ Dr.\ Vera Traub}  % leave empty if none
\newcommand{\thesisdepartment}{Department of Computer Science}
\newcommand{\thesisinstitution}{ETH Z\"urich}

% ============================================================
% Title page
% ============================================================
\newcommand{\maketitlepage}{%
\begin{titlepage}
\begin{center}

\vspace*{0.5cm}
{\large\textbf{ETH} Z\"urich}\\[0.15cm]
{\small Eidgen\"ossische Technische Hochschule Z\"urich}\\[0.05cm]
{\small Swiss Federal Institute of Technology Zurich}

\vspace{4cm}

{\LARGE\bfseries \thesistitle}

\vspace{2.5cm}

{\Large Master Thesis}

\vspace{0.5cm}

{\Large \thesisauthor}

\vspace{0.5cm}

{\Large \thesisdate}

\vfill

\ifthenelse{\isempty{\thesissecondadvisor}}{%
{\large Advisor: \thesisadvisor}\par
}{%
{\large Advisors: \thesisadvisor, \thesissecondadvisor}\par
}
\vspace{0.3cm}
{\large \thesisdepartment, \thesisinstitution}\par

\end{center}
\end{titlepage}
}

\hypersetup{
  pdftitle={\thesistitle},
  pdfauthor={\thesisauthor}
}
\maketitlepage
\pagenumbering{roman}

% --- Abstract ---
\chapter*{Abstract}
\addcontentsline{toc}{chapter}{Abstract}

We study the stochastic traveling salesperson problem, where each non-root
vertex is active independently according to a given probability and all active
vertices must be visited at minimum expected travel cost.

The classical a~posteriori and a~priori models represent two information
extremes.  An a~posteriori solution knows the complete active set before
choosing its tour, whereas an a~priori solution commits to a fixed cyclic
ordering before any activations are known.  The gap between these two
benchmarks is constant for symmetric metrics but polynomial in the number of
vertices for asymmetric metrics.  This thesis introduces an adaptive model
between these extremes, in which vertices are probed sequentially and the next
probe may depend on previous outcomes, with the goal of understanding whether
adaptivity produces better stochastic tours.

We establish new separations between these models.  For symmetric metrics, we
construct instances with adaptivity gap approaching~$2$ and clairvoyance gap
approaching~$10/9$; for instances with one semi-active vertex, we also prove the
uniform upper bound $(2+\sqrt6)/4$ on the clairvoyance gap.  For asymmetric
metrics, we give a layered-poset construction showing that the clairvoyance
gap is at least~$3/2$.

On the algorithmic side, for asymmetric metrics we construct an adaptive
policy with expected cost
$O(n^{1/3})\optpost$, proving an $O(n^{1/3})$ upper bound on the
clairvoyance gap.  For every fixed
$\varepsilon>0$, we also give a deterministic
$(1{+}\varepsilon)$-approximation algorithm for a~priori TSP with running time
$c_\varepsilon^n\poly(n)$ based on dynamic programming.

\tableofcontents
\clearpage

% ============================================================
% CHAPTER 1: INTRODUCTION
% ============================================================
\pagenumbering{arabic}
\chapter{Introduction}\label{ch:intro}

Many real-world routing and scheduling problems are subject to uncertainty.
In a delivery service, for example, a courier must plan a route through several customer locations, yet some customers may have no packages to receive or their requests may be canceled.
A classical abstraction is the \emph{stochastic traveling salesperson problem} (\StochTSP{}), where each client~$v$ comes with a specified probability~$p_v$ and becomes \emph{active} (requires service) independently with probability~$p_v$.
The goal is to find a tour that visits all active vertices and minimizes the expected travel cost, where the cost comes from the underlying metric on vertices.

An important modeling choice is whether the salesperson is allowed to adapt future decisions based on information revealed during the execution of the tour.
In the \emph{nonadaptive} (a~priori) setting, the salesperson commits to a fixed tour before observing the realization; during execution the salesperson follows this tour while shortcutting inactive vertices.
In the \emph{adaptive} setting, decisions may depend on observations revealed along the way, allowing the tour to be adjusted on the fly.
In the \emph{a~posteriori} setting, the salesperson has full knowledge of the realization.

A fundamental question is how much adaptivity can improve performance relative to an optimal nonadaptive tour, and how much full clairvoyance improves over the best adaptive policy.
Note that computing an optimal adaptive policy (or an optimal a~priori tour) for \StochTSP{} is computationally hard, since the setting with $p_v=1$ for all~$v$ corresponds to the classical TSP, which is known to be APX-hard~\cite{Karpinski2015}.
Also, describing an optimal adaptive policy might require exponential space.
Thus, we are interested in understanding the separations between the three benchmarks, and in designing policies or algorithms whose expected cost can be provably bounded.

\section{Our contributions}\label{sec:contributions}

We obtain the following results.

\begin{enumerate}
    \item \textbf{Adaptivity gap lower bound for symmetric metrics} (\cref{ch:sym-adapt}).
    We first construct an explicit family of instances based on cycle-subdivided cliques
    for which $\optprior/\optadapt$ tends to~$5/4$
    (\cref{thm:cycle-subdivision}).  Motivated by the random activation of the
    semi-active vertices in this construction, we then transfer the randomness
    to the metric and use the probabilistic method to obtain deterministic
    symmetric instances whose adaptivity ratio tends to~$2$
    (\cref{thm:random-metric-gap}).  Since
    $\optpost\le\optadapt$, the same factor-two lower bound applies to the
    obliviousness gap.

    \item \textbf{Clairvoyance gap for symmetric metrics} (\cref{ch:sym-clair}).
    We give a three-chain family with one semi-active vertex of
    activation probability~$1/2$ whose ratios converge to~$10/9$.  We also
    prove the upper bound
    \[
        \frac{\optadapt}{\optpost}\le\frac{2+p}{2+p^2}
    \]
    for every symmetric instance with one semi-active vertex of activation
    probability~$p$.  The right-hand side is maximized at
    $p=\sqrt6-2$, giving the uniform upper bound
    $(2+\sqrt6)/4$ for this class.

    \item \textbf{Clairvoyance gap lower bound for asymmetric metrics} (\cref{ch:asym-clair}).
    We give a four-selector layered-poset construction for which
    \[
        \liminf\frac{\optadapt}{\optpost}\ge\frac32
    \]
    (\cref{thm:layered-poset-gap}).  The construction uses independent
    activations and a strictly positive asymmetric metric.  Before the main
    construction, we record a natural recursive-triangle family whose ratio
    tends to~$4/3$ (\cref{rem:recursive-triangle}); its longer analysis is
    kept in a separate technical note.  We also give an explicit finite
    member whose ratio is strictly larger than~$4/3$.

    \item \textbf{Clairvoyance-gap upper bound for asymmetric metrics}
    (\cref{ch:circulation}).
    We prove that every asymmetric instance admits an adaptive policy
    of expected cost $O(n^{1/3})\optpost$
    (\cref{thm:asymmetric-clairvoyance-upper}).  The proof starts from a
    low-cost backbone circulation.  After assigning each outside client a
    public directed backbone chord, we round an auxiliary fractional flow to
    an integer flow whose mean imbalance is at most one on every cyclic
    interval.  A maximal inequality controls the remaining Bernoulli
    imbalance, and an active-first stack walk turns the rounded object into a
    adaptive policy.  Splitting clients at probability $n^{-1/3}$ gives the
    stated bound.

    \item \textbf{Approximate a~priori TSP} (\cref{ch:dp}).
    We give a $(1{+}\varepsilon)$-approximation algorithm for a~priori TSP on
    both symmetric and asymmetric metrics, running in time
    $c_\varepsilon^n\poly(n)$ (\cref{thm:apriori-ptas}).
    The algorithm discards pairs for which it is very unlikely that every
    vertex lying strictly between their endpoints is inactive, losing only a
    negligible amount, and then uses a dynamic program whose state stores just
    the relevant suffix of the partial order.
\end{enumerate}

\section{Related work}\label{sec:related}

\paragraph{Classical TSP.}
The traveling salesperson problem is one of the most studied problems in combinatorial optimization.
For symmetric metrics, the classical $3/2$-approximation of Christofides~\cite{Christofides2022} stood for decades until Karlin, Klein, and Oveis~Gharan~\cite{Karlin2021} obtained a slightly improved factor.
For asymmetric metrics, Svensson, Tarnawski, and V\'egh~\cite{Svensson2020} gave the first constant-factor approximation.

\paragraph{A~priori TSP.}
The a~priori model was introduced by Jaillet~\cite{Jaillet1985}.
Shmoys and Talwar~\cite{Shmoys2008} proved that the obliviousness gap $\optprior/\optpost$ is at most~$3$ for symmetric metrics, yielding a constant-factor approximation via a master tour.
Schalekamp and Shmoys~\cite{SchalekampShmoys2008} studied the connection between universal and a~priori TSP.
Blauth, Neuwohner, Puhlmann, and Vygen~\cite{Blauth2024} recently obtained improved approximation guarantees for a~priori TSP.
For asymmetric metrics, Christalla, Puhlmann, and Traub~\cite{Christalla2025} showed that the obliviousness gap is $O(\sqrt{n})$ and $\Omega(n^{1/4}/\log n)$; moreover, their construction also implies an adaptivity gap of $\Omega(n^{1/4}/\log n)$.

\paragraph{Adaptive and \StochTSP{}.}
Gupta, Nagarajan, and Ravi~\cite{GuptaNagarajanRavi2017} studied adaptive TSP and gave an $O(\log^2 n)$-approximation for the stochastic version.
Ene, Nagarajan, and Saket~\cite{EneNagarajanSaket2017} studied approximation algorithms for stochastic $k$-TSP.
Jiang, Li, Liu, and Singla~\cite{JiangLiLiuSingla2020} obtained improved algorithms and adaptivity gaps for stochastic $k$-TSP.
Fu, Li, and Xu~\cite{FLX18} gave a PTAS for a class of stochastic dynamic programs.

\paragraph{Adaptivity gaps in stochastic optimization.}
The study of adaptivity gaps---how much adaptivity helps in stochastic optimization---has a rich history beyond TSP.
Dean, Goemans, and Vondr\'ak~\cite{DeanGoemansVondrak2008} studied the benefit of adaptivity for stochastic knapsack.
Gupta and Nagarajan~\cite{GuptaNagarajan2013} introduced the stochastic probing framework and proved constant adaptivity gaps for matching.
Gupta, Nagarajan, and Singla~\cite{GuptaNagarajanSingla2016,GuptaNagarajanSingla2017} extended these results to submodular and XOS objectives.
Brada\v{c}, Singla, and \v{Z}u\v{z}i\'c~\cite{BradacSinglaZuzic2019} obtained near-optimal adaptivity gaps for stochastic multi-value probing.
Levin and Vainer~\cite{LV14} studied adaptivity in the stochastic blackjack knapsack problem.
Barak and Talgam-Cohen~\cite{BarakTalgamCohen2026} recently studied semi-adaptivity gaps for stochastic knapsack.

% ============================================================
% MODEL AND BENCHMARKS
% ============================================================
\section{Model and benchmarks}\label{sec:model}

This section fixes the notation used throughout the thesis.  We first define
the underlying metric space and its tours, then formalize the three stochastic
benchmarks, and finally introduce the ratios that compare their optimal costs.

\subsection{Metric spaces and tours}\label{sec:metric-spaces-tours}

Let $V$ be a finite vertex set of cardinality $n+1$, fix a distinguished
\emph{root} $r\in V$, and write
\[
    V^\circ:=V\setminus\{r\}
\]
for the $n$ non-root vertices.  A \emph{possibly asymmetric metric} on~$V$ is a function
$d:V\times V\to\mathbb R_{\ge0}$ such that $d(x,x)=0$, $d(x,y)>0$ for
distinct $x,y\in V$, and
\[
    d(x,z)\le d(x,y)+d(y,z)
    \qquad\forall x,y,z\in V.
\]
A metric~$d$ is \emph{symmetric} if $d(x,y)=d(y,x)$ for all $x,y\in V$,
and \emph{asymmetric} otherwise.  Unless explicitly stated otherwise, metrics
are allowed to be asymmetric.

A \emph{tour} on a subset $S\subseteq V$ containing the root is a cyclic
sequence $T=(r,u_1,\ldots,u_k,r)$, where $k=|S|-1$ and
$(u_1,\ldots,u_k)$ is an ordering of $S\setminus\{r\}$.  If $k=0$, the
root-only tour $T=(r)$ has cost~$0$.  If $k\ge1$, its \emph{cost} is
\[
    \cost(T)
    :=d(r,u_1)+\sum_{i=1}^{k-1}d(u_i,u_{i+1})+d(u_k,r).
\]
For a subset $A\subseteq V^\circ$, we write $\opt(A)$ for the minimum cost
of a tour on $A\cup\{r\}$.

\subsection{Three benchmarks for \StochTSP{}}\label{sec:stoch-tsp-benchmarks}

In the stochastic setting, the root is always active, so $p_r=1$.  The set of
active non-root vertices is a subset $A\subseteq V^\circ$ generated by
independent activations: each $v\in V^\circ$ is active (that is, $v\in A$)
independently with probability~$p_v$. We denote by~$\mathcal Q$ the resulting product distribution
on~$2^{V^\circ}$, so
\[
    \mathcal Q(A)
    =\prod_{v\in A}p_v
     \prod_{v\in V^\circ\setminus A}(1-p_v).
\]

We consider three benchmarks for \StochTSP{}, distinguished by how much the solution may depend on the realization of the active set. In each case we define the cost of a feasible solution via an overloaded function $\cost(\cdot)$.

\begin{enumerate}
    \item \textbf{A posteriori TSP.}
        An \emph{a~posteriori solution} is a family
        $\mathcal{S}=(T_A)_{A\subseteq V^\circ}$, where $T_A$ is an optimal
        tour on $A\cup\{r\}$ and therefore has cost~$\opt(A)$.
        We define
        \[
        \optpost
        :=\cost(\mathcal{S})
        :=\E_{A\sim\mathcal Q}\bigl[\cost(T_A)\bigr]
        =\E_{A\sim\mathcal Q}\bigl[\opt(A)\bigr].
        \]

    \item \textbf{A priori (nonadaptive) TSP.}
        An \emph{a~priori tour} is a cyclic ordering of all vertices,
        represented as $T=(r,v_1,\ldots,v_n,r)$, where
        $V^\circ=\{v_1,\ldots,v_n\}$.
        For a realization~$A$, the \emph{shortcut} $T\vert_A$ is obtained by retaining only the vertices in $A\cup\{r\}$ in the cyclic order prescribed by~$T$.
        If $A=\emptyset$, set $T\vert_A=(r)$.  Otherwise, if
        $1 \leq i_1 < i_2 < \cdots < i_{|A|} \leq n$ are such that
        $\{v_{i_k} : 1 \leq k \leq |A|\} = A$, then
        $T\vert_A = (r,v_{i_1},v_{i_2},\ldots,v_{i_{|A|}},r)$.
        We overload:
        \[
        \cost(T) := \E_{A\sim\mathcal Q}\bigl[\cost(T\vert_A)\bigr].
        \]
        The optimum is $\optprior := \min_{T}\,\cost(T)$, where the minimum is over all a~priori tours.

    \item \textbf{Adaptive TSP.}
        The salesperson knows all the activation probabilities but not the realization~$A$.
        During execution, the salesperson \emph{probes} non-root vertices one
        by one. A probe is performed remotely and reveals whether the probed
        vertex is active before the salesperson moves.  If the probed vertex~$v$ lies
        in~$A$, the salesperson \emph{must} travel to~$v$ immediately;
        otherwise the salesperson remains in place.  The order of probes may
        depend on the outcomes observed so far.

        Define the set of valid states by
        \[
            \mathcal D
            :=\{(u,S)\in V\times2^{V^\circ}:u\notin S\}.
        \]
        Formally, an adaptive policy is a function
        $\pi:\mathcal D\rightarrow V$, where
        $\pi(u,S)$ denotes the next vertex to probe when the salesperson is
        at~$u$ and the set of unprobed vertices is~$S$.
        We require $\pi(u,\emptyset)=r$, corresponding to the final return to
        the root, and $\pi(u,S)\in S$ whenever $S\ne\emptyset$.

        Restricting the policy to the state $(u,S)$ is without loss of
        generality.  Indeed, independence implies that, conditional on the
        probe outcomes observed so far, the vertices in~$S$ retain their original
        independent activation probabilities.  If $J(u,S)$ denotes the
        minimum expected remaining cost, then backward induction gives
        \[
            J(u,\emptyset)=d(u,r)
        \]
        and, for $S\neq\emptyset$,
        \[
            J(u,S)
            =
            \min_{v\in S}
            \left\{
                (1-p_v)J(u,S\setminus\{v\})
                +p_v\bigl(d(u,v)+J(v,S\setminus\{v\})\bigr)
            \right\}.
        \]
        Thus a deterministic Markov policy attains the optimum; allowing
        dependence on the full history or private randomization does not
        decrease the optimal expected cost.

        Given a realization~$A$, the policy induces a realized tour $T(A;\pi)$ on $A\cup\{r\}$ as follows.
        Initialize $u \leftarrow r$, $S \leftarrow V^\circ$.
        While $S \neq \emptyset$: compute $v = \pi(u,S)$ and probe~$v$; if $v \in A$ then travel to~$v$ and set $u\leftarrow v$; set $S \leftarrow S \setminus \{v\}$.
        Finally, return to~$r$.
        We overload:
        \[
        \cost(\pi) := \E_{A\sim\mathcal Q}\bigl[\cost(T(A;\pi))\bigr].
        \]
        The optimum is $\optadapt := \min_{\pi}\,\cost(\pi)$, where the minimum is over all adaptive policies.
\end{enumerate}

Throughout the thesis, we reserve $\pi$ for adaptive policies and use $T$ for
tours, including fixed a~priori orders.  We use $p$ (or $p_v$) for activation
probabilities and $q$ (or $q_v$) for the corresponding inactivity
probabilities.

\subsection{Gap measures}\label{sec:gaps}

Among the three benchmarks, a~posteriori TSP has the most information (the full realization~$A$), a~priori TSP has the least (only the metric and probabilities), and adaptive TSP lies in between.

We define an \emph{instance} of \StochTSP{} by
\[
    \mathcal{I}:=(V,r,d,(p_v)_{v\in V}),
    \qquad p_r=1.
\]
For any instance~$\mathcal{I}$,
\begin{equation}\label{eq:optineq}
\optpost \le \optadapt \le \optprior
\end{equation}
since from left to right the solution has less information about the realization, resulting in higher optimal cost.

We define the following gap measures:
\begin{itemize}
    \item \emph{Adaptivity gap}: $\sup_{\mathcal{I}} \f{\optprior}{\optadapt}$ --- the worst-case benefit of adaptivity over a fixed a~priori tour.
    \item \emph{Clairvoyance gap}: $\sup_{\mathcal{I}} \f{\optadapt}{\optpost}$ --- the benefit of full clairvoyance (knowing $A$ upfront) over the best adaptive policy.
    \item \emph{Obliviousness gap}: $\sup_{\mathcal{I}} \f{\optprior}{\optpost}$ --- combines both effects.
\end{itemize}
Whenever one of these ratios has both numerator and denominator equal to zero,
we define its value to be~$0$.  Under the strict positivity assumption on the
metric, a zero denominator in any of the three ratios forces the corresponding
numerator to be zero as well.
We remark that the literature on a~priori TSP uses the term adaptivity gap for what we call the obliviousness gap; we deviate from this convention following \cite{Christalla2025} for the sake of clarity.

For symmetric metrics, the obliviousness gap is known to be at most $O(1)$; in particular, $\optprior/\optpost \le 3$ is shown in~\cite{Shmoys2008}.
The same upper bound carries over to the clairvoyance gap and the adaptivity gap by~\cref{eq:optineq}.
For asymmetric metrics, Christalla, Puhlmann, and Traub~\cite{Christalla2025} show that the obliviousness gap is $O(\sqrt{n})$ and $\Omega(n^{1/4} \cdot \log^{-1}n)$.
Moreover, the construction in Section~3.1 of~\cite{Christalla2025} also implies an adaptivity gap of $\Omega(n^{1/4} \cdot \log^{-1}n)$.

\sharat{This subsection should be moved to prelims chapter. It's OK for a chapter to be short.}

\subsection{Probabilistic inequalities}\label{sec:probabilistic-inequalities}

We use the following weighted form of Hoeffding's inequality.  Let
$Y_1,\ldots,Y_m$ be independent random variables with
$Y_i\in[a_i,b_i]$ almost surely.  Then, for every $t>0$,
\begin{equation}\label{eq:weighted-hoeffding}
    \P\left[\sum_{i=1}^m(Y_i-\E[Y_i])\ge t\right]
    \le
    \exp\left(-\frac{2t^2}{\sum_{i=1}^m(b_i-a_i)^2}\right).
\end{equation}
In particular, if the $X_i$ are independent Bernoulli random variables and
$c_i\ge0$, then~\cref{eq:weighted-hoeffding} applies to
$Y_i=c_iX_i\in[0,c_i]$.  We also use the union bound and Markov's inequality
in their standard forms.

\section{Organization}\label{sec:organization}

\Cref{ch:sym-adapt} proves an adaptivity-gap lower bound of~$2$ for symmetric metrics, beginning with an explicit construction based on cycle-subdivided cliques and then strengthening it through a probabilistic construction.
\Cref{ch:sym-clair} proves a symmetric clairvoyance-gap lower bound
of~$10/9$ and establishes a probability-dependent upper bound for
symmetric instances with one semi-active vertex.
\Cref{ch:asym-clair} motivates the asymmetric obstruction through a
recursive-triangle family with ratio tending to~$4/3$, then proves the
stronger lower bound~$3/2$ using a layered-poset construction.
\Cref{ch:circulation} develops the backbone-circulation and integer-flow
rounding framework and uses it to prove an $O(n^{1/3})$ upper bound on the
asymmetric clairvoyance gap.
\Cref{ch:dp} presents the $(1{+}\varepsilon)$-approximation algorithm for a~priori TSP.
\Cref{ch:conclusion} concludes with a summary and directions for future
research.

% ============================================================
% CHAPTER 2: ADAPTIVITY GAP — SYMMETRIC
% ============================================================
\chapter{Adaptivity Gap for Symmetric \StochTSP{}}\label{ch:sym-adapt}

In this chapter we focus on symmetric \StochTSP{}.
The main result of this chapter is that the adaptivity gap is at least~$2$.
We give two qualitatively different constructions.
In \Cref{sec:cycle-subdivided-construction}, we establish the
weaker lower bound of~$5/4$ through an explicit family of instances
whose metrics arise from cycle-subdivided cliques.  In
\Cref{sec:random-metric-gap}, we strengthen the lower bound to~$2$ using
the probabilistic method on $(1,2)$-TSP metrics induced by a random graph~$G_{n,p}$
in the Erdős--Rényi model.

\section{The cycle-subdivided clique construction}
\label{sec:cycle-subdivided-construction}

This section gives an explicit family of instances for which the ratio between
the a~priori and adaptive optima tends to~$5/4$.
The construction begins with a complete graph, selects a
Hamiltonian cycle, and inserts one vertex of activation probability~$1/2$ on
each edge of that cycle.

\begin{definition}[Cycle-subdivided clique]\label{def:cycle-sub-graph}
Fix $n\ge 7$, and let $b_1,\dots,b_n$ be the vertices of the complete graph
$K_n$.  Throughout this construction, vertex indices are interpreted
modulo~$n$ with representatives in~$[n]$; in particular, $n+1$ is identified
with~$1$.  We call
$C_n=(b_1,\dots,b_n,b_1)$ the \emph{base cycle}.  For every $i\in[n]$,
subdivide the edge $b_ib_{i+1}$ of the base cycle by introducing a new
vertex~$s_i$.  Denote the resulting unweighted graph
by~$\Sub(K_n;C_n)$.
\end{definition}
The instance $\mathcal{I}_n^{\mathrm{CSC}}$ based on the cycle-subdivided clique
has root $r:=b_1$ and vertex set
\[
V_n:=\{b_1,\dots,b_n\}\cup\{s_1,\dots,s_n\}.
\]
The underlying distance function is the shortest-path metric induced
by~$\Sub(K_n;C_n)$.  The vertices $b_1,\dots,b_n$ are the
\emph{always-active} vertices: the root~$b_1$ is always active by definition,
and every non-root vertex $b_2,\dots,b_n$ is active with probability~$1$.
The vertices $s_1,\dots,s_n$ are the \emph{semi-active} vertices, each active
independently with probability~$1/2$.

\begin{theorem}\label{thm:cycle-subdivision}
The adaptivity gap is at least $\frac54$.
\end{theorem}

\begin{figure}[htbp]
\centering
\begin{tikzpicture}[
    scale=1.05,
    alwaysactive/.style={rectangle,draw=black,fill=white,minimum size=8mm,inner sep=0pt,font=\small},
    semiactive/.style={circle,draw=black,fill=gray!20,minimum size=7mm,inner sep=0pt,font=\small},
    cycleedge/.style={black,line width=0.6pt},
    chord/.style={black,line width=0.6pt}
]
\foreach \i in {1,...,6} {
  \coordinate (c\i) at ({90-60*(\i-1)}:3.2);
}
\foreach \i in {1,...,6} {
  \pgfmathtruncatemacro{\next}{mod(\i,6)+1}
  \coordinate (m\i) at ($(c\i)!0.5!(c\next)$);
}
\foreach \i in {1,...,6} {
  \foreach \j in {1,...,6} {
    \ifnum\j>\i
      \pgfmathtruncatemacro{\next}{mod(\i,6)+1}
      \ifnum\j=\next\else
        \ifnum\i=1\ifnum\j=6\else\draw[chord] (c\i)--(c\j);\fi
        \else\draw[chord] (c\i)--(c\j);\fi
      \fi
    \fi
  }
}
\foreach \i in {1,...,6} {
  \pgfmathtruncatemacro{\next}{mod(\i,6)+1}
  \draw[cycleedge] (c\i)--(m\i)--(c\next);
}
\foreach \i in {1,...,6} {
  \node[alwaysactive] (b\i) at (c\i) {$b_\i$};
  \node[semiactive] (s\i) at (m\i) {$s_\i$};
}
\node[above=2mm of b1,font=\small] {root $r$};
\end{tikzpicture}
\caption{The cycle-subdivided clique $\Sub(K_6;C_6)$. The always-active
vertices are represented as squares, with the root indicated above; the
semi-active vertices are represented as gray circles.}
\label{fig:cycle-subdivision}
\end{figure}

We compute the exact a~priori optimum by decomposing every tour into
segments between always-active vertices, and then analyze an adaptive
policy that changes its sweep direction after inactive probes.  Together,
these two estimates yield the following lower bound.

\Cref{thm:cycle-subdivision} follows from the next two lemmas.

\begin{lemma}\label{lem:cycle-prior}
For every $n\ge 7$, the instance $\mathcal{I}_n^{\mathrm{CSC}}$ satisfies
$\optprior=\frac{15}{8}n$.
\end{lemma}

\begin{lemma}\label{lem:cycle-adapt}
For every $n\ge 7$, the instance $\mathcal{I}_n^{\mathrm{CSC}}$ satisfies
$\frac32n\le\optadapt\le \frac32n+5$.
\end{lemma}

\begin{proof}[Proof of \cref{thm:cycle-subdivision}]
By \cref{lem:cycle-prior,lem:cycle-adapt},
\[
    \optprior=\frac{15}{8}n,
    \qquad
    \frac32n\le\optadapt\le\frac32n+5.
\]
Therefore
\[
    \frac{15n/8}{3n/2+5}
    \le
    \frac{\optprior}{\optadapt}
    \le
    \frac{15n/8}{3n/2}
    =\frac54.
\]
The lower bound also converges to~$5/4$, so the ratio converges to~$5/4$ as
$n\to\infty$.
\end{proof}


\subsection{The a priori optimum}
\label{sec:cycle-prior}
We prove \cref{lem:cycle-prior} in two steps.  First, we bound the
expected contribution of every segment of an arbitrary a~priori tour between
two always-active vertices.  Summing these local bounds gives a global lower
bound.  We then construct tours attaining it.

We begin by recording the relevant distances in the metric completion.  For
always-active vertices,
\begin{equation}\label{eq:cycle-base-distances}
    d(b_i,b_{i+1})=2,
    \qquad
    d(b_i,b_j)=1
    \quad\text{if }j\not\equiv i\pm1\pmod n.
\end{equation}
For a semi-active vertex and an always-active vertex,
\begin{equation}\label{eq:cycle-mixed-distances}
    d(s_i,b_i)=d(s_i,b_{i+1})=1,
    \qquad
    d(s_i,b_j)=2
    \quad\text{if }b_j\notin\{b_i,b_{i+1}\}.
\end{equation}
Finally, for distinct semi-active vertices,
\begin{equation}\label{eq:cycle-subdivision-distances}
    d(s_i,s_j)
    =
    \begin{cases}
        2, & \text{if the base-cycle edges }b_ib_{i+1}
             \text{ and }b_jb_{j+1}\text{ are adjacent},\\
        3, & \text{otherwise.}
    \end{cases}
\end{equation}

Given always-active vertices $u,v$ and an ordered sequence
$W=(w_1,\dots,w_k)$ of distinct semi-active vertices, let
$\Phi(u;W;v)$ be the expected length of the realized path obtained from
\[
    (u,w_1,\dots,w_k,v)
\]
after deleting the semi-active vertices that are inactive.

\begin{lemma}[Segment lower bound]\label{lem:cycle-segment-bound}
Every such segment satisfies
\[
    \Phi(u;W;v)\ge 1+\frac{7}{8}|W|.
\]
\end{lemma}

\begin{proof}
We first treat the cases $k\le3$.

\begin{claim}\label{clm:short-seg}
For every segment $(u;W;v)$ with $k=|W|\le3$,
\[
    \Phi(u;W;v)
    \ge
    \begin{cases}
        1, & k=0,\\
        2, & k=1,\\
        \frac{11}{4}, & k=2,\\
        \frac{29}{8}, & k=3.
    \end{cases}
\]
Moreover, equality is attained for $k=2$ by
$(b_i,s_i,s_{i+1},b_{i+1})$, and for $k=3$ by
$(b_{i+1},s_{i+1},s_i,s_{i-1},b_i)$.
\end{claim}
\begin{proof}
If $k=0$, then the segment contributes the single edge $uv$, so
$\Phi(u;\varnothing;v)=d(u,v)\ge1$.

If $k=1$, say $W=(w)$, then
\[
    \Phi(u;w;v)
    =
    \tfrac12d(u,v)
    +\tfrac12\bigl(d(u,w)+d(w,v)\bigr).
\]
If $u$ and $v$ are consecutive on the base cycle, then $d(u,v)=2$
and also $d(u,w)+d(w,v)\ge2$.  If they are not consecutive, then
$d(u,v)=1$, and no semi-active vertex is incident with both $u$ and $v$,
so $d(u,w)+d(w,v)\ge3$.  In either case
$\Phi(u;w;v)\ge2$.

Now let $k=2$, say $W=(w_1,w_2)$.  There are four realizations.  If
$d(u,v)=1$, then neither $w_1$ nor $w_2$ is incident with both $u$ and
$v$, so each singleton realization has cost at least~$3$, while the
realization with both active has cost at least~$4$.  Hence
\[
    \Phi(u;W;v)
    \ge\frac{1+3+3+4}{4}
    =\frac{11}{4}.
\]
If $d(u,v)=2$, then exactly one semi-active vertex is incident with both
$u$ and~$v$, so at most one singleton realization has cost~$2$; the two
singleton costs therefore sum to at least~$5$, and the realization with
both active has cost at least~$4$.  Thus
\[
    \Phi(u;W;v)
    \ge\frac{2+5+4}{4}
    =\frac{11}{4}.
\]
Equality is attained, for example, by
$(b_i,s_i,s_{i+1},b_{i+1})$, whose four realization costs are
$2,2,3,4$.

Finally, let $k=3$, say $W=(w_1,w_2,w_3)$.  Suppose first that
$d(u,v)=1$.  Then each singleton realization has cost at least~$3$.
A two-vertex realization can have cost~$4$ only if one of its active vertices
is incident with~$u$, the other is incident with~$v$, and the corresponding
base-cycle edges are adjacent.  For fixed nonconsecutive $u$ and $v$, there
is at most one unordered pair of semi-active vertices with this property.
Hence among the three two-vertex realizations, at most one has cost~$4$,
and the other two have cost at least~$5$.  The realization with all three
active has cost at least~$6$.  Therefore
\[
    \Phi(u;W;v)
    \ge\frac{1+9+14+6}{8}
    =\frac{15}{4}
    >\frac{29}{8}.
\]

Now suppose that $d(u,v)=2$.  Then exactly one semi-active vertex is
incident with both $u$ and~$v$, so among the three singleton realizations
at most one has cost~$2$.  Hence the sum of the singleton costs is at least
$2+3+3=8$.  A two-vertex realization can have cost~$4$ only if one of its
active vertices is the unique semi-active vertex on the edge~$uv$, and the
other comes from one of the two base-cycle edges adjacent to~$uv$.  Thus among the three
two-vertex realizations, at most two have cost~$4$, and the remaining one
has cost at least~$5$; the sum of these three costs is therefore at
least~$13$.  The realization with all three active has cost at least~$6$.
Consequently,
\[
    \Phi(u;W;v)
    \ge\frac{2+8+13+6}{8}
    =\frac{29}{8}.
\]
Equality is attained, for example, by
$(b_{i+1},s_{i+1},s_i,s_{i-1},b_i)$, whose eight realization costs are
\[
    2,\ 3,\ 2,\ 3,\ 4,\ 5,\ 4,\ 6.
\]
\end{proof}

It remains to consider $k\ge4$.  Set
\[
    z_0:=u,\qquad
    z_t:=w_t\ (1\le t\le k),\qquad
    z_{k+1}:=v.
\]
Two retained vertices $z_i,z_j$ become consecutive after shortcutting
precisely when every vertex strictly between them is deleted.  Since each
$w_t$ is retained independently with probability~$1/2$, we obtain
\begin{align*}
    \Phi(u;W;v)
    &=
    2^{-k}d(u,v)
    +\sum_{t=1}^k2^{-t}d(u,w_t)
    +\sum_{t=1}^k2^{-(k-t+1)}d(w_t,v)\\
    &\qquad
    +\sum_{1\le i<j\le k}2^{-(j-i+1)}d(w_i,w_j).
\end{align*}
Introduce the indicators
\[
    \alpha_t:=\1_{\{w_t\text{ is incident with }u\}},
    \qquad
    \beta_t:=\1_{\{w_t\text{ is incident with }v\}},
\]
and, for $1\le i<j\le k$,
\[
    \gamma_{ij}
    :=
    \1_{\{w_i,w_j\text{ correspond to adjacent base-cycle edges}\}}.
\]
By~\crefrange{eq:cycle-base-distances}{eq:cycle-subdivision-distances},
\begin{align*}
    d(u,v)
    &=1+\1_{\{u,v\text{ are consecutive on the base cycle}\}},\\
    d(u,w_t)&=2-\alpha_t,\\
    d(w_t,v)&=2-\beta_t,\\
    d(w_i,w_j)&=3-\gamma_{ij}.
\end{align*}
Substituting these identities gives
\[
    \Phi(u;W;v)
    =
    1+\frac{3k}{2}
    +2^{-k}\1_{\{u,v\text{ consecutive}\}}
    -D,
\]
where
\[
    D
    :=
    \sum_{t=1}^k2^{-t}\alpha_t
    +\sum_{t=1}^k2^{-(k-t+1)}\beta_t
    +\sum_{1\le i<j\le k}2^{-(j-i+1)}\gamma_{ij}.
\]
Since an always-active vertex is incident with exactly two semi-active vertices,
\[
    \sum_{t=1}^k2^{-t}\alpha_t
    \le\frac34,
    \qquad
    \sum_{t=1}^k2^{-(k-t+1)}\beta_t
    \le\frac34.
\]
For the last sum, consider the graph on
$\{w_1,\dots,w_k\}$ in which two vertices are adjacent when the
corresponding semi-active vertices arise from adjacent edges of the base
cycle.  This graph is a subgraph of the base cycle, hence has maximum degree at
most~$2$ and at most $k$ edges.  Every edge contributes weight at
most~$1/4$, so
\[
    \sum_{1\le i<j\le k}
    2^{-(j-i+1)}\gamma_{ij}
    \le\frac{k}{4}.
\]
Therefore
\[
    D\le\frac32+\frac{k}{4},
\]
and hence
\[
    \Phi(u;W;v)
    \ge
    1+\frac{3k}{2}
    -\left(\frac32+\frac{k}{4}\right)
    =
    \frac54k-\frac12.
\]
Since $k\ge4$,
\[
    \frac54k-\frac12
    \ge1+\frac{7k}{8}.
\]
Together with \cref{clm:short-seg}, this proves the lemma.
\end{proof}

\begin{proof}[Proof of \cref{lem:cycle-prior}]
Fix an arbitrary a~priori tour
\[
    T=(b_1,x_1,x_2,\dots,x_{2n-1},b_1)
\]
on~$V_n$, where $x_1,\dots,x_{2n-1}$ is a permutation of
$\{b_2,\dots,b_n,s_1,\dots,s_n\}$.  Deleting all semi-active vertices
from~$T$ leaves a cyclic ordering
\[
    (b_1=u_1,u_2,\dots,u_n,b_1)
\]
of the always-active vertices.  For each $j\in[n]$, let
\[
    W_j=(w_{j,1},\dots,w_{j,k_j})
\]
be the ordered list, possibly empty, of semi-active vertices appearing
between $u_j$ and $u_{j+1}$, where $u_{n+1}:=b_1$.  Every semi-active
vertex belongs to exactly one such segment, so
\[
    \sum_{j=1}^n k_j=n.
\]
Because all always-active vertices are present in every realization, the realized
tour decomposes into these segment contributions.  Taking expectation over
$A\sim\mathcal Q$ and applying
\cref{lem:cycle-segment-bound}, we obtain
\begin{align*}
    \cost(T)
    &=\E_{A\sim\mathcal Q}[\cost(T|_A)]\\
    &=
    \sum_{j=1}^n\Phi(u_j;W_j;u_{j+1})\\
    &\ge
    \sum_{j=1}^n\left(1+\frac{7}{8}|W_j|\right)\\
    &=n+\frac78\sum_{j=1}^nk_j
    =
    \frac{15}{8}n.
\end{align*}
Since $T$ was arbitrary,
\[
    \optprior\ge\frac{15}{8}n.
\]

For the reverse inequality, we construct a tour attaining this bound.

\medskip
\noindent\textbf{Even \(n\).}
Write $n=2m$ and consider
\begin{align*}
T^{\mathrm{even}}
:=(&b_1,s_1,s_2,b_2,
   b_4,s_4,s_3,b_3,
   b_6,s_6,s_5,b_5,
   \dots,\\
  &b_{2m},s_{2m},s_{2m-1},b_{2m-1},b_1).
\end{align*}
This tour has $m$ nonempty segments, namely
\[
    (b_1,s_1,s_2,b_2)
    \quad\text{and}\quad
    (b_{2t},s_{2t},s_{2t-1},b_{2t-1})
    \qquad(2\le t\le m),
\]
each with expected contribution $11/4$, together with $m$ direct edges
between always-active vertices, each of length~$1$.  Hence
\begin{align*}
    \cost(T^{\mathrm{even}})
    &=\E_{A\sim\mathcal Q}[\cost(T^{\mathrm{even}}|_A)]\\
    &=
    m\left(\frac{11}{4}+1\right)
    =
    \frac{15}{8}n.
\end{align*}

\medskip
\noindent\textbf{Odd \(n\).}
Write $n=2m+1$ with $m\ge2$ and consider
\begin{align*}
T^{\mathrm{odd}}
:=(&b_n,b_2,s_2,s_1,s_n,b_1,
   b_4,s_4,s_3,b_3,
   b_6,s_6,s_5,b_5,
   \dots,\\
  &b_{2m},s_{2m},s_{2m-1},b_{2m-1},b_n).
\end{align*}
This tour has one segment
\[
    (b_2,s_2,s_1,s_n,b_1)
\]
with expected contribution $29/8$, together with $m-1$ segments
\[
    (b_{2t},s_{2t},s_{2t-1},b_{2t-1})
    \qquad(2\le t\le m),
\]
each contributing $11/4$, and $m+1$ direct edges between always-active vertices,
each of length~$1$.  Therefore
\begin{align*}
    \cost(T^{\mathrm{odd}})
    &=\E_{A\sim\mathcal Q}[\cost(T^{\mathrm{odd}}|_A)]\\
    &=
    \frac{29}{8}
    +(m-1)\frac{11}{4}
    +(m+1)\\
    &=
    \frac{15}{8}(2m+1)
    =
    \frac{15}{8}n.
\end{align*}

In both parity cases we have constructed an a~priori tour of expected cost
$15n/8$.  Together with the lower bound, this proves
\[
    \optprior=\frac{15}{8}n. \qedhere 
\]
\end{proof}

\subsection{An adaptive policy}
\label{sec:cycle-adaptive-policy}

We now prove \cref{lem:cycle-adapt} by first giving a universal lower bound and
then analyzing an adaptive policy with asymptotically matching cost.

\paragraph{The alternating-sweep policy.}
Let
\[
P \;:=\; (b_2,\,s_2,\,b_3,\,\dots,\,s_{n-1},\,b_n)
\]
be viewed as an ordered path.  The policy first probes $s_1$, then $s_n$, and
then the always-active vertex~$b_2$.  Its current location is now~$b_2$, and it
begins sweeping along~$P$ toward~$b_n$.

At every subsequent step, the unprobed vertices of~$P$ form a contiguous
subpath, and the current location is a probed always-active vertex adjacent to
one end of that subpath.  The policy probes the semi-active vertex at this end.
If it is active, the policy visits it, probes the next always-active vertex in
the current direction if one remains, and continues the sweep.  If it is
inactive, the current location is unchanged; the policy probes the
always-active vertex at the opposite end of the remaining subpath and reverses
the sweep direction.  When no unprobed vertex remains, the policy returns to
the root.  This rule preserves the stated invariant and therefore defines the
policy until every non-root vertex has been probed.

\begin{proof}[Proof of \cref{lem:cycle-adapt}]
Every realization contains the $n$ always-active vertices.  If $Y$ denotes the
number of active semi-active vertices, then every realized tour visits $n+Y$
vertices and therefore uses $n+Y$ metric edges.  Each such edge has length at
least~$1$, so every adaptive policy has expected cost at least
\[
    \E[n+Y]=n+\frac n2=\frac32n.
\]

For the upper bound, fix an activation set~$A$.  The movement
generated by the preliminary probes is a shortcut of the walk
$(r,s_1,s_n,b_2)$ and therefore costs at most
\[
    d(r,s_1)+d(s_1,s_n)+d(s_n,b_2)
    =1+2+2=5.
\]

Let
\[
    X:=\bigl|\{i\in\{2,\dots,n-1\}:s_i\text{ is active}\}\bigr|.
\]
Each of the $X$ semi-active vertices on~$P$ that are active causes at most two
unit-length moves: one to that vertex and one to the next always-active
vertex.  A semi-active vertex that is inactive can cause one switch to the
always-active vertex at the opposite end of the remaining subpath.  Every
such switch costs~$1$, except possibly the last, which may cost~$2$ when its endpoints
are consecutive on the base cycle.  Hence the total switch cost is at most
$(n-2)-X+1$.  Finally, the return to~$r$ costs at most~$2$.  Therefore, for
every realization,
\[
    \cost(T(A;\pi))
    \le
    5+2X+\bigl((n-2)-X+1\bigr)+2
    =n+X+6.
\]
Since $\E[X]=(n-2)/2$, it follows that
\[
    \cost(\pi)=\E_{A\sim\mathcal Q}[\cost(T(A;\pi))]
    \le \frac32n+5.
\]
Together with the preceding lower bound, this proves the lemma.
\end{proof}

\section{Randomizing the metric}
\label{sec:random-metric-gap}

Unlike the preceding explicit construction, we now randomize the metric.  An
Erd\H{o}s--R\'enyi graph determines which vertex pairs have distance~$1$;
nonedges have distance~$2$.  Up to an additive $O(1)$ term, which is
asymptotically negligible because the expected number of active vertices tends
to infinity, every realized tour has cost between~$|A|$ and~$2|A|$.  Thus
factor~$2$ is the largest separation possible within this family.  We show
that, with positive probability over the sampled metric, every a~priori tour
has expected cost $(2-o(1))np$, whereas one adaptive policy has expected cost
$(1+o(1))np$.

The random graph is only a probabilistic-method device.  At the end of the
argument we fix a deterministic realization, which is then known completely
to both benchmarks.  The randomness used to prove existence is distinct from
the activation randomness faced by a policy on the resulting instance.

\begin{theorem}[Probabilistic factor-two construction]
\label{thm:random-metric-gap}
There exists a sequence $(\mathcal{I}_n^{\mathrm{RM}})_n$ of deterministic
symmetric \StochTSP{} instances
for which
\[
    \frac{\optprior(\mathcal{I}_n^{\mathrm{RM}})}
         {\optadapt(\mathcal{I}_n^{\mathrm{RM}})}
    \longrightarrow 2.
\]
In particular, the adaptivity gap for symmetric \StochTSP{} is at
least~$2$.
\end{theorem}

\begin{corollary}\label{cor:random-metric-oblivious}
The obliviousness gap for symmetric \StochTSP{} is at least~$2$.
\end{corollary}

\begin{proof}
For every instance, $\optpost\le\optadapt$.  Hence
\[
    \frac{\optprior}{\optpost}
    \ge
    \frac{\optprior}{\optadapt},
\]
and the claim follows from \cref{thm:random-metric-gap}.
\end{proof}

\subsection{The random metric}\label{sec:random-metric-model}

We now specify the random family.  All logs in this section are to
base~$\mathrm{e}$.  Fix a sufficiently large integer~$n$, label the non-root
vertices by $[n]=\{1,\ldots,n\}$, let $V:=\{r\}\cup[n]$, and set
\[
    p:=\frac{1}{(\log n)^4},
    \qquad
    \theta:=\frac{1}{\log n},
    \qquad
    q:=1-p.
\]
Every non-root vertex is active independently with probability~$p$.

Independently for every unordered pair
$e=\{x,y\}\subseteq V$, sample a Bernoulli random variable
$X_e$ with mean~$\theta$, and define
\[
    D(x,y):=2-X_{\{x,y\}}
    \quad (x\neq y),
    \qquad
    D(x,x):=0.
\]
Thus a pair has distance~$1$ with probability~$\theta$ and distance~$2$
otherwise.  Let $H(D)$ be the graph whose edge set is
$\{e:X_e=1\}$.  Then $H(D)\sim G_{n+1,\theta}$, and for $k\ge1$ we use
$G_{k,\theta}$ for the same distribution on any fixed set of $k$ vertices.
We call the edges of~$H(D)$ \emph{short}.  Every realization of~$D$ is a
symmetric $(1,2)$-TSP metric.
For every deterministic realization~$d$ of~$D$, let
$\mathcal{I}_n^{\mathrm{RM}}(d)$ denote the resulting random-metric
\StochTSP{} instance on~$V$ with root~$r$: its metric is~$d$, and its
non-root vertices are activated independently with probability~$p$.

Let $A$ be the random set of active vertices, and let $K:=|A|$.  Then
$K\sim\operatorname{Binomial}(n,p)$.
Define
\begin{equation}\label{eq:random-metric-baseline}
    W
    :=
    \E\bigl[(K+1)\1_{\{K\ge1\}}\bigr]
    =
    np+1-q^n.
\end{equation}
Since $np=n/(\log n)^4\to\infty$ and $0\le1-q^n\le1$,
\[
    W=np+O(1)=(1+o(1))np.
\]
For a nonempty realization with $K=k$, every feasible tour has exactly
$k+1$ edges.  Since all distances are at least~$1$, $W$ is a
lower bound on the expected cost of every benchmark.  Moreover,
\begin{equation}\label{eq:random-metric-baseline-lower-bound}
    W\ge np=\frac{n}{(\log n)^4}.
\end{equation}

\subsection{A uniform lower bound for a priori tours}
\label{sec:random-metric-prior-bound}

We next show that, with high probability over the metric, every a~priori tour
has nearly the same large expected cost.  For a fixed order, its expected cost
over the activations is a weighted sum of the independent random distances.
The weights are the probabilities that the corresponding pairs become
consecutive in the shortcut.  We identify these weights, bound their squared
sum, and then apply Hoeffding's inequality.  The random graph is invariant
under relabeling the non-root vertices, so a sufficiently strong tail bound for
the canonical order $r,1,\ldots,n,r$ applies to every order by symmetry.  A
union bound over the at most $n!$ orders then gives a simultaneous lower bound.

By the definition of a~priori cost in~\cref{sec:stoch-tsp-benchmarks}, for a
deterministic realization~$d$ of~$D$ and an a~priori tour
\[
    T=(r,v_1,v_2,\ldots,v_n,r)
\]
we write
\[
    C_d(T):=\E_{A\sim\mathcal Q}[\cost_d(T|_A)]=\cost_d(T).
\]

Recall that $q=1-p$.

\begin{lemma}[Expected-cost expansion]\label{lem:random-metric-expansion}
For $i\in[n]$, define
\[
    \alpha_i:=p\bigl(q^{i-1}+q^{n-i}\bigr),
\]
and, for $1\le i<j\le n$, define
\[
    \beta_{ij}:=p^2q^{j-i-1}.
\]
Then
\begin{equation}\label{eq:random-metric-expansion}
    C_d(T)
    =
    \sum_{i=1}^n\alpha_i d(r,v_i)
    +
    \sum_{1\le i<j\le n}\beta_{ij}d(v_i,v_j).
\end{equation}
The coefficients in~\cref{eq:random-metric-expansion} sum to~$W$.
\end{lemma}

\begin{proof}
The edge from $r$ to $v_i$ is used when $v_i$ is active and
$v_1,\ldots,v_{i-1}$ are inactive, which occurs with probability $pq^{i-1}$.
Similarly, the edge from $v_i$ to~$r$ is used with probability
$pq^{n-i}$.  For $i<j$, the vertices $v_i$ and $v_j$ are consecutive
active vertices in the shortcut precisely when both are active and all vertices
strictly between them are inactive.  This event has probability
$p^2q^{j-i-1}$, proving~\cref{eq:random-metric-expansion}.

If every distance in the expansion is replaced by~$1$, the right-hand side
is the expected number of edges in the shortcut tour.  This number is
$(K+1)\1_{\{K\ge1\}}$, whose expectation is~$W$.
\end{proof}

\begin{lemma}[Coefficient-square bound]
\label{lem:random-metric-coefficients}
For every a~priori order and all sufficiently large~$n$,
\[
    \sum_{i=1}^n\alpha_i^2
    +
    \sum_{1\le i<j\le n}\beta_{ij}^2
    \le 2np^3.
\]
\end{lemma}

\begin{proof}
Using $(a+b)^2\le2a^2+2b^2$,
\begin{align*}
    \sum_{i=1}^n\alpha_i^2
    &\le
    2p^2\sum_{i=1}^n
        \bigl(q^{2(i-1)}+q^{2(n-i)}\bigr)\\
    &\le
    \frac{4p^2}{1-q^2}
    =
    \frac{4p}{2-p}
    \le4p.
\end{align*}
Also,
\begin{align*}
    \sum_{1\le i<j\le n}\beta_{ij}^2
    &=
    p^4\sum_{g=0}^{n-2}(n-g-1)q^{2g}\\
    &\le
    \frac{np^4}{1-q^2}
    =
    \frac{np^3}{2-p}
    \le np^3.
\end{align*}
Since $np^2\to\infty$, we have $4p\le np^3$ for all sufficiently
large~$n$.  Adding the two estimates proves the claim.
\end{proof}

\begin{lemma}[Uniform a~priori lower bound]
\label{lem:random-metric-prior}
With probability $1-o(1)$ over the random metric~$D$,
\[
    \optprior(\mathcal{I}_n^{\mathrm{RM}}(D))
    \ge
    \left(2-\frac{2}{\log n}\right)W.
\]
\end{lemma}

\begin{proof}
Consider the canonical tour $T_0=(r,1,\ldots,n,r)$, and let $c_e$ be the
coefficient of the unordered pair~$e$ in~\cref{eq:random-metric-expansion}.
Since $D(e)=2-X_e$, $\E[X_e] = \theta$, and the
coefficients sum to~$W$,
\[
    C_D(T_0)=2W-\sum_e c_eX_e,
    \qquad
    \E_D[C_D(T_0)]=(2-\theta)W.
\]
Thus the canonical tour has unusually low cost only if the total coefficient
weight of distance-$1$ shortcut pairs is unusually large.

Set $\varepsilon:=1/\log n$.  The event
\[
    C_D(T_0)\le(2-\theta-\varepsilon)W
\]
is equivalent to
\[
    \sum_e c_e(X_e-\theta)\ge\varepsilon W.
\]
The variables $X_e$ are independent and each $c_eX_e$ lies in~$[0,c_e]$.
Hence the weighted Hoeffding bound~\cref{eq:weighted-hoeffding} gives
\[
    \P_D\left[
        C_D(T_0)\le(2-\theta-\varepsilon)W
    \right]
    \le
    \exp\left(
        -\frac{2\varepsilon^2W^2}{\sum_ec_e^2}
    \right).
\]
For all sufficiently large~$n$, the coefficient bound
\cref{lem:random-metric-coefficients} applies; together with $W\ge np$, it gives
\begin{align*}
    \P_D\left[
        C_D(T_0)\le(2-\theta-\varepsilon)W
    \right]
    &\le
    \exp\left(-\frac{\varepsilon^2W^2}{np^3}\right)\\
    &\le
    \exp\left(-\frac{\varepsilon^2n}{p}\right)\\
    &=
    \exp\bigl(-n(\log n)^2\bigr).
\end{align*}

For any other a~priori tour~$T$, a permutation of the non-root labels maps
$T$ to~$T_0$.  Because the distribution of~$D$ is invariant under every such
permutation, the same tail bound holds for~$C_D(T)$.  There are at most $n!$
master tours with the root fixed, so a union bound shows that the probability
that some tour violates the claimed lower bound is at most
\[
    n!\exp\bigl(-n(\log n)^2\bigr)
    \le
    \exp\bigl(n\log n-n(\log n)^2\bigr)
    =
    o(1).
\]
The conclusion follows because
$\theta=\varepsilon=1/\log n$.
\end{proof}

\subsection{A low-cost adaptive policy}\label{sec:random-metric-adaptive-policy}

To complement the uniform lower bound, we construct a policy that exhausts
distance-$1$ unprobed neighbors before making a distance-$2$ move.  If $s$
active vertices remain when the policy reaches a fresh vertex, the probability
that at least one of them is a distance-$1$ neighbor is
$1-(1-\theta)^s$.  We bound the average cost of this policy over the random
metric and then fix a metric
realization for which this bound and the preceding a~priori bound hold
simultaneously.

Fix a deterministic realization~$d$ in the support of~$D$.  Its short-edge
graph $H(d)$, and hence the whole metric, is known before any activations are
probed.  Define the \emph{nearest-unprobed-neighbor policy}~$\pi_d$ as follows.
At a state $(u,S)$ with $S\ne\varnothing$, choose a vertex~$v\in S$ that
minimizes $d(u,v)$, breaking ties by the labels $1,\ldots,n$, and probe~$v$.
If $S=\varnothing$, return to~$r$.  Equivalently, the policy probes every
distance-$1$ unprobed neighbor of its current location before it probes a
vertex at distance~$2$.

The metric is sampled in advance and is known to the policy.  For the
analysis, however, we may reveal its independent edge variables only when
the policy uses them.  The following invariant formalizes this deferred
exposure.

\begin{lemma}[Fresh-edge invariant]\label{lem:random-metric-fresh-edges}
Consider the joint experiment over the random metric~$D$ and the active
set~$A$, and expose only the metric edges incident to the policy's current
location.  Initially, when $u=r$, and immediately after every arrival at a
new active vertex~$u$, let $S$ be the set of unprobed vertices.  Conditional
on the complete exposure and probe history so far, no edge induced by
$\{u\}\cup S$ has been exposed.  Consequently, the conditional distribution
of the short-edge graph induced by $\{u\}\cup S$ is
$G_{1+|S|,\theta}$, independently of the activation states on~$S$.
\end{lemma}

\begin{proof}
At the initial state no metric edge has been exposed.  Suppose the invariant
holds at a current vertex~$u$.  Reveal the edges $\{u,v\}$ for $v\in S$;
these are the only metric variables needed to determine the probes made while
the policy remains at~$u$.  If an active probed vertex~$v$ is reached,
let $S'\subseteq S\setminus\{v\}$ be the vertices still unprobed.  The edges
induced by $\{v\}\cup S'$ were not among the edges just exposed, and they
were not exposed earlier by the induction hypothesis.  Independence of the
edge variables gives their conditional law $G_{1+|S'|,\theta}$.  The metric
and activation processes are independent, so this law is also independent
of the activation states on~$S'$.  This proves the invariant by induction.
\end{proof}

\begin{lemma}[Cost of the adaptive policy]
\label{lem:random-metric-adaptive-average}
Let
\[
    U(d):=\cost_d(\pi_d)
    =\E_{A\sim\mathcal Q}[\cost_d(T(A;\pi_d))].
\]
Then $U(d)\ge W$ for every~$d$, and
\[
    \E_D[U(D)]
    \le
    np+\left(\frac1\theta+1\right)(1-q^n)
    \le
    W+\frac1\theta.
\]
\end{lemma}

\begin{proof}
Fix $k\ge1$ and condition on $K=k$.  Consider the execution jointly over
the random metric and the active set.  Immediately after the policy
arrives at an active vertex~$u$, let $S$ be the set of unprobed vertices and
suppose that exactly $s$ vertices in~$S$ are active.

By \cref{lem:random-metric-fresh-edges}, the edges from~$u$ to these $s$
active vertices are conditionally independent, and each is short with
probability~$\theta$.  The next edge used has length~$2$ exactly when
none of them is short.  Its conditional expected length is therefore
\[
    1+(1-\theta)^s.
\]
As the policy visits the $k$ active vertices, the value of~$s$ before the
successive edges used is $k,k-1,\ldots,1$.  The final return to the root
costs at most~$2$.  Hence
\begin{align*}
    \E_{D,A}[\cost_D(T(A;\pi_D))\mid K=k]
    &\le
    \sum_{s=1}^k\bigl(1+(1-\theta)^s\bigr)+2\\
    &\le
    k+\frac{1-\theta}{\theta}+2\\
    &=
    k+\frac1\theta+1.
\end{align*}
For $K=0$ the cost is~$0$.  Averaging over~$K$ gives
\[
    \E_D[U(D)]
    \le \E[K]
       +\left(\frac1\theta+1\right)\P[K\ge1]
    =W+\frac{\P[K\ge1]}{\theta}
    \le W+\frac1\theta.
\]

For the pointwise lower bound, if $K=k\ge1$, the realized tour uses exactly
$k+1$ edges, each of length at least~$1$.  Taking expectation over the active
set gives $U(d)\ge W$.
\end{proof}

Since $W\ge np$,
\[
    \theta W\ge\frac{n}{(\log n)^5}\longrightarrow\infty.
\]
Set $\delta_n:=1/\sqrt{\theta W}=o(1)$.  The following lemma shows that the
adaptive optimum is at most $(1+o(1))np$ for most sampled metrics.

\begin{lemma}[High-probability adaptive upper bound]
\label{lem:random-metric-adaptive-whp}
\[
    \P_D\bigl[
        \optadapt(\mathcal{I}_n^{\mathrm{RM}}(D))
        \le(1+\delta_n)W
    \bigr]
    \ge1-\delta_n.
\]
\end{lemma}

\begin{proof}
The nonnegative random variable $U(D)-W$ satisfies
\[
    \E_D[U(D)-W]\le\frac1\theta
\]
by \cref{lem:random-metric-adaptive-average}.  Markov's inequality gives
\begin{align*}
    \P_D[U(D)>(1+\delta_n)W]
    &=
    \P_D[U(D)-W>\delta_nW]\\
    &\le
    \frac{1/\theta}{\delta_nW}
    =
    \frac1{\sqrt{\theta W}}
    =
    \delta_n.
\end{align*}
Since $\optadapt(\mathcal{I}_n^{\mathrm{RM}}(d))\le U(d)$ for every~$d$,
the lemma follows.
\end{proof}

\begin{proof}[Proof of \cref{thm:random-metric-gap}]
By \cref{lem:random-metric-prior}
and~\cref{lem:random-metric-adaptive-whp}, with probability $1-o(1)$ the
sampled metric satisfies both
\[
    \optprior(\mathcal{I}_n^{\mathrm{RM}}(D))
    \ge
    \left(2-\frac{2}{\log n}\right)W
\]
and
\[
    \optadapt(\mathcal{I}_n^{\mathrm{RM}}(D))
    \le(1+\delta_n)W.
\]
The two events therefore have a nonempty intersection for all sufficiently
large~$n$.  Choose a deterministic realization~$d_n$ in this intersection
and set
\[
    \mathcal{I}_n^{\mathrm{RM}}
    :=\mathcal{I}_n^{\mathrm{RM}}(d_n).
\]
Then
\[
    \frac{\optprior(\mathcal{I}_n^{\mathrm{RM}})}
         {\optadapt(\mathcal{I}_n^{\mathrm{RM}})}
    \ge
    \frac{2-2/\log n}{1+\delta_n}
    =
    2-o(1).
\]

It remains to verify the stated convergence, rather than only a lower limit.
For every metric in the support of the construction, every shortcut tour on a
nonempty active set of size~$k$ uses $k+1$ edges of length at most~$2$.
Consequently, $\optprior(\mathcal{I}_n^{\mathrm{RM}})\le2W$.  On the other
hand, every adaptive policy has expected cost at least~$W$.  Thus
\[
    \frac{\optprior(\mathcal{I}_n^{\mathrm{RM}})}
         {\optadapt(\mathcal{I}_n^{\mathrm{RM}})}
    \le2.
\]
Combining the two bounds proves that the ratio converges to~$2$.
\end{proof}

% ============================================================
% CHAPTER 3: CLAIRVOYANCE GAP — SYMMETRIC
% ============================================================
\chapter{Clairvoyance Gap for Symmetric Metrics}\label{ch:sym-clair}

\providecommand{\Gammaone}{\Gamma_{1}}

In this chapter, we prove that the clairvoyance gap for symmetric \StochTSP{}
is at least~$10/9$ by constructing a family of instances with one semi-active
vertex $s$ of activation probability~$1/2$ for which
\[
  \frac{\optadapt}{\optpost}
  =\frac{10m+2}{9m+2}
  \longrightarrow\frac{10}{9}.
\]
Throughout this chapter, all instances discussed are symmetric.

The optimal visit order when $s$ is inactive conflicts with the optimal order
when $s$ is active.  An adaptive policy, however, must begin probing before it
learns which case applies.

As a complementary result, we prove that every instance with one
semi-active vertex of activation probability~$p\in[0,1]$ satisfies
\[
  \frac{\optadapt}{\optpost}\le\frac{2+p}{2+p^2}.
\]
The right-hand side is maximized at $p=\sqrt6-2$.  Therefore the following
uniform bound holds for instances with exactly one semi-active vertex:
\[
  \frac{\optadapt}{\optpost}\le\frac{2+\sqrt6}{4}.
\]

\section{The three-chain construction}\label{sec:three-chain-construction}

We use a family of instances with one semi-active vertex to obtain a lower
bound of~$10/9$ on the clairvoyance gap.  When $s$ is inactive, one optimal
tour uses the order $A,B,C$; when $s$ is active, one uses the conflicting order
$A,C,s,B$.  An a posteriori tour can choose the appropriate order after seeing
whether $s$ is active, but an adaptive policy must commit to part of its route
before it has this information.

Fix an integer $m\ge3$ and set
\begin{equation}\label{eq:HUP}
  H:=2m,
  \qquad
  U:=3m+2,
  \qquad
  P:=4m=2H.
\end{equation}

\begin{construction}\label{con:three-chain}
The graph contains the root $r=a_0$, one semi-active vertex $s$, and three
paths
\[
  A=(a_0,a_1,\ldots,a_m),
  \qquad
  B=(b_0,b_1,\ldots,b_m),
  \qquad
  C=(c_0,c_1,\ldots,c_m).
\]
Every path edge has length $2$, so each chain has total length $H=2m$.
Add four edges of length $U$, which we call the \emph{$U$-edges},
\begin{equation}\label{eq:u-edges}
  a_0b_m,
  \qquad
  a_0c_m,
  \qquad
  a_mb_0,
  \qquad
  a_mc_0,
\end{equation}
and three edges of length $P$, which we call the \emph{$P$-edges},
\begin{equation}\label{eq:p-edges}
  b_mc_0,
  \qquad
  b_0s,
  \qquad
  c_ms.
\end{equation}
The $P$-edges provide two ways to connect $B$ and $C$: directly through
$b_mc_0$, or through $s$ using $b_0s$ and $sc_m$.  The $U$-edges are the
remaining connections between chains, joining $A$ to either $B$ or $C$.
The metric $d$ is the shortest-path metric of this graph.  Every non-root
vertex other than $s$ is always active, while $s$ is active with probability
$1/2$.  Denote the resulting three-chain instance by
$\mathcal{I}_m^{\mathrm{TC}}$.
\end{construction}

\Cref{fig:three-chain} shows the topology of the construction.

\begin{figure}[htbp]
\centering
\begin{tikzpicture}[
    x=1cm,y=1cm,
    deterministic/.style={rectangle,draw=black,fill=white,minimum width=8.5mm,
      minimum height=7mm,inner sep=0pt,font=\scriptsize},
    root/.style={deterministic,fill=gray!15},
    stochastic/.style={circle,draw=black,fill=gray!30,minimum size=7.5mm,
      inner sep=0pt,font=\scriptsize},
    internal/.style={black,line width=0.75pt},
    uedge/.style={black!70,densely dashed,line width=0.8pt},
    pedge/.style={black,dash pattern=on 4pt off 2pt on 1pt off 2pt,
      line width=0.8pt},
    rowlabel/.style={font=\small\bfseries,anchor=east}
]
\node[root]          (A0) at (0, 2.4) {$a_0$};
\node[deterministic] (A1) at (1.5, 2.4) {$a_1$};
\node[deterministic] (Am) at (6, 2.4) {$a_m$};

\node[deterministic] (Bm) at (0, 0) {$b_m$};
\node[deterministic] (B1) at (4.5, 0) {$b_1$};
\node[deterministic] (B0) at (6, 0) {$b_0$};

\node[deterministic] (C0) at (0,-2.4) {$c_0$};
\node[deterministic] (C1) at (1.5,-2.4) {$c_1$};
\node[deterministic] (Cm) at (6,-2.4) {$c_m$};
\node[stochastic]    (S)  at (7.3,-1.2) {$s$};

\draw[uedge] (A0)--(Bm);
\draw[uedge] (A0)--(Cm);
\draw[uedge] (Am)--(B0);
\draw[uedge] (Am)--(C0);

\draw[pedge] (Bm)--(C0);
\draw[pedge] (B0)--(S);
\draw[pedge] (Cm)--(S);

\draw[internal] (A0)--(A1)--(Am);
\draw[internal] (Bm)--(B1)--(B0);
\draw[internal] (C0)--(C1)--(Cm);

\node[fill=white,inner sep=3pt] at (3.45, 2.4) {$\cdots$};
\node[fill=white,inner sep=3pt] at (2.1, 0) {$\cdots$};
\node[fill=white,inner sep=3pt] at (3.45,-2.4) {$\cdots$};

\node[rowlabel] at (-0.65, 2.4) {chain $A$};
\node[rowlabel] at (-0.65, 0) {chain $B$};
\node[rowlabel] at (-0.65,-2.4) {chain $C$};
\node[above=1mm of A0,font=\scriptsize] {root $r$};
\end{tikzpicture}
\caption{The three-chain graph.  Solid chain edges have length~$2$, dashed
$U$-edges have length~$U$, and dash-dotted $P$-edges have length~$P$.  The
gray circle is the only semi-active vertex.}
\label{fig:three-chain}
\end{figure}


To compute the a posteriori benchmark, we identify an optimal tour in each
case.  When $s$ is inactive, an optimal tour is $T_0$, which visits all
vertices on the three always-active chains in the order $A,B,C$:
\begin{equation}\label{eq:T0}
\begin{split}
  T_0:\quad
  &a_0,a_1,\ldots,a_m,
  b_0,b_1,\ldots,b_m,\\
  &c_0,c_1,\ldots,c_m,a_0,
\end{split}
\end{equation}
When $s$ is active, an optimal tour is $T_1$, with order $A,C,s,B$:
\begin{equation}\label{eq:T1}
\begin{split}
  T_1:\quad
  &a_0,a_1,\ldots,a_m,
  c_0,c_1,\ldots,c_m,\\
  &s,b_0,b_1,\ldots,b_m,a_0.
\end{split}
\end{equation}
Their costs are
\begin{align}
  L_0&:=\cost(T_0)=3H+2U+P=16m+4,
  \label{eq:L0}\\
  L_1&:=\cost(T_1)=3H+2U+2P=20m+4.
  \label{eq:L1}
\end{align}

Intuitively, the conflict lies in which chain is visited first after~$A$:
$T_0$ visits $B$ before $C$, whereas $T_1$ visits $C$ before $B$.

Here is how this conflict gives a lower bound for every adaptive policy.  Let
$C_0$ and $C_1$ be the costs of a policy when $s$ is inactive and active,
respectively.  Until the policy probes~$s$, it receives the same information
in both cases and must therefore follow the same initial route up to that
probe.  We show that when $s$ is inactive, a route with cost $C_0<L_1$ must have the same chain order
as~$T_0$; while when $s$ is active, a route with cost $C_1<L_1+P$ must have the same chain order as
$T_1$, up to reversing the entire tour.  In the second
case, the policy must use either $a_mc_0$ or $a_0b_m$ before probing~$s$, but then neither edge can occur in a route of cost below~$L_1$ when $s$ is inactive.
Thus the two inequalities cannot hold simultaneously.

Consequently, either $C_0\ge L_1$ or $C_1\ge L_1+P$.  We also always have
$C_0\ge L_0$ and $C_1\ge L_1$, because $T_0$ and $T_1$ are optimal in their
respective cases.  If $C_0\ge L_1$, then both costs are at least~$L_1$.  If
$C_1\ge L_1+P$, then $C_0+C_1\ge L_0+L_1+P=2L_1$.  Since the two cases each
have probability~$1/2$, every adaptive policy has expected cost at least
$L_1$.  A policy that follows~$T_1$ costs at most~$L_1$ whether $s$ is active
or inactive, so the optimal adaptive cost is exactly~$L_1$.
\Cref{lem:inactive-low-cost} characterizes all routes of cost below~$L_1$
when $s$ is inactive, and \cref{lem:active-low-cost} characterizes all routes
of cost below~$L_1+P$ when $s$ is active.
\Cref{lem:adaptive-lower-bound} combines these two characterizations to prove
that every adaptive policy has expected cost at least~$L_1$.

\subsection{Adaptive policy lower bound}
\label{sec:three-chain-low-cost-routes}

Because $d$ is the shortest-path metric, every tour move can be realised by moving along
graph edges without changing its cost.  We therefore treat tours as
closed walks in the graph.  If a shortest path is not unique, fix one choice
and use its reverse for the reverse move.  Merely passing through an
intermediate vertex does not probe it, while probing an inactive vertex adds
no movement to the walk.
We call each traversal of one of the seven non-chain edges in
\cref{eq:u-edges,eq:p-edges} a \emph{crossing}.  \Cref{lem:chain-contribution}
provides the cost bounds used throughout this section: it shows how much the
chain edges contribute and when connecting to a chain at the same endpoint
forces additional cost.  These bounds are then used to determine how a
low-cost tour can move between the chains.

\begin{lemma}\label{lem:chain-contribution}
Let $W$ be the closed graph walk representing a tour that visits all vertices
on $A,B,C$.
\begin{enumerate}[label=(\roman*)]
  \item The internal edges of each chain contribute at least $H$ to
  $\cost(W)$.
  \item If a chain has exactly two crossing incidences and both are at the
  same endpoint, the internal edges of that chain contribute at least $2H$.
  \item Suppose all four $U$-edges in \cref{eq:u-edges} are traversed
  exactly once.  Then $A$ has two external incidences at each endpoint, and
  its internal contribution is at least
  \[
    2(H-2)=2H-4.
  \]
\end{enumerate}
\end{lemma}

\begin{proof}
Each chain is a path connected to the rest of the graph only at its two
endpoints.  If
two internal edges were unused, the vertices in a component strictly between
two unused edges would be disconnected from $W$.  Hence at most
one internal edge can be unused.

If every internal edge is used, the contribution is at least $H$.  If exactly
one length-$2$ edge is unused, each of the two remaining path components is
attached to the rest of the walk at only one endpoint.  Since the walk is
closed, every used edge in either component is traversed at least twice.
Consequently the contribution is at least $2(H-2)\ge H$, proving (i).

For (ii), both external incidences occur at one endpoint.  To visit the other
endpoint and return, every internal edge must be traversed at least twice, for
a contribution of at least $2H$.

For (iii), the number of external incidences at each endpoint of $A$ is even.
Because $W$ is closed, every vertex has even degree in its edge multiset.
Thus every used edge of the path $A$ has even multiplicity.  At most one
internal edge can be unused, so the contribution is at least $2(H-2)$.
\end{proof}

We next determine which connections between the chains are possible when the
tour cost is smaller than the thresholds stated above.  The first lemma
considers the case in which $s$ is inactive, and the second considers the case
in which it is active.

\begin{lemma}\label{lem:inactive-low-cost}
When $s$ is inactive, the optimal tour cost is $L_0$.  Moreover, every tour of
cost strictly less than $L_1$ has, up to reversal, exactly the following three
crossings:
\begin{equation}\label{eq:inactive-crossings}
  a_mb_0,
  \qquad
  b_mc_0,
  \qquad
  c_ma_0.
\end{equation}
\end{lemma}

\begin{proof}
Represent a tour visiting all always-active vertices as a closed graph walk.
If this walk passes through $s$ on one of its chosen shortest paths, then it
has at least two crossings incident to $A$, each of cost $U$, and at least
two crossings incident to $s$, each of cost $P$.  By
\cref{lem:chain-contribution}(i), its cost is at least
\[
  3H+2U+2P=L_1.
\]
Thus a tour of cost below $L_1$ does not traverse $s$.

Contract each of $A,B,C$ to one vertex.  Let $x,y,z$ be the respective numbers
of $A$--$B$, $A$--$C$, and $B$--$C$ crossings.  The contracted multigraph is
connected and Eulerian, so
\[
  x+y\equiv x+z\equiv y+z\equiv0\pmod2.
\]
Hence $x,y,z$ have the same parity.

Suppose first that they are odd.  Then $x,y,z\ge1$.  If one is at least $3$,
the crossing cost is at least $4U+P$, and the total cost is at least
\[
  3H+4U+P>L_1,
\]
because $2U>P$.  Therefore a tour below $L_1$ must have
$x=y=z=1$.  Its contracted walk is the triangle through $A,B,C$ and its
crossing cost is $2U+P$.

The $B$--$C$ crossing is necessarily $b_mc_0$.  For $B$ to use opposite
endpoints, the $A$--$B$ crossing must be $a_mb_0$; for $C$ to use opposite
endpoints, the $A$--$C$ crossing must be $a_0c_m$.  These choices also make
$A$ use opposite endpoints.  Conversely, if one chain uses the same endpoint
twice, at least one other chain does as well.  By
\cref{lem:chain-contribution}, a noncanonical port choice has internal
contribution at least $2H+2H+H=5H$.  Its total cost is therefore at least
\[
  5H+2U+P=L_1,
\]
where we used $P=2H$.  Hence every tour below $L_1$ in the odd case has the
three crossings in \cref{eq:inactive-crossings}, up to reversal.  Its cost is
at least $3H+2U+P=L_0$, and $T_0$ attains this bound.

Now suppose $x,y,z$ are even.  Connectivity implies that at least two of them
are positive.  If $z>0$, then there are at least two $B$--$C$ crossings and at
least two crossings incident to $A$, so the total cost is at least
\[
  3H+2U+2P=L_1.
\]
Assume therefore that $z=0$.  Then $x,y\ge2$.  If $x+y\ge6$, the total cost is
at least $3H+6U>L_1$.  The only remaining possibility is
\[
  x=y=2,
  \qquad
  z=0.
\]
If $B$ or $C$ uses the same endpoint twice, the total internal contribution is
at least $4H$, and hence the total cost is at least
\[
  4H+4U=L_1+4>L_1.
\]
Otherwise both $B$ and $C$ use opposite endpoints.  Hence each of the four
$U$-edges is used exactly once.  By
\cref{lem:chain-contribution}(iii), the internal contribution is at least
\[
  H+H+(2H-4)=4H-4.
\]
Thus the total cost is at least
\[
  4H-4+4U
  =20m+4
  =L_1.
\]
This excludes the even case below $L_1$ and completes the proof.
\end{proof}

\begin{lemma}\label{lem:active-low-cost}
When $s$ is active, the optimal tour cost is $L_1$.  Moreover, every tour of
cost strictly less than $L_1+P$ has, up to reversal, exactly the following
four crossings:
\begin{equation}\label{eq:active-crossings}
  a_mc_0,
  \qquad
  c_ms,
  \qquad
  sb_0,
  \qquad
  b_ma_0.
\end{equation}
\end{lemma}

\begin{proof}
Let $W$ be the closed graph walk representing any tour that visits all active
vertices.  It has at least two crossings incident to $A$, each of cost $U$,
and at least two crossings incident to $s$, each of cost $P$.  Therefore
\[
  \cost(W)\ge3H+2U+2P=L_1
\]
by \cref{lem:chain-contribution}(i), and $T_1$ attains this value.

Now consider a tour of cost below $L_1+P$, and let $N_\times$ be its total number of
crossings.  After contracting $A,B,C$ and retaining $s$, we obtain a connected
Eulerian multigraph on four vertices.  Hence every contracted vertex has
positive even degree and $N_\times\ge4$.

If $N_\times=4$, every contracted degree is $2$, so the quotient is the four-cycle
\[
  A-B-s-C-A
\]
up to reversal.  The two edges incident to $s$ are necessarily $b_0s$ and
$c_ms$.  In order that $B$ and $C$ use opposite endpoints, the remaining two
crossings must be $a_0b_m$ and $a_mc_0$; these choices also make $A$ use
opposite endpoints.  Any other port choice makes at least two chains use the
same endpoint twice.  Its internal contribution is then at least $5H$, so its
cost is at least
\[
  5H+2U+2P=L_1+P.
\]
Thus strict inequality forces exactly the four crossings in
\cref{eq:active-crossings}, up to reversal.

If $N_\times=5$, let $N_A$ and $N_s$ be the numbers of crossings incident to
$A$ and $s$, respectively, and let $N_{BC}$ count the $B$--$C$ crossings.
The first two numbers are positive and even.  Every crossing belongs to
exactly one of these three categories, so
$N_\times=5$ forces
\[
  N_A=2,
  \qquad
  N_s=2,
  \qquad
  N_{BC}=1.
\]
The crossing cost is at least $2U+3P$, and hence the total cost is at least
\[
  3H+2U+3P=L_1+P.
\]

Finally, if $N_\times\ge6$, the two mandatory $A$-crossings and two mandatory
$s$-crossings are supplemented by at least two further crossings.  Every
additional crossing costs at least $U$, so
\[
  \cost(W)
  \ge3H+4U+2P
  >3H+2U+3P
  =L_1+P,
\]
where the strict inequality is $2U>P$.  This proves the claim.
\end{proof}

We now combine the two descriptions of low-cost tours to obtain the adaptive
lower bound.  The executions when $s$ is active and inactive are identical
until $s$ is probed, so an adaptive policy cannot follow both low-cost
crossing patterns.

\begin{lemma}\label{lem:adaptive-lower-bound}
Every adaptive policy for $\mathcal{I}_m^{\mathrm{TC}}$ has expected cost at
least
\[
  L_1=20m+4.
\]
\end{lemma}

\begin{proof}
Let $C_0$ and $C_1$ be the costs of an arbitrary adaptive policy when $s$ is
inactive and active, respectively.  Suppose, for contradiction, that its
expected cost is smaller than~$L_1$, so that
\[
  \frac{C_0+C_1}{2}<L_1.
\]
By \cref{lem:inactive-low-cost,lem:active-low-cost}, every valid tour in the
two cases costs at least $L_0$ and $L_1$, respectively.  Hence
\[
  C_0<L_1,
  \qquad
  C_1<2L_1-L_0=L_1+P.
\]

Represent every movement by the fixed shortest paths chosen above.  Before
the policy probes $s$, all probed vertices are always active.  Therefore the
execution when $s$ is active and the execution when $s$ is inactive make
exactly the same probes, occupy the same locations, and traverse the same
graph edges up to the moment at which $s$ is probed.

By \cref{lem:inactive-low-cost}, the inactive execution uses only the three
crossings in \cref{eq:inactive-crossings}.  By
\cref{lem:active-low-cost}, the active execution uses only the four crossings
in \cref{eq:active-crossings}.

No movement in the active execution can pass through $s$ before it is probed:
such a transit would already use two crossings incident to~$s$, while the
mandatory post-probe movement to the active vertex~$s$ would use another;
the four crossings in \cref{eq:active-crossings} include only two incident
to~$s$.  After contracting $A,B,C$, the active execution follows a four-cycle.
It starts in $A$, so the first visit to $s$ is preceded by one of the quotient
paths
\[
  A-C-s
  \qquad\text{or}\qquad
  A-B-s.
\]
The corresponding chain $X\in\{B,C\}$ has exactly two crossing incidences in
the entire walk: one from $A$ and one to $s$.  The $X$--$s$ crossing occurs in
the movement immediately after the remote probe of the active vertex~$s$.
The $A$--$X$ crossing must occur earlier.  Otherwise the execution has not
entered $X$ before probing~$s$, so none of the always-active vertices of~$X$
has been probed: probing any one of them would force an immediate movement
into~$X$.  After the movement to~$s$, both crossings incident to $X$ would
already have been used, and visiting any remaining vertex of~$X$ would require
an additional crossing, contradicting \cref{lem:active-low-cost}.

Thus, before probing $s$, the common prefix has already used either
\[
  a_mc_0
  \qquad\text{or}\qquad
  a_0b_m,
\]
depending on the orientation of the active-case tour.  Neither crossing occurs
among the three crossings in \cref{eq:inactive-crossings}.  This contradicts
the fact that the execution when $s$ is inactive has the same prefix before
the probe of $s$.  Therefore every adaptive policy has expected cost at
least~$L_1$.
\end{proof}

\subsection{The adaptive optimum}\label{sec:three-chain-adaptive-optimum}

\begin{theorem}[Cost of the three-chain instance]\label{thm:three-chain-cost}
For every integer $m\ge3$, the instance $\mathcal{I}_m^{\mathrm{TC}}$ in
\cref{con:three-chain} satisfies
\[
  \optpost(\mathcal{I}_m^{\mathrm{TC}})=18m+4,
  \qquad
  \optadapt(\mathcal{I}_m^{\mathrm{TC}})=20m+4.
\]
Therefore
\[
  \frac{\optadapt(\mathcal{I}_m^{\mathrm{TC}})}
       {\optpost(\mathcal{I}_m^{\mathrm{TC}})}
  =\frac{20m+4}{18m+4}
  =\frac{10m+2}{9m+2},
\]
and the ratio converges to $10/9$ as $m\to\infty$.  Consequently, the
clairvoyance gap is at least~$10/9$.
\end{theorem}

\begin{proof}
By \cref{lem:inactive-low-cost,lem:active-low-cost}, the two posterior optimum
costs are $L_0$ and $L_1$.  Since $p_s=1/2$,
\[
  \optpost(\mathcal{I}_m^{\mathrm{TC}})
  =\frac{L_0+L_1}{2}=18m+4.
\]

For an adaptive upper bound, follow the order of $T_1$ and probe $s$ upon
reaching $c_m$.  If $s$ is active, the realized tour is exactly $T_1$.  If it
is inactive, probe $b_0$ next.  The metric distance satisfies
\[
  d(c_m,b_0)\le d(c_m,s)+d(s,b_0)=2P,
\]
which is exactly the cost of the two $P$-edges used by the active branch.
Thus both realized costs are at most $L_1$, and
$\optadapt(\mathcal{I}_m^{\mathrm{TC}})\le L_1$.

The reverse inequality follows from \cref{lem:adaptive-lower-bound}.  Hence
$\optadapt(\mathcal{I}_m^{\mathrm{TC}})=L_1$.
\end{proof}

\subsection{Why these edge lengths?}\label{sec:three-chain-edge-lengths}

The choices of $U$ and $P$ are dictated by the two cost thresholds used above.
First consider a three-crossing walk that uses the wrong pair of chain
endpoints.  Relative to $T_0$, at least two chains
must each be traversed to an endpoint and back.  Its internal chain
contribution therefore increases from $3H$ to at least $5H$.  Choosing
\[
  P=2H
\]
makes this extra chain cost equal to~$P$.  Thus every such wrong-port walk has
cost at least~$L_1$, exactly the threshold needed in
\cref{lem:inactive-low-cost}.

The remaining critical competitor when $s$ is inactive uses all four
$U$-edges.  By \cref{lem:chain-contribution}(iii), its cost is at least
\[
  4H-4+4U.
\]
We need this value to be at least the cost of $T_1$,
\[
  L_1=3H+2U+2P.
\]
To make the $U$-edges as short as possible without allowing this competitor
below~$L_1$, we impose equality:
\[
  4H-4+4U=3H+2U+2P.
\]
Together with $P=2H$, this gives
\[
  U=\frac{3H}{2}+2,
\]
which is exactly $U=3m+2$ when $H=2m$.

If $U$ were smaller, the route using all four $U$-edges would have cost below the
threshold~$L_1$, so the preceding property of low-cost tours would no
longer hold.  If $U$ were larger, both posterior tour costs would increase
while their difference~$P$ stayed fixed, moving the ratio closer
to~$1$.  Hence the chosen values are at the balance for this construction.

\section{Upper bound with one semi-active vertex}
\label{sec:one-stochastic-upper-bound}

For $p\in[0,1]$, let $\Gammaone(p)$ denote the worst-case clairvoyance gap
among symmetric instances with exactly one semi-active vertex whose activation
probability is~$p$.  The subscript~$1$ refers to this single semi-active
vertex.  Thus,
\[
  \Gammaone(p)
  :=\sup_{\mathcal I}
  \frac{\optadapt(\mathcal I)}{\optpost(\mathcal I)},
\]
where the supremum ranges over all such instances, with any number of
always-active non-root vertices.

We now consider an arbitrary instance with exactly one semi-active vertex $s$,
active with probability $p\in(0,1)$.  Let $D$ be the
set of always-active non-root vertices.  If $D=\varnothing$, then probing
$s$ at the root is posterior-optimal whether $s$ is active or inactive, so the
clairvoyance gap is $1$.  We therefore assume $D\ne\varnothing$.  Let $T_0$ be
an optimal tour on
$\{r\}\cup D$, and let $T_1$ be an optimal tour on
$\{r\}\cup D\cup\{s\}$.  Write
\[
  \ell_0:=\cost(T_0),
  \qquad
  \ell_1:=\cost(T_1),
  \qquad
  \Delta:=\ell_1-\ell_0.
\]
Shortcutting $s$ from $T_1$ shows that $\ell_0\le\ell_1$, so $\Delta\ge0$.  The
posterior optimum is
\begin{equation}\label{eq:post-general}
  \optpost=(1-p)\ell_0+p\ell_1=\ell_0+p\Delta.
\end{equation}

We need two local exchange quantities.

First, define the cheapest insertion surcharge of $s$ into $T_0$ by
\begin{equation}\label{eq:alpha}
  \alpha
  :=\min_{uv\in E(T_0)}
  \bigl(d(u,s)+d(s,v)-d(u,v)\bigr).
\end{equation}
Inserting $s$ into the minimizing edge gives a full tour of cost $\ell_0+\alpha$,
so
\begin{equation}\label{eq:delta-alpha}
  \Delta\le\alpha.
\end{equation}

Second, let $x,y$ be the two neighbors of $s$ in $T_1$, and define the saving
from shortcutting $s$ by
\begin{equation}\label{eq:beta}
  \beta:=d(x,s)+d(s,y)-d(x,y).
\end{equation}
The triangle inequality gives $\beta\ge0$, and deleting $s$ from $T_1$ gives
a tour on the always-active vertices of cost $\ell_1-\beta$.  Hence
\begin{equation}\label{eq:beta-delta}
  0\le\beta\le\Delta.
\end{equation}

The key point is that the cheapest insertion surcharge into $T_0$ is
controlled by the deletion saving in $T_1$.

\begin{lemma}\label{lem:exchange}
The quantities in \cref{eq:alpha,eq:beta} satisfy
\[
  \alpha\le\frac{\ell_0}{2}+\beta.
\]
\end{lemma}

\begin{proof}
The vertices $x,y$ split the cycle $T_0$ into two $x$--$y$ arcs.  Let $P$ be
the shorter arc.  Then
\[
  \cost(P)\le\frac{\ell_0}{2},
  \qquad
  d(x,y)\le\cost(P)\le\frac{\ell_0}{2}.
\]
By the definition of $\beta$,
\[
  d(x,s)+d(s,y)=d(x,y)+\beta\le\frac{\ell_0}{2}+\beta.
\]
Therefore one of $x,y$, say $z$, satisfies
\[
  d(z,s)\le\frac{\ell_0}{4}+\frac\beta2.
\]
Let $z'$ be either neighbor of $z$ in $T_0$.  Inserting $s$ into the edge
$zz'$ costs
\[
\begin{aligned}
  d(z,s)+d(s,z')-d(z,z')
  &\le d(z,s)+\bigl(d(s,z)+d(z,z')\bigr)-d(z,z')\\
  &=2d(z,s)\\
  &\le\frac{\ell_0}{2}+\beta.
\end{aligned}
\]
Since $\alpha$ is the minimum insertion surcharge, the claim follows.
\end{proof}

\subsection{Two adaptive policies}\label{sec:one-stochastic-policies}

We translate the geometric quantities above into an adaptive upper bound by
considering two complementary policies.  One follows the optimal tour on the
always-active vertices and inserts $s$ when it is active; the other follows
the active-realization tour and shortcuts $s$ when it is inactive.  Taking the
better of the two gives the following estimate.

\begin{lemma}\label{lem:two-policies}
For every instance with one semi-active vertex,
\begin{equation}\label{eq:two-policy-bound}
  \optadapt
  \le
  \min\left\{
    \ell_0+p\left(\frac{\ell_0}{2}+\beta\right),
    \;\ell_0+\Delta-(1-p)\beta
  \right\}.
\end{equation}
\end{lemma}

\begin{proof}
For the first policy, orient $T_0$ and follow it until reaching one endpoint
$u$ of an edge $uv$ attaining the minimum in \cref{eq:alpha}.  Probe $s$ at
$u$.  If $s$ is inactive, continue directly to $v$ and complete $T_0$, at cost
$\ell_0$.  If $s$ is active, traverse $u,s,v$ and then complete $T_0$, at cost
$\ell_0+\alpha$.  Thus the expected cost is
\[
  \ell_0+p\alpha
  \le \ell_0+p\left(\frac{\ell_0}{2}+\beta\right)
\]
by \cref{lem:exchange}.

For the second policy, orient $T_1$ so that the local order is
$x,s,y$.  Follow $T_1$ to $x$ and probe $s$.  If $s$ is active, complete
$T_1$, at cost $\ell_1$.  If $s$ is inactive, move directly from $x$ to $y$ and
complete the shortcut tour, at cost $\ell_1-\beta$.  Its expected cost is
\[
  p\ell_1+(1-p)(\ell_1-\beta)
  =\ell_1-(1-p)\beta
  =\ell_0+\Delta-(1-p)\beta.
\]
Taking the better policy gives \cref{eq:two-policy-bound}.
\end{proof}

\subsection{Optimization of the bound}\label{sec:one-stochastic-optimization}

It remains to optimize the two-policy estimate over the insertion and
shortcut savings permitted by the metric.  After normalizing these quantities
by the cost of the tour on the always-active vertices, the problem reduces to
comparing two affine bounds in one variable.  This yields the general ratio
below.

\begin{theorem}[One semi-active vertex]\label{thm:general-upper}
For every instance with exactly one semi-active vertex of activation
probability $p\in[0,1]$,
\[
  \frac{\optadapt}{\optpost}
  \le\frac{2+p}{2+p^2}.
\]
Equivalently,
\[
  \Gammaone(p)\le\frac{2+p}{2+p^2}.
\]
Moreover, the right-hand side is maximized at $p=\sqrt6-2$, and hence
\[
  \Gammaone(p)\le\frac{2+\sqrt6}{4}
  \qquad\text{for every }p\in[0,1].
\]
\end{theorem}

\begin{proof}
The cases $p=0$ and $p=1$ have no information gap.  When the common
optimum is positive, the ratio is $1$; when it is zero, our convention gives
ratio $0$.  Thus the claimed bound holds, so assume $p\in(0,1)$.  Because
$D\ne\varnothing$ and $d$ is a metric, $\ell_0>0$.  Normalize
\[
  \delta:=\frac{\Delta}{\ell_0},
  \qquad
  b:=\frac{\beta}{\ell_0}.
\]
By \cref{eq:post-general,eq:two-policy-bound},
\begin{equation}\label{eq:normalized-bound}
  \frac{\optadapt}{\ell_0}
  \le
  \min\left\{
    1+\frac p2+pb,
    \;1+\delta-(1-p)b
  \right\},
  \qquad
  \frac{\optpost}{\ell_0}=1+p\delta.
\end{equation}

Suppose first that $\delta\le p/2$.  The second policy and $b\ge0$ give
\[
  \frac{\optadapt}{\ell_0}\le1+\delta.
\]
The function $(1+\delta)/(1+p\delta)$ is increasing in $\delta$ because
$p<1$.  Hence
\[
  \frac{\optadapt}{\optpost}
  \le\frac{1+\delta}{1+p\delta}
  \le\frac{1+p/2}{1+p^2/2}
  =\frac{2+p}{2+p^2}.
\]

Now suppose that $\delta\ge p/2$ and set
\[
  b_0:=\delta-\frac p2\ge0.
\]
If $b\le b_0$, the first term in \cref{eq:normalized-bound} gives
\[
\begin{aligned}
  \frac{\optadapt}{\ell_0}
  &\le1+\frac p2+pb\\
  &\le1+\frac p2+p\left(\delta-\frac p2\right)\\
  &=1+p\delta+\frac{p(1-p)}2.
\end{aligned}
\]
If $b\ge b_0$, the second term gives
\[
\begin{aligned}
  \frac{\optadapt}{\ell_0}
  &\le1+\delta-(1-p)b\\
  &\le1+\delta-(1-p)\left(\delta-\frac p2\right)\\
  &=1+p\delta+\frac{p(1-p)}2.
\end{aligned}
\]
Thus in either case,
\[
  \frac{\optadapt}{\optpost}
  \le
  1+\frac{p(1-p)/2}{1+p\delta}.
\]
Since $\delta\ge p/2$,
\[
  \frac{\optadapt}{\optpost}
  \le
  1+\frac{p(1-p)/2}{1+p^2/2}
  =\frac{1+p/2}{1+p^2/2}
  =\frac{2+p}{2+p^2}.
\]
This proves the probability-dependent bound.

Finally, for
\[
  f(p):=\frac{2+p}{2+p^2},
\]
the sign of
\[
  f'(p)=\frac{2-4p-p^2}{(2+p^2)^2}
\]
changes from positive to negative at the unique point
$p=\sqrt6-2$ in~$[0,1]$.  Therefore $f$ is maximized there, with
\[
  f(\sqrt6-2)=\frac{2+\sqrt6}{4}.
\]
\end{proof}

At $p=1/2$, the bound in \cref{thm:general-upper} equals~$10/9$, and
\cref{thm:three-chain-cost} shows that it is tight within the
one-semi-active-vertex class.

% ============================================================
% CHAPTER 4: CLAIRVOYANCE GAP — ASYMMETRIC
% ============================================================
\chapter{Clairvoyance Gap for Asymmetric Metrics}\label{ch:asym-clair}

We now turn to asymmetric (directed) metrics and present two constructions.
The first is a small subdivided-triangle example and its natural recursive
generalization.  At each recursive level, a clairvoyant tour chooses the
cheaper orientation after seeing whether a new midpoint is active, whereas an
adaptive policy must make part of this choice before learning that activation
bit.  This family makes the local order conflict especially clear, but its
ratio tends only to~$4/3$ and its full analysis is more complicated than that
of the stronger construction, so we do not expand the full analysis here.

Our main construction proves a clairvoyance-gap lower bound of~$3/2$ by
embedding a layered poset into an asymmetric metric.  A clairvoyant tour needs
essentially two increasing chains, while every adaptive policy needs at least
three increasing runs with high probability.  The analysis allows a policy to
interleave probes across different recursive blocks or poset layers.

\section{The recursive-triangle construction}
\label{sec:recursive-construction}

\begin{remark}[From one subdivided arc to a recursive family]
\label{rem:recursive-triangle}
A tiny instance in \cref{fig:recursive-triangle-idea} is constructed in this way:
start from the bidirected unit triangle on
$r=a,b,c$, keep the arc $b\to c$, and subdivide the reverse arc by a new
vertex~$d$. So in summary we have edges:
\[
  a\leftrightarrow b,\qquad
  a\leftrightarrow c,\qquad
  b\to c, \qquad
  c\to d\to b.
\]
The vertices $a,b,c$ are always active, while $d$ is active with probability~$1/2$.

If $d$ is inactive, the posterior tour $a\to b\to c\to a$ costs~$3$.
If $d$ is active, the posterior tour
$a\to c\to d\to b\to a$ costs~$4$.  Hence
\[
  \optpost=\frac12\cdot3+\frac12\cdot4=\frac72.
\]
Clearly, we have conflicting optimal orders when $d$ is active and when $d$ is inactive.
A one-step Bellman calculation over the three possible first probes
$b,c,d$ gives $\optadapt=4$, attained for example by probing in the order $c,d,b$, and the
clairvoyance ratio is
\[
  \frac{\optadapt}{\optpost}=\frac{4}{7/2}=\frac87.
\]
More generally, if $d$ is active with probability~$p$, then
$\optpost=3+p$ and $\optadapt=\min\{3+2p,4\}$; the ratio is maximized at
$p=1/2$.

The recursive family in \cref{fig:recursive-triangle-idea} is the direct
generalization of this example.  At level~$k$, the vertices $b$ and~$c$ are
replaced by two copies $B_{k-1}^{\mathrm L}$ and
$B_{k-1}^{\mathrm R}$, the vertex $d$ becomes a new semi-active midpoint
$m_k$, and each displayed unit arc becomes a complete bundle of port arcs of
length $w_k=2^{k-1}$.  The leaves are always active and all midpoints
are independently active with probability~$1/2$.  With root arcs of length
$w_l$ at depth~$l$, the exact ratio is
\[
  \frac{\optadapt(\mathcal I_l^{\mathrm{REC}})}
       {\optpost(\mathcal I_l^{\mathrm{REC}})}
  =\frac{4(l+1)}{3l+4}
  \longrightarrow\frac43.
\]
The depth-one member is precisely the four-vertex instance above.
\end{remark}

\begin{figure}[htbp]
\centering
\begin{minipage}[c]{0.39\textwidth}
\centering
\begin{tikzpicture}[
    >=Stealth,
    always/.style={rectangle,draw=black,fill=white,minimum size=7mm,
      inner sep=1pt,font=\small},
    semi/.style={circle,draw=black,fill=gray!25,minimum size=7mm,
      inner sep=0pt,font=\small},
    arc/.style={->,line width=0.8pt}
]
  \node[always] (A) at (0,2.1) {$r=a$};
  \node[always] (B) at (-1.7,0) {$b$};
  \node[always] (C) at (1.7,0) {$c$};
  \node[semi]   (D) at (0,-1.65) {$d$};

  \draw[arc] (A) to[bend right=10] (B);
  \draw[arc] (B) to[bend right=10] (A);
  \draw[arc] (A) to[bend left=10] (C);
  \draw[arc] (C) to[bend left=10] (A);
  \draw[arc] (B) -- (C);
  \draw[arc] (C) -- (D);
  \draw[arc] (D) -- (B);
\end{tikzpicture}

\smallskip
\textbf{(a)} The four-vertex seed
\end{minipage}
\hfill
\begin{minipage}[c]{0.57\textwidth}
\centering
\begin{tikzpicture}[
    >=Stealth,
    block/.style={draw,dashed,rounded corners,minimum width=2.45cm,
      minimum height=1.15cm,align=center,font=\small},
    root/.style={rectangle,draw=black,fill=white,minimum size=7mm,
      inner sep=1pt,font=\small},
    semi/.style={circle,draw=black,fill=gray!25,minimum size=7mm,
      inner sep=0pt,font=\small},
    arc/.style={->,line width=0.8pt}
]
  \node[root]  (R0) at (0,2.15) {$r$};
  \node[block] (L)  at (-2.0,0)
    {$B_{k-1}^{\mathrm L}$\\[-0.15em]\scriptsize always-active ports};
  \node[block] (R)  at (2.0,0)
    {$B_{k-1}^{\mathrm R}$\\[-0.15em]\scriptsize always-active ports};
  \node[semi]  (M)  at (0,-1.75) {$m_k$};

  \draw[arc] (L) -- node[above,font=\scriptsize] {$w_k$} (R);
  \draw[arc] (R.south) to[bend left=9]
    node[right,font=\scriptsize] {$w_k$} (M.east);
  \draw[arc] (M.west) to[bend left=9]
    node[left,font=\scriptsize] {$w_k$} (L.south);
  \draw[arc] (R0) to[bend right=12]
    node[left,font=\scriptsize] {$w_l$} (L.north);
  \draw[arc] (L.north) to[bend right=12] (R0);
  \draw[arc] (R0) to[bend left=12]
    node[right,font=\scriptsize] {$w_l$} (R.north);
  \draw[arc] (R.north) to[bend left=12] (R0);
\end{tikzpicture}

\smallskip
\textbf{(b)} Recursive substitution
\end{minipage}
\caption{The recursive-triangle idea.  All displayed arcs in panel~(a) have
unit length.  Panel~(b) replaces $b,c,d$ by
$B_{k-1}^{\mathrm L},B_{k-1}^{\mathrm R},m_k$; an arrow incident with a
child denotes the complete bundle of arcs incident with its always-active
ports.
The root arcs occur only at the top level $k=l$.}
\label{fig:recursive-triangle-idea}
\end{figure}

\section{A layered-poset lower bound}
\label{sec:layered-poset-gap}

The construction is best viewed as routing two increasing chains through a
sequence of stages.  At one stage, a chain enters through one of two gates
and leaves through one of two gates, so its transition is a cell of a
$2\times2$ matrix.  Four semi-active selectors represent the two rows and
the two columns of this matrix.  After seeing the realization, any two active
selectors can be assigned to the two complementary cells.  An adaptive
policy, however, may have to choose the first cell before it knows which cell
the second chain will need.

The posterior solution needs an
extra chain only when at least three selectors in a stage are active, an
event of order~$p^3$, whereas an exact active pair, an event of order~$p^2$,
can invalidate an adaptive routing decision.  We therefore want $Lp^2$ to be
large but $Lp^3$ to be small.  The asymmetric metric below makes its tour
cost essentially equal to the number of increasing runs, turning this local
conflict into a clairvoyance gap.

\subsection{The construction}\label{sec:four-star-layered-poset}

\begin{construction}\label{con:layered-poset}
Fix an integer $L\ge1$ and parameters $p,\varepsilon\in(0,1)$.  For each
$0\le i\le L$, introduce two \emph{gate vertices}
\[
    G_i:=\{g_i^1,g_i^2\},
\]
which are always active.  For every stage $1\le i\le L$, introduce four
\emph{selector vertices} $A_i,B_i,C_i,D_i$, each independently active with
probability~$p$.

Define a poset~$P_L$ by the following generating comparabilities:
\begin{equation}\label{eq:four-star-comparabilities}
\begin{array}{c|c|c}
\text{selector}
    &\text{predecessor gates}
    &\text{successor gates}\\ \hline
A_i&\{g_{i-1}^1\}&\{g_i^1,g_i^2\}\\
B_i&\{g_{i-1}^2\}&\{g_i^1,g_i^2\}\\
C_i&\{g_{i-1}^1,g_{i-1}^2\}&\{g_i^1\}\\
D_i&\{g_{i-1}^1,g_{i-1}^2\}&\{g_i^2\}.
\end{array}
\end{equation}
Every gate in the middle column is below the selector, and the selector is
below every gate in the last column.  Take the transitive closure and add no
other comparabilities.

Add a root~$r$.  For distinct vertices $x,y$ of~$P_L$, set
\begin{equation}\label{eq:layered-poset-metric}
    d_\varepsilon(x,y)
    :=
    \begin{cases}
        \varepsilon,&x<y\text{ in }P_L,\\
        1,&x\not<y\text{ in }P_L,
    \end{cases}
    \qquad
    d_\varepsilon(r,x)=d_\varepsilon(x,r):=\frac12,
\end{equation}
and put $d_\varepsilon(x,x)=0$.  Denote the resulting \StochTSP{} instance
by $\mathcal{I}^{\mathrm{LP}}(L,p,\varepsilon)$.
\end{construction}

To read one stage, identify a transition from $G_{i-1}$ to~$G_i$ with a
cell~$(r,c)$: the chain enters through $g_{i-1}^r$ and leaves through
$g_i^c$.  The allowed cells of the four selectors are
\begin{equation}\label{eq:four-star-matrix}
\begin{array}{c|cc}
&c=1&c=2\\ \hline
r=1&A_i,C_i&A_i,D_i\\
r=2&B_i,C_i&B_i,D_i.
\end{array}
\end{equation}
Thus $A_i,B_i$ are the two row selectors and $C_i,D_i$ are the two column
selectors.  The four selectors at one stage are pairwise incomparable, as
are the two gates in one layer, while every gate in $G_{i-1}$ is below every
gate in~$G_i$.

\begin{figure}[!htbp]
\centering
\begin{minipage}[c]{0.54\textwidth}
\centering
\begin{tikzpicture}[
    >=Stealth,
    gate/.style={rectangle,draw=black,fill=white,minimum width=8mm,
      minimum height=6.5mm,inner sep=1pt,font=\scriptsize},
    selector/.style={circle,draw=black,fill=gray!20,minimum size=7mm,
      inner sep=0pt,font=\scriptsize},
    order/.style={->,black!65,line width=0.65pt}
]
  \node[gate] (L1) at (-1.5,0) {$g_{i-1}^1$};
  \node[gate] (L2) at ( 1.5,0) {$g_{i-1}^2$};
  \node[selector] (A) at (-2.55,1.7) {$A_i$};
  \node[selector] (C) at (-0.85,1.7) {$C_i$};
  \node[selector] (D) at ( 0.85,1.7) {$D_i$};
  \node[selector] (B) at ( 2.55,1.7) {$B_i$};
  \node[gate] (R1) at (-1.5,3.4) {$g_i^1$};
  \node[gate] (R2) at ( 1.5,3.4) {$g_i^2$};

  \draw[order] (L1)--(A);
  \draw[order] (L1)--(C);
  \draw[order] (L1)--(D);
  \draw[order] (L2)--(B);
  \draw[order] (L2)--(C);
  \draw[order] (L2)--(D);
  \draw[order] (A)--(R1);
  \draw[order] (A)--(R2);
  \draw[order] (B)--(R1);
  \draw[order] (B)--(R2);
  \draw[order] (C)--(R1);
  \draw[order] (D)--(R2);
\end{tikzpicture}

\smallskip
\textbf{(a)} One layered-poset stage
\end{minipage}
\hfill
\begin{minipage}[c]{0.42\textwidth}
\centering
\begin{tikzpicture}[font=\scriptsize]
  \fill[gray!25] (0,1) rectangle (1.65,2);
  \draw[line width=1pt] (0,1) rectangle (1.65,2);
  \draw[densely dashed,line width=1pt] (1.65,0) rectangle (3.3,1);
  \draw (0,0) rectangle (3.3,2);
  \draw (1.65,0)--(1.65,2);
  \draw (0,1)--(3.3,1);
  \node at (0.825,1.5) {$\mathbf{A_i,C_i}$};
  \node at (2.475,1.5) {$A_i,D_i$};
  \node at (0.825,0.5) {$B_i,C_i$};
  \node at (2.475,0.5) {$B_i,D_i$};
  \node[above] at (0.825,2) {$g_i^1$};
  \node[above] at (2.475,2) {$g_i^2$};
  \node[left,align=right] at (0,1.5) {$g_{i-1}^1$};
  \node[left,align=right] at (0,0.5) {$g_{i-1}^2$};
  \node[align=center] at (1.65,-0.52)
    {solid: first transition\\dashed: complementary transition};
\end{tikzpicture}

\smallskip
\textbf{(b)} A premature commitment
\end{minipage}
\caption{One stage of the layered-poset construction.  Arrows in panel~(a)
point upward in the poset.  In panel~(b), a first chain entering through
$g_{i-1}^1$ assigns $A_i$ to cell~$(1,1)$, forcing the second chain to use
cell~$(2,2)$.  If exactly $A_i,C_i$ are active, this commitment fails,
although a posterior solution can assign $A_i$ to~$(1,2)$ and $C_i$
to~$(2,1)$.}
\label{fig:four-star-stage}
\end{figure}

The example in \cref{fig:four-star-stage} panel (b) contains the entire local obstruction.  If the
first chain enters in row~$1$, sees $A_i$ active, and leaves in column~$1$,
then a later active $C_i$ cannot use the complementary cell~$(2,2)$.  Had the
policy instead left in column~$2$, the exact pair $A_i,D_i$ would create the
same problem.  A posterior solution avoids either conflict because it chooses
both cells after seeing the pair.

\subsection{From increasing runs to tour cost}

The metric in \cref{con:layered-poset} makes the poset interpretation exact.
Forward moves within an increasing run are cheap, whereas every break between
runs costs one unit.

\begin{lemma}
\label{lem:layered-poset-metric}
The function $d_\varepsilon$ is a strictly positive directed metric.  If a
realization has $N$ active vertices and a service order has $K$ maximal
increasing runs, then its closed-tour cost is exactly
\begin{equation}\label{eq:layered-poset-run-cost}
    \varepsilon N+(1-\varepsilon)K.
\end{equation}
In particular,
\begin{equation}\label{eq:layered-poset-posterior-identity}
    \opt(A)
    =
    \varepsilon N(A)
    +(1-\varepsilon)\operatorname{width}(P_L[A]).
\end{equation}
\end{lemma}

\begin{proof}
For the triangle inequality, consider first three non-root vertices.  If two
consecutive arcs have length~$\varepsilon$, then $x<y<z$, so transitivity
gives $x<z$ and the direct arc also has length~$\varepsilon$.  If a direct
arc has length~$1$, a two-arc path through a distinct non-root vertex cannot
consist of two $\varepsilon$-arcs and therefore has length at least
$1+\varepsilon$.  A detour through the root has length~$1$.  The triangle
inequalities involving a root arc follow in the same way.  Thus all directed
triangle inequalities hold, and all distances between distinct vertices are
positive.

The two root legs contribute~$1$.  Within the service order, the $N-K$
transitions inside increasing runs cost~$\varepsilon$, and the $K-1$ breaks
between runs cost~$1$.  Hence the total is
\[
    1+\varepsilon(N-K)+(K-1)
    =\varepsilon N+(1-\varepsilon)K.
\]
By Dilworth's theorem~\cite{Dilworth1950}, the minimum number of increasing chains covering a
finite poset equals its width.  Concatenating the chains gives a service order
with that many runs, while the runs of every service order form a chain
cover.  Minimizing the preceding expression proves
\cref{eq:layered-poset-posterior-identity}.
\end{proof}

The remainder of the proof compares these two quantities.  After seeing the
realization, two chains suffice except for the excess selectors in stages
with at least three active selectors, yielding expected width
$2+O(Lp^3)$.  For an adaptive policy, repeated exact-pair conflicts make the
probability of using only two runs at most
$(1-p^2(1-p)^2)^L$.

\subsection{The posterior chain cover}\label{sec:posterior-chain-cover}

We first evaluate what a posteriori solution can achieve after all active
selectors are known.  An explicit chain cover shows that the width is two
unless some stage contains at least three active selectors, so the expected
excess width is governed by a binomial tail of order~$p^3$.

Let $Z_i$ be the number of active selectors in stage~$i$, so
$Z_i\sim\operatorname{Binomial}(4,p)$.  For a realization~$A$, write
$P_L[A]$ for the subposet induced by its active vertices.

\begin{lemma}
\label{lem:layered-poset-width}
Every realization satisfies
\[
    2
    \le\operatorname{width}(P_L[A])
    \le2+\sum_{i=1}^L(Z_i-2)_+.
\]
Consequently,
\[
    2
    \le\E[\operatorname{width}(P_L[A])]
    \le2+L(4p^3-2p^4)
    \le2+4Lp^3.
\]
\end{lemma}

\begin{proof}
Every gate layer is an always-active two-element antichain, which proves the lower
bound.

For the upper bound, consider a stage with exactly two active selectors.
There are only six possible pairs.  The following table assigns each selector
an allowed cell from \cref{eq:four-star-matrix}, making sure the two assigned cells
are in different row and different column. For each assignment, we
write down the cell~$(r,c)$ and the corresponding chain transition
$g_{i-1}^r\to X_i\to g_i^c$.  Each case uses both gates in~$G_{i-1}$ and both
gates in~$G_i$, so the two gates are all covered and two chains continue to the next stage.
\[
\begin{array}{c|cc}
\text{active pair}
    &\text{transition from }g_{i-1}^1
    &\text{transition from }g_{i-1}^2\\ \hline
\{A_i,B_i\}&(1,1):\ g_{i-1}^1\to A_i\to g_i^1&(2,2):\ g_{i-1}^2\to B_i\to g_i^2\\
\{C_i,D_i\}&(1,1):\ g_{i-1}^1\to C_i\to g_i^1&(2,2):\ g_{i-1}^2\to D_i\to g_i^2\\
\{A_i,C_i\}&(1,2):\ g_{i-1}^1\to A_i\to g_i^2&(2,1):\ g_{i-1}^2\to C_i\to g_i^1\\
\{A_i,D_i\}&(1,1):\ g_{i-1}^1\to A_i\to g_i^1&(2,2):\ g_{i-1}^2\to D_i\to g_i^2\\
\{B_i,C_i\}&(1,1):\ g_{i-1}^1\to C_i\to g_i^1&(2,2):\ g_{i-1}^2\to B_i\to g_i^2\\
\{B_i,D_i\}&(1,2):\ g_{i-1}^1\to D_i\to g_i^2&(2,1):\ g_{i-1}^2\to B_i\to g_i^1.
\end{array}
\]
If at most one selector is active, simply use any one of its allowed cells.

Start two chains at $g_0^1$ and~$g_0^2$.  At stage~$i$, if $Z_i\le2$, put all
active selectors on the chain transitions as shown by the cell
assignment above.  If $Z_i>2$, put any two active selectors on those transitions
and give each remaining selector a singleton chain. This gives a chain cover of size at most
\[
    2+\sum_{i=1}^L(Z_i-2)_+.
\]
Finally,
\begin{align*}
    \E[(Z_i-2)_+]
    &=\P(Z_i=3)+2\P(Z_i=4)\\
    &=4p^3(1-p)+2p^4\\
    &=4p^3-2p^4,
\end{align*}
and Dilworth's theorem proves the remaining claims.
\end{proof}

\subsection{The adaptive obstruction}\label{sec:adaptive-obstruction}

We next show why an adaptive policy cannot reproduce the posterior chain
cover.  Writing the service order as increasing runs turns a two-run execution
into a sequence of local routing decisions.  At each untouched stage, a probe
can force one such decision before the policy knows which compatible routing
will be required.

Fix an adaptive policy and one realization.  List the active vertices in the
order in which the policy serves them.  An \emph{increasing run} is a maximal
contiguous part of this list in which every vertex is below the next vertex
in~$P_L$.  Let $K_\pi$ be the number of such runs.

Every gate layer forces $K_\pi\ge2$.  Moreover, on the event $K_\pi\le2$,
the two runs form a two-chain cover of all always-active gates.  Each run
therefore contains exactly one gate from every layer and passes through the
layers in increasing order.  In each stage the first run can contain at most
one selector.  If it enters the stage in row~$r$ and leaves in column~$c$,
any other active selector must fit the complementary cell $(3-r,3-c)$ on the
second run.

This observation also explains why remote probes do not evade the
construction.  If a policy probes a selector in a future stage before the
first run has served the intervening gates, an active answer jumps past one
gate in each intervening layer.  The second run cannot contain both gates of
such a layer, so $K_\pi\le2$ is then impossible.  Remote probes with inactive outcomes cause
no movement and will be allowed in the product argument below.

Put
\[
    h:=p^2(1-p)^2.
\]

\begin{lemma}
\label{lem:four-star-local}
Fix the input row of the first run at one stage.  Every deterministic adaptive
strategy fails to preserve a two-run routing on at least one realization
having exactly two active selectors.  Its one-stage success probability is
therefore at most~$1-h$.  The same statement holds if the strategy is given
an arbitrary advice string independent of the four local activation bits.
\end{lemma}

\begin{proof}
Follow the strategy on the branch on which every selector probed before the
output gate is inactive.  If the strategy chooses an output gate without
probing any selector, every exact two-active realization is fatal.  Probing
one of the two selectors before the second run reaches its predecessor gate
skips an unserved gate and forces a third run; waiting leaves two incomparable
selectors for the single second run.

Otherwise, let $X$ be the first probed selector.  Follow the branch on which
$X$ is active and every later selector probed before the output gate is
inactive.  If $X$ is incompatible with the input row, this branch already
cannot belong to a two-run routing, and any exact pair containing~$X$ is
fatal.  We may therefore assume that the strategy assigns~$X$ to a cell
$(r,c)$ compatible with its input row and chosen output column.

Exactly two of the four selectors allow cell~$(r,c)$: the row-$r$ selector
and the column-$c$ selector.  Let $Y\ne X$ be the other one.  Consider the realization in which
exactly $X$ and~$Y$ are active.  If the strategy probes~$Y$ before taking
the output gate, then it sees two incomparable active selectors before the
first run leaves the stage, so the two-run event has already failed.  Suppose
instead that it probes~$Y$ later.  If the second run has not yet reached its
predecessor gate in stage~$i$, probing the active vertex~$Y$ jumps past an unserved
gate of that run and forces a third run.  Once the second run has reached
that predecessor, it must use the complementary cell $(3-r,3-c)$, but neither
selector that allows $(r,c)$ allows this cell.  Probing~$Y$ then
or after leaving the stage also forces a third run.  Since every selector
must eventually be probed, the exact $\{X,Y\}$ realization is fatal in all
cases.

Thus one exact-pair atom is always excluded from the success event.  Every
such atom has probability $p^2(1-p)^2=h$.  Conditioning on independent advice
merely selects a deterministic local strategy, so it does not change the
argument.
\end{proof}

\paragraph{Handling interleaved probes.}
The policy need not finish one stage before probing another.  The next lemma
shows that this freedom does not remove the local loss: the conditional
success probabilities can be multiplied even under arbitrary interleaving.

\begin{lemma}
\label{lem:four-star-product}
Every adaptive policy, including a randomized policy making arbitrary remote
probes, satisfies
\[
    \P(K_\pi\le2)\le(1-h)^L.
\]
Consequently,
\[
    \E[K_\pi]\ge3-(1-h)^L.
\]
\end{lemma}

\begin{proof}
We first fix a deterministic policy.  View the preservation of two runs as a
finite reachability problem with one independent four-bit component per
stage.  A component becomes \emph{ready} when the first run reaches its input
gate and remains unresolved until all its selectors have been probed and the
second-run routing, if needed, has been completed.  For a ready unresolved
stage, let $\Psi(s)$ be the optimal conditional probability of completing a
valid local routing from its current local state~$s$.  The state records the
input row, all local answers seen so far, any selector placed on the first
run, its output column, whether the first run has left the stage, and whether
the second run has reached its complementary input gate.  In particular, a
probe with an active selector outcome after the first-run output but before that complementary
input gate is reached is a failed state, by the argument in the proof of
\cref{lem:four-star-local}.  The Bellman definition gives
\begin{equation}\label{eq:four-star-bellman}
    \E[\Psi(s_{\mathrm{next}})\mid s]\le \Psi(s)
\end{equation}
for every local action.  At an untouched stage,
\cref{lem:four-star-local} gives $\Psi(s)\le1-h$, for either input row and
conditional on any observations outside the stage.

For a future stage that is not yet ready, assign an upper potential as
follows.  Its potential is $1-h$ before any premature local probe, it is~$1$
after one or more such probes have all answered inactive, and it is~$0$ if
one answers active.  The last case makes $K_\pi\le2$ impossible by the gate
argument preceding \cref{lem:four-star-local}.  The first premature probe
changes the expected potential by at most
\[
    p\cdot0+(1-p)\cdot1
    =1-p
    \le1-h,
\]
where $h\le p$.  Further premature probes start from potential~$1$ and have
expected next potential at most~$1$.

Take the product of the $L$ local potentials, using factor~$1$ only after all
local selectors have been resolved and a valid two-chain routing has been
completed, and factor~$0$ for a failed stage.  After the first run leaves a
stage, its exact Bellman factor remains in the product until the second-run
part of that stage is resolved.  A selector probe changes only one factor.
When the first run reaches a new input gate, an untouched factor becomes an
exact Bellman value at most~$1-h$, while a prematurely probed factor becomes
an exact conditional value at most~$1$.  Gate probes that violate the forced
layer order make the two-run event impossible and may also set the product
to zero.  Hence
\cref{eq:four-star-bellman} and the preceding premature-probe estimate make
the product a nonnegative supermartingale under arbitrary interleaving.

Its initial value is at most $(1-h)^L$.  At termination it is the indicator
that every stage admits the required local two-chain routing.  The event
$K_\pi\le2$ implies this terminal event, so taking expectations proves
\[
    \P(K_\pi\le2)\le(1-h)^L.
\]
For a randomized policy, condition on its private seed and average the same
bound.  Finally, $K_\pi\ge2$ implies
\[
    \E[K_\pi]
    \ge2+\P(K_\pi\ge3)
    =3-\P(K_\pi\le2)
    \ge3-(1-h)^L.
\]
\end{proof}

\begin{theorem}[Layered-poset lower bound]
\label{thm:layered-poset-gap}
The clairvoyance gap for asymmetric \StochTSP{} is at least~$3/2$.
More precisely, there is a sequence of \StochTSP{} instances
$(\mathcal{I}_t^{\mathrm{LP}})_{t\ge2}$ such that
\[
    \liminf_{t\to\infty}
    \frac{\optadapt(\mathcal{I}_t^{\mathrm{LP}})}
         {\optpost(\mathcal{I}_t^{\mathrm{LP}})}
    \ge\frac32.
\]
\end{theorem}

\begin{proof}
For an integer $t\ge2$, choose
\[
    L:=t^5,
    \qquad
    p:=t^{-2},
    \qquad
    \varepsilon:=t^{-10}=L^{-2},
\]
and define the resulting instance by
\[
    \mathcal{I}_t^{\mathrm{LP}}
    :=\mathcal{I}^{\mathrm{LP}}(L,p,\varepsilon).
\]
Its expected number of active vertices is
\[
    \E[N]=2(L+1)+4Lp=O(L),
\]
so $\varepsilon\E[N]=o(1)$.  \Cref{lem:layered-poset-width} and
$Lp^3=t^{-1}$ give
\[
    \optpost(\mathcal{I}_t^{\mathrm{LP}})
    \le
    \varepsilon\E[N]
    +(1-\varepsilon)(2+4Lp^3)
    =2+o(1).
\]

On the adaptive side,
\[
    Lh
    =Lp^2(1-p)^2
    =t(1-t^{-2})^2
    \longrightarrow\infty.
\]
Hence $(1-h)^L\le\e^{-Lh}=o(1)$.
\Cref{lem:four-star-product,lem:layered-poset-metric} imply, for every
policy~$\pi$,
\[
    \E[\cost(\pi)]
    \ge
    \varepsilon\E[N]
    +(1-\varepsilon)\bigl(3-(1-h)^L\bigr)
    =3-o(1).
\]
Taking the infimum over policies and combining the two estimates yields
\[
    \liminf_{t\to\infty}
    \frac{\optadapt(\mathcal{I}_t^{\mathrm{LP}})}
         {\optpost(\mathcal{I}_t^{\mathrm{LP}})}
    \ge\frac32.
\]
\end{proof}

% ============================================================
% CHAPTER 5: BACKBONE CIRCULATIONS AND ROUNDING
% ============================================================
\chapter{Backbone Circulations and Rounding for Asymmetric \StochTSP{}}
\label{ch:circulation}

This chapter proves the following upper bound on the value of clairvoyance in
asymmetric \StochTSP{}.

\begin{theorem}[Asymmetric clairvoyance-gap upper bound]
\label{thm:asymmetric-clairvoyance-upper}
There is a universal constant~$C$ such that every asymmetric \StochTSP{}
instance on $n$ non-root vertices admits an adaptive policy~$\pi$ with
\[
    \cost(\pi)\le C(1+n^{1/3})\optpost.
\]
Consequently, the clairvoyance gap for asymmetric metrics is
$O(n^{1/3})$.
\end{theorem}

The proof is existential; we do not claim that the policy can be computed in
polynomial time.  The bound improves the $O(\sqrt n)$ estimate inherited from
the asymmetric obliviousness-gap bound of
Christalla, Puhlmann, and Traub~\cite{Christalla2025}.

The starting point is the backbone method for symmetric a~priori
TSP~\cite{Shmoys2008}.  Sample clients according to their activation
probabilities, route the sampled clients, and attach the others to this
backbone.  In an asymmetric metric, a nearest-backbone detour is not adequate
because entering and leaving a client can have very different costs.  We
therefore encode the attachments by a directed circulation.

The proof has two stages.  We first develop the structural circulation and
endpoint-rounding framework.  It gives a low-cost fractional circulation and
shows how its outside-client paths can be rounded to directed chords while
controlling imbalance on every interval of the backbone order.  This also
yields an offline bounded-multiplicity walk and isolates the obstacle to an
adaptive execution.

We then remove that obstacle by changing what is rounded.  Publicly sample one
backbone chord for each outside client, correct the resulting weighted
endpoint divergence along the backbone cycle, and round the auxiliary flow to
an integer flow.  Its \emph{mean} imbalance is at most one on every cyclic
interval.  During execution, only the independent activation indicators
fluctuate; a one-dimensional maximal inequality bounds their maximum interval
imbalance.  An active-first stack walk implements the chords adaptively.  This
gives a low-mass bound of $O(\log n+\sqrt\mu)\optpost$, where~$\mu$ is the
expected number of active clients.  Finally, separating clients at activation
probability $n^{-1/3}$ balances the low-mass policy against a fixed tour for
the high-probability clients and proves
\cref{thm:asymmetric-clairvoyance-upper}.

\section{The backbone circulation relaxation}
\label{sec:circulation-lp}

We first formalize the fractional object that attaches vertices outside a
sampled backbone to a tour on that backbone.  Each outside vertex receives
the prescribed expected in- and out-degree, while forbidding arcs between two
outside vertices makes every such attachment pass through the backbone.

Independently include each non-root vertex $v\in V^\circ$ in a random set~$S$ with
probability~$p_v$.  We call
\[
    R:=S\cup\{r\}
    \qquad\text{and}\qquad
    U:=V^\circ\setminus S
\]
the \emph{backbone} and the set of \emph{outside vertices}, respectively.
For a nonnegative vector~$z$ on directed pairs, use the notation
\[
    z(\delta^+(w)):=\sum_v z_{wv},
    \qquad
    z(\delta^-(w)):=\sum_u z_{uw}.
\]
For a fixed choice of~$S$, consider the following circulation relaxation:
\begin{equation}\label{eq:backbone-circulation-lp}
\begin{aligned}
\min\quad
    &\sum_{u,v\in V}d(u,v)z_{uv}\\
\text{subject to}\quad
    &z(\delta^+(w))=z(\delta^-(w))
        &&\forall w\in V,\\
    &z(\delta^+(i))=z(\delta^-(i))=p_i
        &&\forall i\in U,\\
    &z_{ij}=0
        &&\forall i,j\in U,\\
    &z_{uv}\ge0
        &&\forall u,v\in V.
\end{aligned}
\tag{$\mathrm{CIRC}(S)$}
\end{equation}
For a feasible vector~$z$, write
\[
    \cost(z):=\sum_{u,v\in V}d(u,v)z_{uv},
\]
and let $c(S)$ denote the optimum value of
$\mathrm{CIRC}(S)$.  The prohibition on arcs between two
outside vertices makes the intended structure explicit: flow that visits an
outside vertex enters from the backbone and immediately returns to the
backbone.  Thus the relaxation is a fractional directed analogue of assigning
a detour to each outside vertex.

\section{A low-cost fractional circulation}
\label{sec:low-cost-circulation}

We next show that the backbone tour and the circulation can be chosen to have
small total cost.  The logarithmic factor comes from the longest consecutive
run of outside vertices on an optimal tour through two independent
realizations.

Let $A\subseteq V^\circ$ be an activation set sampled independently of~$S$,
using the same probabilities~$(p_v)_{v\in V^\circ}$, and put $C:=A\cup S$.
Fix an optimal tour~$T_C$ on $C\cup\{r\}$.  Cutting this tour at the vertices
of~$R$ produces directed segments
\[
    P_j=(s_j,u_{j,1},\ldots,u_{j,k_j},s_{j+1}),
\]
where $s_j,s_{j+1}\in R$ and every internal vertex belongs to
$A\setminus S$.  Indices on the vertices of~$R$ are cyclic.  Define
\[
    M(A,S):=\max_j k_j.
\]

For every segment~$P_j$, send one unit of flow along
$s_j\to u_{j,\ell}\to s_{j+1}$ for each $1\le\ell\le k_j$, and send
$M(A,S)-k_j$ units directly from $s_j$ to~$s_{j+1}$.  Denote the resulting
flow by~$z^{A,S}$.  Every outside vertex in $A\setminus S$ has one unit of
incoming and outgoing flow.  Moreover, every backbone vertex has
$M(A,S)$ units entering from the preceding segment and the same amount
leaving along the next segment, so $z^{A,S}$ is a circulation.

By the directed triangle inequality, the cost of each two-arc detour
$s_j\to u_{j,\ell}\to s_{j+1}$, as well as the direct arc
$s_j\to s_{j+1}$, is at most the length of~$P_j$.  Consequently,
\begin{equation}\label{eq:realization-circulation-cost}
    \cost(z^{A,S})
    \le M(A,S)\cost(T_C)
    =M(A,S)\opt(C).
\end{equation}

For a fixed backbone~$S$, average this construction over~$A$:
\[
    z^S:=\E_A[z^{A,S}].
\]
If $i\in U$, then $i$ receives one unit of flow exactly when $i\in A$.
Hence its expected incoming and outgoing flow is~$p_i$.  All other
constraints in~\cref{eq:backbone-circulation-lp} are preserved under
averaging, so $z^S$ is feasible for $\mathrm{CIRC}(S)$.
The fixed-pair construction also gives the same degree at every backbone
vertex:
\begin{equation}\label{eq:ordinary-backbone-degree}
    z^S(\delta^+(s))=z^S(\delta^-(s))
    =Z(S)
    \quad\forall s\in R,
    \qquad
    Z(S):=\E_A[M(A,S)\mid S].
\end{equation}

\begin{lemma}[Longest outside run]\label{lem:circulation-longest-run}
For every integer $t\ge1$ and every fixed set~$C$ with
$\P(A\cup S=C)>0$,
\[
    \P\bigl(M(A,S)\ge t\mid A\cup S=C\bigr)
    \le n2^{-t}.
\]
In particular,
\[
    \E\bigl[M(A,S)\mid A\cup S=C\bigr]
    =O(\log n).
\]
\end{lemma}

\begin{proof}
Condition on $A\cup S=C$.  The tour~$T_C$ is then fixed.  For each
$v\in C$, the conditional probability that it is active but outside the
backbone is
\[
    \P(v\in A\setminus S\mid v\in A\cup S)
    =
    \frac{p_v(1-p_v)}{1-(1-p_v)^2}
    =
    \frac{1-p_v}{2-p_v}
    \le\frac12.
\]
The conditioning factors over vertices, so these classifications remain
independent.  A run of length~$t$ can start in at most~$n$ positions of the
fixed cyclic tour.  A union bound therefore gives the claimed tail estimate.

Let $L:=\lceil\log_2(\max\{n,2\})\rceil$.  The tail-sum formula yields
\[
    \E[M(A,S)\mid A\cup S=C]
    \le
    L+\sum_{t>L}n2^{-t}
    =O(\log n).
\]
\end{proof}

\begin{proposition}[Low-cost backbone and circulation]
\label{prop:low-cost-backbone-circulation}
If $S$ is sampled independently with $\P(v\in S)=p_v$, then
\[
    \E_S\bigl[\opt(S)+c(S)\bigr]
    \le O(\log n)\optpost.
\]
Consequently, there exists a fixed backbone~$S$ for which
\[
    \opt(S)+c(S)
    \le O(\log n)\optpost.
\]
\end{proposition}

\begin{proof}
The random set~$S$ has the same distribution as an activation set, and hence
\[
    \E_S[\opt(S)]=\optpost.
\]
Using the feasible solution~$z^S$, linearity of cost, and
\cref{eq:realization-circulation-cost,lem:circulation-longest-run}, we obtain
\begin{align*}
    \E_S[c(S)]
    &\le \E_S[\cost(z^S)]\\
    &=\E_{A,S}[\cost(z^{A,S})]\\
    &\le\E_{A,S}[M(A,S)\opt(A\cup S)]\\
    &\le O(\log n)\,\E_{A,S}[\opt(A\cup S)].
\end{align*}
For every pair~$A,S$, concatenate optimal tours on
$A\cup\{r\}$ and $S\cup\{r\}$ and shortcut repeated vertices.  The directed
triangle inequality gives
\[
    \opt(A\cup S)\le\opt(A)+\opt(S).
\]
Both terms on the right have expectation~$\optpost$, so
\[
    \E_S[c(S)]\le O(\log n)\optpost.
\]
Adding the expected backbone-tour cost proves the first statement.  The
existence statement follows by averaging.
\end{proof}

The final adaptive construction needs a slightly more careful consequence of
the same averaging argument.  In particular, the backbone tour itself must
cost $O(\optpost)$ rather than merely $O(\log n)\optpost$.

\begin{lemma}[Ordinary-backbone certificate]
\label{lem:ordinary-backbone-certificate}
Fix an induced subinstance on a client set $W\subseteq V^\circ$, and write
\[
    B:=\optpost(W):=\E[\opt(A\cap W)].
\]
If $B>0$, there are a set $S\subseteq W$, a tour~$T_R$ on
$R:=S\cup\{r\}$, and a feasible circulation~$z$ for the analogue of
$\mathrm{CIRC}(S)$ on $W\cup\{r\}$ such that
\[
    \cost(T_R)\le C_1B,
    \qquad
    \cost(z)\le C_1\max\{1,\log(n+1)\}B
\]
for a universal constant~$C_1$.
\end{lemma}

\begin{proof}
Apply the preceding random-backbone construction to the induced subinstance
on $W\cup\{r\}$ and retain its averaged witness~$z^S$.  Because~$S$ has the
posterior distribution on~$W$,
\[
    \E_S[\opt(S)]=B.
\]
The proof of \cref{prop:low-cost-backbone-circulation}, with the ambient
value~$n$ kept in the longest-run estimate, gives
\[
    \E_S[\cost(z^S)]\le C_2\max\{1,\log(n+1)\}B
\]
for a universal constant~$C_2$.  Hence the nonnegative random variable
\[
    \frac{\opt(S)}{B}
    +\frac{\cost(z^S)}{C_2\max\{1,\log(n+1)\}B}
\]
has expectation at most~$2$.  Some outcome makes both summands at most a
universal constant.  Fix that outcome, take~$T_R$ to be an optimal tour on
$R$, and enlarge~$C_1$ if necessary.
\end{proof}

\subsection{A boosted backbone}
\label{sec:boosted-backbone}

For the sharper rounding and the offline construction below, we need two
additional properties: bounded circulation degree at backbone vertices and
small activation probabilities outside the backbone.  Both can be obtained
by boosting the sampling probabilities.

\begin{proposition}[Boosted backbone]\label{prop:boosted-backbone}
There are absolute constants $c,D_0,K_0>0$ such that the following holds.
Put
\[
    \lambda:=\left\lceil c\log n\right\rceil
    \qquad\text{and}\qquad
    \widehat p_v:=\min\{1,\lambda p_v\}.
\]
There exist a set~$S\subseteq V^\circ$, a tour~$T_R$ on
$R:=S\cup\{r\}$, and a feasible solution~$z$ to
$\mathrm{CIRC}(S)$ such that
\begin{align}
    \cost(T_R)+\cost(z)
    &\le K_0\log n\optpost,
    \label{eq:boosted-cost}\\
    z(\delta^+(s))=z(\delta^-(s))
    &\le D_0\log n
    \qquad\forall s\in R,
    \label{eq:boosted-degree}\\
    p_i&<\frac1\lambda
    \qquad\forall i\in V^\circ\setminus S.
    \label{eq:boosted-outside-probability}
\end{align}
The set~$S$ may be chosen from the product distribution with inclusion
probabilities~$(\widehat p_v)_{v\in V^\circ}$.
\end{proposition}

\begin{proof}
First consider one ordinary sample~$B$, in which each non-root vertex~$v$ is included
with probability~$p_v$, and the averaged witness~$z^B$ constructed above.
\Cref{eq:ordinary-backbone-degree,lem:circulation-longest-run} give
\[
    \E_B[Z(B)]=O(\log n).
\]
The calculation in the proof of
\cref{prop:low-cost-backbone-circulation} also gives
\[
    \E_B[\cost(z^B)]
    =O(\log n)\optpost.
\]
By Markov's inequality, after choosing absolute constants $D_0$ and~$K_0$
large enough, both
\[
    Z(B)\le D_0\log n
    \quad\text{and}\quad
    \cost(z^B)\le K_0\log n\optpost
\]
hold with probability at least~$3/4$.  Call such a sample \emph{good}.

Take $\lambda$ independent ordinary samples
$B^1,\ldots,B^\lambda$ and let $S_0:=\bigcup_jB^j$.  The probability that
none is good is at most $4^{-\lambda}$.  Moreover,
\[
    \P(v\in S_0)=1-(1-p_v)^\lambda
    \le\min\{1,\lambda p_v\}=\widehat p_v.
\]
Thus a product sample~$S$ with marginals~$\widehat p_v$ can be coupled with~$S_0$ so
that $S_0\subseteq S$.

Suppose $B^j$ is good.  The same vector~$z^{B^j}$ remains feasible when the
backbone is enlarged from~$B^j$ to~$S$.  Indeed, every vertex still outside
$S$ was outside $B^j$ and retains incoming and outgoing degree~$p_v$.
Every new backbone vertex $s\in S\setminus B^j$ previously had degree
$p_s\le1$, while the old backbone vertices have degree at most
$D_0\log n$.  The support restriction can only become weaker when the
outside set shrinks.  Hence~\cref{eq:boosted-degree} holds and the circulation
cost is unchanged.

It remains to control the backbone tour.  After sampling~$S$, independently
give each $v\in S$ a uniform label in~$[\lambda]$, and let $S_j$ contain the
vertices carrying label~$j$.  The sets~$S_j$ partition~$S$, and, for each
fixed~$j$, vertex memberships are independent with
\[
    \P(v\in S_j)=\frac{\widehat p_v}{\lambda}
    =\min\left\{\frac1\lambda,p_v\right\}
    \le p_v.
\]
By a monotone coupling, $\E[\opt(S_j)]\le\optpost$.  Concatenating tours for
the sets~$S_j$ at the root and shortcutting gives
\[
    \opt(S)\le\sum_{j=1}^{\lambda}\opt(S_j),
    \qquad
    \E[\opt(S)]\le\lambda\optpost.
\]
Markov's inequality therefore gives
$\opt(S)\le4\lambda\optpost$ with probability at least~$3/4$.
For a sufficiently large choice of~$c$, this event and the existence of a
good~$B^j$ have positive joint probability.  Fix an outcome in their
intersection and let~$T_R$ be an optimal tour on~$R$.  Enlarging~$K_0$ absorbs
the term $4\lambda\optpost$ and proves~\cref{eq:boosted-cost}.

Finally, if $i\notin S$, then $\widehat p_i<1$.  By the definition of~$\widehat p_i$, this
implies $\lambda p_i<1$, which is
\cref{eq:boosted-outside-probability}.
\end{proof}

\section{Endpoint rounding with interval discrepancy}
\label{sec:circulation-rounding}

Fix a backbone~$R=S\cup\{r\}$ and a feasible circulation~$z$.  Since there
are no arcs between outside vertices, the flow through each $i\in U$ can be
coupled into paths with both endpoints in~$R$.  More precisely, there are
nonnegative variables
\[
    x_{i,ab}
    \qquad(i\in U,\ a,b\in R)
\]
such that
\begin{equation}\label{eq:coupled-marginals}
    \sum_{b\in R}x_{i,ab}=z_{ai},
    \qquad
    \sum_{a\in R}x_{i,ab}=z_{ib}.
\end{equation}
Such a coupling exists because both marginal sums equal~$p_i$; it can be
obtained, for example, by greedily matching the incoming and outgoing flow
at~$i$.  Put $y_{ab}:=z_{ab}$ for $a,b\in R$.  Then
\begin{equation}\label{eq:coupled-outside-constraint}
    \sum_{a,b\in R}x_{i,ab}=p_i
    \qquad\forall i\in U,
\end{equation}
and the coupling preserves cost:
\begin{align}
    &\sum_{i\in U}\sum_{a,b\in R}
       \bigl(d(a,i)+d(i,b)\bigr)x_{i,ab}
      +\sum_{a,b\in R}d(a,b)y_{ab}
      \notag\\
    &\hspace{35mm}=
      \sum_{u,v\in V}d(u,v)z_{uv}.
    \label{eq:coupled-cost}
\end{align}

Let
\[
    T_R=(s_0,s_1,\ldots,s_{m-1},s_0)
\]
be a tour on the backbone.  A \emph{cyclic interval} is a consecutive set of
vertices in this cyclic order.  For $I\subseteq R$ and $a,b\in R$, define
\[
    \chi_I(a,b):=\1_{\{a\in I\}}-\1_{\{b\in I\}},
\]
and define the weighted outgoing-minus-incoming imbalance of~$I$ by
\begin{equation}\label{eq:interval-imbalance}
    \Delta_{x,y}(I)
    :=
    \sum_{i\in U}\sum_{a,b\in R}\chi_I(a,b)x_{i,ab}
    +\sum_{a,b\in R}\chi_I(a,b)y_{ab}.
\end{equation}
Summing the original flow-conservation equations over the backbone vertices
in~$I$ shows that
\[
    \Delta_{x,y}(I)=0
    \qquad\forall I\subseteq R.
\]

The goal is to replace every $x_{i,ab}$ by either~$0$ or~$p_i$.  By
\cref{eq:coupled-outside-constraint}, this assigns every outside vertex of
positive probability to exactly one ordered pair~$(a,b)$ of backbone
vertices.  Preserving all subset-balance equations while doing this need not
leave a direction in which to move.  We instead preserve only a laminar
family of interval equations and discard an equation once few fractional
variables still touch it.

\begin{lemma}[Endpoint rounding]\label{lem:circulation-endpoint-rounding}
Let $p_{\max}:=\max_{i\in U}p_i$, with $p_{\max}:=0$ when $U=\emptyset$.
There is a vector~$\widehat x$ satisfying
\[
    \widehat x_{i,ab}\in\{0,p_i\}
    \qquad\text{and}\qquad
    \sum_{a,b\in R}\widehat x_{i,ab}=p_i
    \quad\forall i\in U
\]
such that
\begin{align}
    \sum_{i\in U}\sum_{a,b\in R}
       \bigl(d(a,i)+d(i,b)\bigr)\widehat x_{i,ab}
    &\le
    \sum_{i\in U}\sum_{a,b\in R}
       \bigl(d(a,i)+d(i,b)\bigr)x_{i,ab},
    \label{eq:rounding-cost}\\
    |\Delta_{\widehat x,y}(I)|
    &\le O\bigl(p_{\max}\log^2(n+1)\bigr).
    \label{eq:rounding-interval-balance}
\end{align}
The second inequality holds for every cyclic interval~$I$ of~$T_R$.
\end{lemma}

\begin{proof}
The claim is immediate when $|R|=1$, so assume $m:=|R|\ge2$.  Pad the linear
order $s_0,\ldots,s_{m-1}$ with empty positions to the next power of two and
let $\mathfrak D$ be the standard dyadic interval family represented by the
nodes of the resulting binary interval tree.  Put
\[
    h:=\lceil\log_2 m\rceil,
    \qquad
    L:=8(h+1).
\]
Each backbone position belongs to at most $h+1$ members of~$\mathfrak D$.
Consequently, a variable~$x_{i,ab}$ has nonzero coefficient in at most
$2(h+1)$ of the equations
\[
    \Delta_{x,y}(D)=0,
    \qquad D\in\mathfrak D.
\]

Call a variable \emph{live} when $0<x_{i,ab}<p_i$, and let~$F$ be the set of
live variables.  Throughout the procedure we preserve
\cref{eq:coupled-outside-constraint}.  We also preserve a dyadic interval
equation as long as it touches at least~$L$ live variables; once fewer than
$L$ live variables touch it, we discard it permanently.

Let $c_1$ be the number of outside equations that contain a live variable.
No such equation can contain exactly one live variable: after all other
variables in the equation have values in~$\{0,p_i\}$, the equality with
right-hand side~$p_i$ would force the remaining variable to an endpoint as
well.  Hence
\[
    c_1\le\frac{|F|}{2}.
\]
Let $c_2$ be the number of currently preserved dyadic equations.  Counting
incidences between these equations and the live variables gives
\[
    c_2L
    \le 2(h+1)|F|,
\]
and therefore $c_2\le |F|/4$.  Thus
\[
    c_1+c_2<|F|.
\]
The homogeneous system associated with the preserved equalities has a
nonzero direction supported on~$F$.

Move in both signs of this direction until a live variable first reaches
either~$0$ or its upper bound~$p_i$.  The current point is a convex
combination of the two resulting endpoints.  Since the objective in
\cref{eq:rounding-cost} is linear, at least one endpoint has cost no larger
than the current point; choose such an endpoint.  This preserves all current
equalities, never increases cost, and makes at least one additional variable
nonlive.  Repeating the step terminates with every variable in
$\{0,p_i\}$, while all outside equations still hold.  This proves
\cref{eq:rounding-cost} and the first displayed properties of~$\widehat x$.

Consider a dyadic equation at the moment it is discarded.  Its imbalance is
still zero, and fewer than~$L$ of its variables are live.  Only those
variables can subsequently change.  Each changes in absolute value by at
most $p_i\le p_{\max}$ and has coefficient in~$\{-1,0,1\}$.  It follows that
at the end of the procedure,
\[
    |\Delta_{\widehat x,y}(D)|\le Lp_{\max}
    \qquad\forall D\in\mathfrak D.
\]
Every linear interval has a canonical decomposition into $O(h+1)$ disjoint
dyadic intervals.  A cyclic interval is the union of at most two linear
intervals, so it has a decomposition into $O(h+1)$ members of~$\mathfrak D$.
The imbalance is additive over disjoint unions.  Therefore
\[
    |\Delta_{\widehat x,y}(I)|
    \le O((h+1)Lp_{\max})
    =O\bigl(p_{\max}\log^2(n+1)\bigr)
\]
for every cyclic interval~$I$, proving
\cref{eq:rounding-interval-balance}.
\end{proof}

Combining the low-cost construction and endpoint rounding gives the first
structural result of the chapter.

\begin{theorem}[Backbone circulation and rounding]
\label{thm:backbone-circulation-rounding}
For every \StochTSP{} instance on an asymmetric metric, there exist a backbone
$R=S\cup\{r\}$, a tour
$T_R=(s_0,\ldots,s_{m-1},s_0)$, nonnegative backbone variables
$(y_{ab})_{a,b\in R}$, and, for every
$i\in V^\circ\setminus S$ with $p_i>0$, an
ordered pair $(a_i,b_i)\in R\times R$ such that
\begin{align}
    \cost(T_R)
    +\sum_{a,b\in R}d(a,b)y_{ab}
    +\sum_{\substack{i\in V^\circ\setminus S\\p_i>0}}
       p_i\bigl(d(a_i,i)+d(i,b_i)\bigr)
    \le O(\log n)\optpost.
    \label{eq:partial-circulation-total-cost}
\end{align}
Moreover,
\begin{equation}\label{eq:rounded-backbone-degree}
    \sum_{b\in R}y_{ab}\le O(\log n),
    \qquad
    \sum_{b\in R}y_{ba}\le O(\log n)
    \qquad\forall a\in R,
\end{equation}
and every outside vertex satisfies
\begin{equation}\label{eq:rounded-small-probability}
    p_i\le O\left(\frac1{\log n}\right).
\end{equation}
If $\widehat x_{i,a_i b_i}:=p_i$ and all other $\widehat x$-variables are
zero, then for every cyclic interval~$I$ of~$T_R$,
\[
    |\Delta_{\widehat x,y}(I)|
    \le O(\log n).
\]
\end{theorem}

\begin{proof}
Choose $S$ and a feasible circulation~$z$ as in
\cref{prop:boosted-backbone}.  Couple the flow through each outside
vertex according to~\cref{eq:coupled-marginals}, and apply
\cref{lem:circulation-endpoint-rounding}.  The coupling preserves the
circulation cost by~\cref{eq:coupled-cost}, and the rounding does not
increase it by~\cref{eq:rounding-cost}.  Thus
\cref{eq:partial-circulation-total-cost} follows from
\cref{eq:boosted-cost}.  The backbone variables~$y$ form a subvector of~$z$,
so~\cref{eq:rounded-backbone-degree} follows from
\cref{eq:boosted-degree}; \cref{eq:rounded-small-probability} is
\cref{eq:boosted-outside-probability}.  Finally,
$p_{\max}=O(1/\log n)$, and hence
\cref{eq:rounding-interval-balance} improves to
$O(\log n)$.
\end{proof}

\section{An offline bounded-multiplicity closed walk}
\label{sec:offline-backbone-walk}

We now use the rounded structure after the active set has been revealed.
The construction in this section is therefore offline.  Its main tool is the
following classical form of Hoffman's circulation theorem~\cite{Hoffman1960}.

\begin{theorem}[Hoffman's circulation theorem]\label{thm:hoffman-circulation}
Let $G=(W,E)$ be a directed multigraph and let
$\ell,u:E\to\mathbb R_{\ge0}$ satisfy $\ell_e\le u_e$ for every arc~$e$.
There is a circulation~$g$ with
\[
    \ell_e\le g_e\le u_e
    \qquad\forall e\in E
\]
if and only if, for every $X\subseteq W$,
\begin{equation}\label{eq:hoffman-condition}
    \ell(\delta^-(X))\le u(\delta^+(X)).
\end{equation}
If $\ell$ and~$u$ are integral, then~$g$ may be chosen integral.
\end{theorem}

Fix the objects supplied by
\cref{thm:backbone-circulation-rounding}.  For an activation
set~$A$, every active outside vertex $i\in A\cap U$ supplies the directed
chord
\[
    e_i=(a_i,b_i)
\]
on~$R$; after constructing a circulation on these collapsed chords, we will
expand $e_i$ back to the path $a_i\to i\to b_i$.  Parallel chords are kept as
distinct arcs.

For $X\subseteq R$, write
\begin{align*}
    F_A^-(X)
    &:=
    \bigl|\{i\in A\cap U:a_i\notin X,\ b_i\in X\}\bigr|,\\
    F_A^+(X)
    &:=
    \bigl|\{i\in A\cap U:a_i\in X,\ b_i\notin X\}\bigr|,
\end{align*}
and let their expectations be
\begin{align*}
    \mu^-(X)
    &:=
    \sum_{\substack{i\in U\\a_i\notin X,\ b_i\in X}}p_i,\\
    \mu^+(X)
    &:=
    \sum_{\substack{i\in U\\a_i\in X,\ b_i\notin X}}p_i.
\end{align*}
Also put
\[
    y^+(X):=\sum_{\substack{a\in X\\b\notin X}}y_{ab},
    \qquad
    y^-(X):=\sum_{\substack{a\notin X\\b\in X}}y_{ab}.
\]

Let $C_R$ denote the directed backbone cycle given by~$T_R$, and define
\[
    \iota(X):=|\delta_{C_R}^+(X)|.
\]
For a nonempty proper subset~$X$, this is the number of cyclic intervals in
the decomposition of~$X$; it also equals
$|\delta_{C_R}^-(X)|$.  Set $\iota(\emptyset)=\iota(R)=0$.

\begin{lemma}[Uniform cut concentration]\label{lem:offline-cut-concentration}
There is an absolute constant $C_0$ such that, with probability at least
$1-(n+1)^{-5}$ over the activation set~$A$, every $X\subseteq R$ satisfies
\begin{align}
    F_A^-(X)
    &\le2\mu^-(X)+C_0\iota(X)\log n,
    \label{eq:cut-concentration-in}\\
    \mu^+(X)
    &\le2F_A^+(X)+C_0\iota(X)\log n.
    \label{eq:cut-concentration-out}
\end{align}
\end{lemma}

\begin{proof}
Fix $X$ with $k:=\iota(X)\ge1$.  The variable $F_A^-(X)$ is a sum of independent
Bernoulli variables with mean~$\mu^-(X)$.  A standard Chernoff bound gives,
for every $t\ge0$,
\[
    \P\bigl(F_A^-(X)>2\mu^-(X)+t\bigr)\le\e^{-t/3}.
\]
For the other direction, if $\mu^+(X)\le t$, then
\cref{eq:cut-concentration-out} holds automatically with~$t$ in place of
its additive term.  If $\mu^+(X)>t$, a violation implies
$F_A^+(X)<\mu^+(X)/2$, whose probability is at most
\[
    \e^{-\mu^+(X)/8}\le\e^{-t/8}.
\]

There are at most $|R|^{2k}\le(n+1)^{2k}$ subsets of a cyclic order having
$k$ interval components: it is enough to choose their $2k$ boundary
positions.  Set $t=C_0k\log n$, choose $C_0$ sufficiently large, and take
a union bound over all $k\ge1$ and both inequalities.  The resulting failure
probability is at most $(n+1)^{-5}$.  The two sets with $k=0$ satisfy both
inequalities trivially.
\end{proof}

\begin{theorem}[Offline bounded-multiplicity walk]
\label{thm:offline-backbone-walk}
Let $A$ be an independent activation set.  With probability at least
$1-(n+1)^{-5}$, there is an offline closed walk~$W_A$ with the following
properties:
\begin{enumerate}[label=(\roman*)]
    \item for every active outside vertex $i\in A\cap U$, the walk traverses
    both arcs of $a_i\to i\to b_i$ between one and four times;
    \item every backbone-cycle edge is traversed at least once and at most
    $O(\log n)$ times; and
    \item every auxiliary backbone arc $a\to b$ with $y_{ab}>0$ is traversed
    at most $O(\log n)$ times.
\end{enumerate}
Thus every active outside chord is used at least once and every edge used by
the offline walk has multiplicity $O(\log n)$.
\end{theorem}

\begin{proof}
Fix a realization satisfying
\cref{lem:offline-cut-concentration}.  By
\cref{thm:backbone-circulation-rounding}, there is an absolute constant
$C_\Delta$ such that every cyclic interval~$I$ satisfies
\[
    |\Delta_{\widehat x,y}(I)|
    \le C_\Delta\log n.
\]
If $X$ is the disjoint union of $\iota(X)$ cyclic intervals, additivity of the
imbalance gives
\[
    \bigl|
       \mu^+(X)-\mu^-(X)+y^+(X)-y^-(X)
    \bigr|
    \le C_\Delta\iota(X)\log n.
\]
In particular,
\begin{equation}\label{eq:mean-cut-comparison}
    \mu^-(X)
    \le\mu^+(X)+y^+(X)+C_\Delta\iota(X)\log n.
\end{equation}
Combining
\cref{eq:cut-concentration-in,eq:mean-cut-comparison,eq:cut-concentration-out}
yields
\begin{equation}\label{eq:realized-cut-comparison}
    F_A^-(X)
    \le
    4F_A^+(X)+2y^+(X)+
    (2C_\Delta+3C_0)\iota(X)\log n
    \qquad\forall X\subseteq R.
\end{equation}

We apply \cref{thm:hoffman-circulation} to a directed multigraph on
vertex set~$R$.  It has three types of arcs, with the following lower and
upper bounds:
\begin{enumerate}
    \item each active outside chord~$e_i=(a_i,b_i)$ has
    $\ell_{e_i}=1$ and $u_{e_i}=4$;
    \item each auxiliary backbone arc $a\to b$ with $y_{ab}>0$ has
    $\ell_{ab}=0$ and $u_{ab}=\lceil2y_{ab}\rceil$; and
    \item each edge of the directed cycle~$C_R$ has lower bound~$1$ and
    upper bound
    \[
        K:=
        \left\lceil
        (2C_\Delta+3C_0)\log n
        \right\rceil+1.
    \]
\end{enumerate}
Coincident arcs of different types are treated as parallel arcs.

For every $X\subseteq R$, the total lower bound on arcs entering~$X$ is
\[
    \ell(\delta^-(X))=F_A^-(X)+\iota(X).
\]
The total upper bound on arcs leaving~$X$ satisfies
\[
    u(\delta^+(X))
    \ge4F_A^+(X)+2y^+(X)+K\iota(X).
\]
\Cref{eq:realized-cut-comparison} and the definition of~$K$
therefore imply Hoffman's condition~\cref{eq:hoffman-condition}.  All bounds
are integral, so there is an integral circulation~$g$ within them.

Every auxiliary backbone arc has
\[
    g_{ab}\le\lceil2y_{ab}\rceil
    \le 2D_0\log n+1
\]
by~\cref{eq:rounded-backbone-degree}.  The cycle-edge multiplicity is at
most~$K=O(\log n)$, while every active outside chord has multiplicity
between one and four.

Finally, expand every chord~$e_i$ carrying $g_{e_i}$ units into
$g_{e_i}$ copies of the path $a_i\to i\to b_i$.  The resulting directed
multigraph is balanced.  Its support is strongly connected: it contains the
whole directed backbone cycle, and every active outside vertex has a path
from and back to that cycle.  It therefore has an Euler tour, which is the
claimed closed walk~$W_A$.
\end{proof}

\section{From the offline circulation to an adaptive policy}
\label{sec:adaptive-circulation}

For $W\subseteq V^\circ$, write
\[
    \optpost(W):=\E[\opt(A\cap W)]
\]
and let $\optadapt(W)$ denote the adaptive optimum of the induced subinstance
on $W\cup\{r\}$.

The preceding results identify the obstruction but still choose the final
circulation after seeing~$A$.  We now close this gap.  The key change is to
return to the ordinary-backbone certificate of
\cref{lem:ordinary-backbone-certificate}, whose backbone tour costs only
$O(\optpost)$, and to make all endpoint and auxiliary-flow choices before any
activation is probed.  The integer auxiliary flow will cancel the
\emph{expected} divergence of the client chords on every cyclic interval up
to one unit.  The realization then contributes only a centered sum of
independent Bernoulli variables.

\subsection{Endpoint sampling and maximal inequalities}
\label{sec:adaptive-endpoint-sampling}

Fix an induced subinstance on $W\subseteq V^\circ$ and write
\[
    \mu:=\sum_{i\in W}p_i,
    \qquad
    B:=\optpost(W),
    \qquad
    L:=\max\{1,\log(n+1)\}.
\]
Assume for now that $B>0$.  Fix the set~$S$, backbone
$R=S\cup\{r\}$, tour~$T_R$, and circulation~$z$ supplied by
\cref{lem:ordinary-backbone-certificate}, and put $U:=W\setminus S$.

Because $z$ has no arcs between two outside clients, its incoming and
outgoing flow at each $i\in U$ can be coupled into variables
\[
    x_{i,ab}\ge0
    \qquad(a,b\in R)
\]
such that
\begin{equation}\label{eq:adaptive-coupling}
    \sum_{b\in R}x_{i,ab}=z_{ai},
    \qquad
    \sum_{a\in R}x_{i,ab}=z_{ib},
    \qquad
    \sum_{a,b\in R}x_{i,ab}=p_i.
\end{equation}
Let $y_{ab}:=z_{ab}$ for $a,b\in R$.  The coupling preserves cost:
\begin{align}
    &\sum_{i\in U}\sum_{a,b\in R}
       \bigl(d(a,i)+d(i,b)\bigr)x_{i,ab}
      +\sum_{a,b\in R}d(a,b)y_{ab}
      =\cost(z).
    \label{eq:adaptive-coupling-cost}
\end{align}

Write $T_R=(s_0,s_1,\ldots,s_{m-1},s_0)$.  For a cyclic interval~$I$ in
this order and a pair $(a,b)\in R\times R$, define
\[
    \chi_I(a,b):=\1_{\{a\in I\}}-\1_{\{b\in I\}}.
\]
We use the following one-dimensional maximal estimates.

\begin{lemma}[Interval maximal inequalities]
\label{lem:adaptive-interval-maximal}
Let all ordered pairs have endpoints in a fixed cyclically ordered set.
\begin{enumerate}[label=\textup{(\roman*)}]
    \item If $(a_i,b_i)$ are fixed and $X_i$ are independent centered random
    variables, then
    \[
        \E\max_I\left|\sum_iX_i\chi_I(a_i,b_i)\right|
        \le8\left(\sum_i\E[X_i^2]\right)^{1/2}.
    \]
    \item If $P_i=(a_i,b_i)$ are independent random pairs and $w_i\ge0$ are
    fixed, then
    \[
        \E\max_I\left|\sum_iw_i
        \bigl(\chi_I(P_i)-\E[\chi_I(P_i)]\bigr)\right|
        \le16\left(\sum_iw_i^2\right)^{1/2}.
    \]
\end{enumerate}
The maxima range over all cyclic intervals.
\end{lemma}

\begin{proof}
Let $J_k:=\{s_0,\ldots,s_k\}$ be a prefix.  A linear interval is the
difference of two prefixes, and a wrapped cyclic interval is the complement
of a linear interval.  Since $\chi_R(a,b)=0$, the maximum over cyclic
intervals is at most twice the maximum over prefixes.

For part~\textup{(i)}, split the prefix process as
\[
    \sum_iX_i\chi_{J_k}(a_i,b_i)
    =\sum_iX_i\1_{\{a_i\in J_k\}}
     -\sum_iX_i\1_{\{b_i\in J_k\}}.
\]
Order the indices by the positions of their tails.  The tail sums at prefix
boundaries form a subset of the partial sums in this order; ties may be
ordered arbitrarily.  The $L^2$ maximal inequality gives
\[
    \E\max_j\left|\sum_{\ell\le j}X_\ell\right|
    \le2\left(\sum_i\E[X_i^2]\right)^{1/2}.
\]
The same argument applies after ordering the heads.  The triangle inequality
and the factor two for cyclic intervals prove~\textup{(i)}.

For part~\textup{(ii)}, first consider the random tails.  Symmetrization with
independent Rademacher signs~$\varepsilon_i$ gives
\begin{align*}
    &\E\max_k\left|\sum_iw_i
      \bigl(\1_{\{a_i\in J_k\}}-\P(a_i\in J_k)\bigr)\right|\\
    &\qquad\le2\E\max_k\left|\sum_i\varepsilon_iw_i
      \1_{\{a_i\in J_k\}}\right|
    \le4\left(\sum_iw_i^2\right)^{1/2}.
\end{align*}
For the last inequality, condition on all tails, order them, and apply the
same $L^2$ maximal inequality to the partial sums of
$\varepsilon_iw_i$.  Repeat for the heads, then use the triangle inequality
and the factor two for cyclic intervals.
\end{proof}

For every $i\in U$ with $p_i>0$, independently sample one public endpoint
pair $P_i=(a_i,b_i)$ according to
\begin{equation}\label{eq:adaptive-endpoint-sampling}
    \P(P_i=(a,b))=\frac{x_{i,ab}}{p_i}.
\end{equation}
If $p_i=0$, put $P_i=(r,r)$.  Define zero-sum divergence vectors on~$R$ by
\begin{align}
    \bar b_s
    &:=\sum_{i\in U}\sum_{a,b\in R}x_{i,ab}
       \bigl(\1_{\{a=s\}}-\1_{\{b=s\}}\bigr),
    \label{eq:adaptive-mean-divergence}\\
    b_s(P)
    &:=\sum_{i\in U}p_i
       \bigl(\1_{\{a_i=s\}}-\1_{\{b_i=s\}}\bigr).
    \label{eq:adaptive-sampled-divergence}
\end{align}
For a flow~$f$ on~$R$, let
\[
    \operatorname{div}(f)_s:=\sum_{t\in R}f_{st}-\sum_{t\in R}f_{ts}.
\]
Since~$z$ is a circulation,
\begin{equation}\label{eq:adaptive-y-divergence}
    \operatorname{div}(y)=-\bar b.
\end{equation}
By \cref{lem:adaptive-interval-maximal}\textup{(ii)},
\begin{equation}\label{eq:adaptive-gamma-bound}
    \E_P[\Gamma(P)]
    \le16\left(\sum_{i\in U}p_i^2\right)^{1/2}
    \le16\sqrt\mu,
\end{equation}
where
\[
    \Gamma(P):=
    \max_I\left|\sum_{s\in I}\bigl(b_s(P)-\bar b_s\bigr)\right|.
\]
Moreover, \cref{eq:adaptive-endpoint-sampling} gives
\begin{align}
    &\E_P\left[\sum_{i\in U}p_i
      \bigl(d(a_i,i)+d(i,b_i)\bigr)\right] \notag\\
    &\qquad=
      \sum_{i\in U}\sum_{a,b\in R}x_{i,ab}
      \bigl(d(a,i)+d(i,b)\bigr).
    \label{eq:adaptive-sampled-detour-cost}
\end{align}

\subsection{Deterministic auxiliary-flow rounding}
\label{sec:integer-auxiliary-rounding}

We next round the auxiliary backbone flow.  Unlike the Hoffman circulation in
\cref{sec:offline-backbone-walk}, this integer flow is fixed from the public
endpoint outcome~$P$ and does not depend on the activation set.

\begin{lemma}[Integer interval balancing]
\label{lem:integer-interval-balancing}
Let $C=(s_0,\ldots,s_{m-1},s_0)$ be a directed cycle, and let~$y$ be a
nonnegative fractional flow on the complete directed graph on its vertices
with $\operatorname{div}(y)=-\bar b$ for a real zero-sum vector~$\bar b$.
For another real zero-sum vector~$b$, put
\[
    \Gamma:=\max_I\left|\sum_{s\in I}(b_s-\bar b_s)\right|.
\]
There is a nonnegative integer flow~$N$ such that
\begin{align}
    \cost(N)&\le\cost(y)+(\Gamma+1)\cost(C),
    \label{eq:integer-balancing-cost}\\
    \left|\sum_{s\in I}
      \bigl(b_s+\operatorname{div}(N)_s\bigr)\right|&\le1
    \qquad\text{for every cyclic interval }I.
    \label{eq:integer-balancing-interval}
\end{align}
\end{lemma}

\begin{proof}
If $m=1$, both divergence vectors vanish and $N=0$ works.  Assume $m\ge2$.
Put $\delta:=\bar b-b$ and
\[
    E_j:=\sum_{h=0}^j\delta_{s_h}.
\]
Thus $E_{m-1}=0$.  The difference of two $E$-values is the sum of~$\delta$
over a cyclic interval, and hence
$\max_jE_j-\min_jE_j\le\Gamma$.  Put
\[
    h_j:=E_j-\min_\ell E_\ell
\]
units of flow on the cycle edge $s_j\to s_{j+1}$.  Then
$0\le h_j\le\Gamma$ and, with cyclic indices,
\[
    \operatorname{div}(h)_{s_j}
    =h_j-h_{j-1}=\delta_{s_j}.
\]
Consequently, $w:=y+h$ is nonnegative,
$\operatorname{div}(w)=-b$, and
\begin{equation}\label{eq:first-divergence-correction}
    \cost(w)\le\cost(y)+\Gamma\cost(C).
\end{equation}

Write $t:=\operatorname{div}(w)=-b$ and
$T_j:=\sum_{h=0}^jt_{s_h}$.  For $j<m-1$, round~$T_j$ to a nearest
integer~$Q_j$; put $Q_{m-1}=0$ and $Q_{-1}:=Q_{m-1}$.  Define
\[
    q_{s_j}:=Q_j-Q_{j-1},
    \qquad
    e_j:=Q_j-T_j.
\]
Then $q$ is an integer zero-sum vector, $|e_j|\le1/2$, and
\begin{equation}\label{eq:rounded-divergence-error}
    q_{s_j}-t_{s_j}=e_j-e_{j-1}.
\end{equation}
Put $k_j:=e_j-\min_\ell e_\ell$ units on $s_j\to s_{j+1}$.  Then
$0\le k_j\le1$, \cref{eq:rounded-divergence-error} gives
$\operatorname{div}(k)=q-t$, and $w':=w+k$ is a nonnegative flow of integer
divergence~$q$ satisfying
\begin{equation}\label{eq:second-divergence-correction}
    \cost(w')\le\cost(w)+\cost(C).
\end{equation}

Consider the bounded flow polytope
\[
    \lfloor w'_e\rfloor\le f_e\le\lceil w'_e\rceil,
    \qquad
    \operatorname{div}(f)=q.
\]
It is nonempty because it contains~$w'$.  The node--arc incidence matrix is
totally unimodular, and all bounds and right-hand sides are integral.  Thus a
minimum-cost solution~$N$ is integral.  Since $w'$ itself is feasible,
$\cost(N)\le\cost(w')$, which proves
\cref{eq:integer-balancing-cost}.

Finally, $\operatorname{div}(N)=q$ and $t=-b$, so
\[
    \sum_{s\in I}\bigl(b_s+\operatorname{div}(N)_s\bigr)
    =\sum_{s\in I}(q_s-t_s).
\]
By \cref{eq:rounded-divergence-error}, the last sum is a difference of two
$e$-values for every cyclic interval, and its absolute value is at most~$1$.
This proves \cref{eq:integer-balancing-interval}.
\end{proof}

For each public endpoint outcome~$P$, apply
\cref{lem:integer-interval-balancing} with $C=T_R$, $b=b(P)$, $\bar b$ from
\cref{eq:adaptive-mean-divergence}, and $y$ from
\cref{eq:adaptive-y-divergence}.  Denote the resulting finite integer multiset
of auxiliary backbone arcs by~$N(P)$.  It is chosen before any activation is
probed, and every auxiliary copy will be used.

\subsection{Adaptive traversal}
\label{sec:adaptive-stack-traversal}

For a multiset~$F$ of directed chords between vertices of a directed cycle,
put
\[
    \Lambda(F):=\max_I\left|\sum_{e\in F}\chi_I(e)\right|.
\]

\begin{lemma}[Active-first stack walk]
\label{lem:active-first-stack-walk}
Starting at the root, repeatedly use an unused outgoing chord at the current
cycle vertex if one exists; otherwise use the next directed cycle edge.  Once
all chords have been used, continue along the cycle to the root.  The
resulting closed walk uses every chord exactly once and every cycle edge at
most $\Lambda(F)+1$ times.
\end{lemma}

\begin{proof}
Finiteness and the directed cycle ensure that every chord is eventually used.
Let $c_e$ be the number of traversals of cycle edge~$e$.  Flow balance of the
completed closed walk over any cyclic interval says that the difference of
its two boundary cycle-edge counts is the negative of the chord imbalance.
Hence $|c_e-c_{e'}|\le\Lambda(F)$ for any two cycle edges.

If some cycle edge is unused, the conclusion follows.  Otherwise consider the
cycle edge whose first use occurs last.  Before a cycle edge is first used,
all outgoing chords at its tail have been exhausted.  By the time of this
last first use, this has happened at every cycle vertex.  No chord can appear
later, so the walk subsequently follows the cycle to the root without using
this same edge again.  Its final multiplicity is one, and all other
multiplicities are at most $\Lambda(F)+1$.
\end{proof}

For fixed $P$, $N(P)$, and activation set~$A$, let
\[
    F(A,P):=N(P)\cup\{a_i\to b_i:i\in A\cap U\}.
\]

\begin{lemma}[Adaptive virtual execution]
\label{lem:adaptive-virtual-execution}
For fixed public $P$ and~$N(P)$, there is an adaptive policy whose
actual cost on every realization~$A$ is at most
\begin{align}
    &\sum_{a,b\in R}N_{ab}(P)d(a,b)
    +\sum_{i\in A\cap U}\bigl(d(a_i,i)+d(i,b_i)\bigr)\notag\\
    &\hspace{32mm}
    +(\Lambda(F(A,P))+2)\cost(T_R).
    \label{eq:adaptive-virtual-cost}
\end{align}
Every client is probed, and every active client is visited immediately after
it is probed.
\end{lemma}

\begin{proof}
Simulate the stack walk at a virtual backbone position.  On the first virtual
arrival at a backbone client~$s\in S$, probe~$s$.  At a virtual vertex~$a$,
probe its as-yet-unprobed outside clients with tail~$a_i=a$ in a fixed order
until an active client is found or the list is exhausted.  If active client
$i$ is found, serve it immediately and let the virtual walk take the chord
$a_i\to i\to b_i$.  Once the outside list at~$a$ is exhausted, take an unused
auxiliary chord leaving~$a$, if one exists; otherwise advance along the next
backbone-cycle edge.  Continue until all clients have been probed, all
auxiliary copies have been used, and the virtual walk has returned to~$r$.
Every choice uses only public data and outcomes already revealed.

Maintain a pending directed walk from the actual position to the virtual
backbone position.  A virtual cycle or auxiliary arc is appended without
actual movement.  If an active backbone client at the virtual vertex is
probed, the direct actual move to it costs at most the pending length, after
which the pending walk is empty.  If an active outside client~$i$ is found
while the virtual position is~$a_i$, the direct move to~$i$ costs at most the
old pending length plus $d(a_i,i)$; the new pending walk is $i\to b_i$.
Inactive probes cause no movement.  Inductively,
\[
    \text{actual cost paid}+\text{pending length}
    \le\text{virtual length charged}.
\]

After the last chord of~$F(A,P)$ is used, some inactive clients may remain
unprobed after the last active client in a tail list.  One additional complete
backbone scan exhausts all such residual lists; every remaining probe is then
inactive.  At termination the virtual position is~$r$, so the actual return
to~$r$ is a shortcut of the pending walk.  Apply
\cref{lem:active-first-stack-walk} to obtain
\cref{eq:adaptive-virtual-cost}.
\end{proof}

\subsection{The low-mass policy}
\label{sec:low-mass-adaptive-policy}

\begin{lemma}[Zero posterior value]
\label{lem:zero-posterior-policy}
If $\optpost(W)=0$, then the induced subinstance on~$W$ has an adaptive policy
of expected cost zero.
\end{lemma}

\begin{proof}
Let $W^+:=\{v\in W:p_v>0\}$.  The event $A\cap W=W^+$ has positive
probability, so $\opt(W^+)=0$.  Probe the clients of~$W^+$ in the order of a
zero-cost rooted tour through~$W^+$.  Every realized route is a directed
shortcut of that tour and therefore has cost zero.  Clients of probability
zero can be probed last and are certainly inactive.
\end{proof}

\begin{proposition}[Low-mass bound]
\label{prop:adaptive-low-mass}
For every $W\subseteq V^\circ$ with
$\mu=\sum_{i\in W}p_i$,
\[
    \optadapt(W)
    \le C_3\bigl(L+\sqrt{\mu+1}\bigr)\optpost(W),
    \qquad L=\max\{1,\log(n+1)\},
\]
where $C_3$ is universal.
\end{proposition}

\begin{proof}
If $\optpost(W)=0$, use \cref{lem:zero-posterior-policy}.  Otherwise use the
ordinary-backbone certificate, sample the public endpoint pairs~$P$, and
construct~$N(P)$ as above.  Condition on~$P$.  For every cyclic interval~$I$,
\cref{eq:integer-balancing-interval} gives
\[
    \left|\E_A\sum_{e\in F(A,P)}\chi_I(e)\right|\le1.
\]
Apply \cref{lem:adaptive-interval-maximal}\textup{(i)} to the independent
centered variables $X_i:=\1_{\{i\in A\}}-p_i$.  Since
$\sum_i\E[X_i^2]=\sum_i p_i(1-p_i)\le\mu$,
\begin{equation}\label{eq:adaptive-active-imbalance}
    \E_A[\Lambda(F(A,P))]\le1+8\sqrt\mu.
\end{equation}

Take activation expectation in \cref{eq:adaptive-virtual-cost}, and then average
over~$P$.  By \cref{eq:integer-balancing-cost},
\cref{eq:adaptive-sampled-detour-cost}, and
\cref{eq:adaptive-gamma-bound},
\begin{align*}
    &\E_P[\cost(N(P))]
    +\E_P\left[\sum_{i\in U}p_i
      \bigl(d(a_i,i)+d(i,b_i)\bigr)\right]\\
    &\qquad\le
      \cost(z)+(\E_P[\Gamma(P)]+1)\cost(T_R).
\end{align*}
Together with \cref{eq:adaptive-active-imbalance}, this bounds the joint
expected cost of the policy by
\[
    \cost(z)+C(1+\sqrt\mu)\cost(T_R)
    \le C'\bigl(L+\sqrt{\mu+1}\bigr)B,
\]
where the final inequality uses
\cref{lem:ordinary-backbone-certificate}.

The calculation averages over public endpoint preprocessing.  Some finite
outcome~$P$ has conditional expected cost no larger than this joint average;
fix it, together with~$N(P)$, before probing.  Thus the policy may be taken to
be deterministic.  No object used by the policy is selected from the
activation realization.
\end{proof}

\section{Threshold optimization and the clairvoyance-gap bound}
\label{sec:asymmetric-threshold-optimization}

We finish by combining the low-mass policy with a fixed tour for clients of
large activation probability.

\begin{lemma}[Concatenating subinstance policies]
\label{lem:adaptive-policy-concatenation}
Let $W_1,\ldots,W_k$ be disjoint client sets.  If $W_j$ has an adaptive policy
of expected cost at most $\alpha_j\optpost(W_j)$, then executing these
policies successively, each from and back to~$r$, gives an adaptive policy of
expected cost at most
\[
    \left(\sum_{j=1}^k\alpha_j\right)\optpost.
\]
\end{lemma}

\begin{proof}
The costs add.  For every realization and every~$j$,
$A\cap W_j\subseteq A$, so directed shortcutting gives
$\opt(A\cap W_j)\le\opt(A)$.  Take expectations and sum.
\end{proof}

\begin{lemma}[High-probability fixed tour]
\label{lem:high-probability-fixed-tour}
Suppose $p_v\ge\tau$ for every $v\in H$, where $0<\tau\le1$.  Then the
induced subinstance on~$H$ has an adaptive policy of expected cost at most
\[
    \left\lceil\frac1\tau\right\rceil\optpost(H).
\]
\end{lemma}

\begin{proof}
Put $k:=\lceil1/\tau\rceil$.  Independently assign every client of~$H$
uniformly to one of $k$ bins $B_1,\ldots,B_k$.  For every realized partition,
concatenating optimal tours through the bins and shortcutting gives
\[
    \opt(H)\le\sum_{j=1}^k\opt(B_j).
\]
For fixed~$j$, the random set~$B_j$ is a product sample with marginal
$1/k\le p_v$ at client~$v$.  Couple it as a subset of a posterior sample
$A_H$: first sample~$A_H$, then retain each active~$v$ independently with
probability $(1/k)/p_v$.  Directed shortcutting gives
$\E[\opt(B_j)]\le\optpost(H)$.  Averaging over the random partition yields
$\opt(H)\le k\optpost(H)$.  Probe the clients in the order of an optimal tour
through all of~$H$; each realized route is a shortcut of that fixed tour.
\end{proof}

\begin{proof}[Proof of \cref{thm:asymmetric-clairvoyance-upper}]
Set $\tau:=n^{-1/3}$ and split the client set into
\[
    H:=\{v\in V^\circ:p_v\ge\tau\},
    \qquad
    W:=\{v\in V^\circ:p_v<\tau\}.
\]
By \cref{lem:high-probability-fixed-tour}, the high-probability clients have a
    adaptive policy of expected cost at most
\[
    \left\lceil\frac1\tau\right\rceil\optpost(H)
    \le2n^{1/3}\optpost.
\]
For the low-probability set,
\[
    \mu(W)=\sum_{v\in W}p_v\le n\tau=n^{2/3}.
\]
Thus \cref{prop:adaptive-low-mass} supplies an adaptive policy of expected cost at
most
\[
    C\bigl(L+n^{1/3}\bigr)\optpost(W)
    \le C'\bigl(1+n^{1/3}\bigr)\optpost,
\]
where $\log(n+1)=O(1+n^{1/3})$.  Execute the two policies successively and
apply \cref{lem:adaptive-policy-concatenation}.
\end{proof}

The polynomial exponent comes from balancing the two dominant terms.  A
threshold~$\tau$ costs $1/\tau$ for the high-probability fixed tour and
$\sqrt{n\tau}$ for the low-mass policy; equality occurs at
$\tau=n^{-1/3}$.  The deterministic auxiliary-flow rounding is what leaves
only the centered activation process at runtime, so the maximal inequality
incurs no union-bound logarithm.  Improving the exponent within this
framework would therefore require a better treatment of at least one of
these two threshold terms.


% ============================================================
% CHAPTER 6: ALMOST-OPTIMAL A PRIORI TOURS
% ============================================================
\chapter{Almost-Optimal A Priori Tours via Dynamic Programming}\label{ch:dp}

Our main focus in this chapter is to compute almost-optimal a~priori tours for
\StochTSP{}.  A polynomial-time approximation scheme is impossible unless
$\mathrm{P}=\mathrm{NP}$: when every non-root vertex is active with
probability~$1$, the a~priori and a~posteriori optima both coincide with the
optimum of classical metric TSP, which is APX-hard~\cite{Karpinski2015}.

We instead allow exponential running time.  The main result is a deterministic
$(1+\varepsilon)$-approximation algorithm for a~priori TSP on both symmetric
and asymmetric metrics.  For every fixed $\varepsilon>0$, the algorithm runs
in time $c_\varepsilon^n\poly(n)$, where $c_\varepsilon$ is a constant depending
only on $\varepsilon$.

An a~priori solution fixes a tour before the active vertices are known, and
its realization shortcuts the inactive vertices.  The main difficulty is
that vertices far apart in the fixed order may therefore become consecutive.
In \Cref{sec:dp-pairwise-objective}, we
decompose the expected cost into pairwise interactions and show that the
rarest of them can be ignored with only a small loss.  After grouping the
vertices into blocks, the remaining interactions are local enough to be
handled by dynamic programming; this is carried out in
\Cref{sec:dp-block-objective,sec:dp-dynamic-program}.  The parameter choices and
the final approximation bound are given in \Cref{sec:dp-approximation}.

\section{Pairwise decomposition and truncation}
\label{sec:dp-pairwise-objective}
\label{sec:dp-truncation}

We begin by rewriting the expected shortcut cost as a sum over pairs of
vertices in the a~priori order.  A pair contributes precisely when its two
endpoints are active and every vertex lying strictly between them is inactive.  This exact
decomposition identifies the low-probability terms that can later be
discarded.

Let $T=(v_0=r,v_1,\dots,v_n,v_{n+1}=r)$ be an a~priori tour,
where $V^\circ=\{v_1,\dots,v_n\}$.  For $0\le i<j\le n+1$, use the
shorthands $p_i:=p_{v_i}$ and $d_{ij}:=d(v_i,v_j)$, and define
$Q(i,j):=\prod_{i<k<j}(1-p_k)$, with the convention that an empty product
equals~$1$.  In particular, $p_0=p_{n+1}=p_r=1$.
Thus, $Q(i,j)$ is the probability that every vertex strictly between $v_i$ and
$v_j$ is inactive.  The expected cost of $T$ is
\begin{equation}\label{eq:apriori-pairwise}
    \cost(T)
    =
    \sum_{0\le i<j\le n+1}p_i p_j Q(i,j)d_{ij}.
\end{equation}
Indeed, the term indexed by $(i,j)$ is paid exactly when $v_i$ and $v_j$ are
consecutive active vertices in the shortcut of~$T$.

For a threshold parameter $\eta\in(0,1)$, define the truncated cost of an
a~priori tour~$T$ to be the right-hand side of
\cref{eq:apriori-pairwise} without the pairs $(i,j)$ for which
$Q(i,j)\le\eta$.  That is,
\begin{equation}\label{eq:local-pair-cost}
    \cost^{\loc}_{\eta}(T)
    :=
    \sum_{\substack{0\le i<j\le n+1\\Q(i,j)>\eta}}
    p_i p_j Q(i,j)d_{ij}.
\end{equation}

The following lemma shows that a polynomially small threshold suffices to
control the approximation loss between the full and truncated costs.

\begin{lemma}[Pairwise truncation bound]
\label{lem:long-gap-truncation}
Let $\rho\in(0,1]$ be an error parameter and suppose that
\begin{equation}\label{eq:long-gap-truncation}
    0<\eta\le\frac{\rho^2}{16}.
\end{equation}
For any a~priori tour~$T$,
\[
    \cost^{\loc}_{\eta}(T)
    \le \cost(T)
    \le(1+\rho)\cost^{\loc}_{\eta}(T).
\]
\end{lemma}

\begin{proof}
The first inequality follows immediately because the truncated cost omits
only nonnegative terms.  For the second inequality, put
\[
    \beta(\eta):=
    \frac{\eta\bigl(3+\ln(1/\eta)\bigr)}{1-\eta}.
\]
The discarded-pair charging bound proved in
\cref{lem:discarded-pair-charging} gives
\[
    \cost(T)-\cost^{\loc}_{\eta}(T)
    \le\beta(\eta)\cost(T).
\]
The function $\beta$ is increasing on $(0,1)$.  Since
$\eta\le\rho^2/16\le1/16$, we have
\begin{align*}
    \beta(\eta)
    &\le\frac{\rho^2}{15}
        \left(3+\ln16+2\ln\frac1\rho\right)\\
    &\le\frac{\rho^2}{15}\left(4+\frac2\rho\right)
    \le\frac{2\rho}{5}
    \le\frac{\rho}{1+\rho}.
\end{align*}
Here we used $\ln(1/\rho)\le1/\rho-1$, $\ln16<3$, and
$0<\rho\le1$.  Therefore
\[
    \cost^{\loc}_{\eta}(T)
    \ge\bigl(1-\beta(\eta)\bigr)\cost(T)
    \ge\frac{\cost(T)}{1+\rho},
\]
which proves the second inequality.
\end{proof}

\section{The block objective}\label{sec:dp-block-objective}

The truncated objective still depends on the relative order of individual
vertices.  Our goal is to compress most of these order decisions into an
order of blocks, while summarizing the internal order of each block by a
single precomputed value.  The remaining problem is that the exact coefficient of a
pair depends on the inactive probabilities of all vertices between its
endpoints, including vertices in the two endpoint blocks. And enumerating the set
of vertices in between is inefficient.

We solve this problem by placing every vertex of large activation
probability in a singleton block and requiring every remaining block to have
small total activation probability.  Within such a small block, every
inactive product $Q$ is close to~$1$, so we can approximate it by~$1$ and optimize a
simpler internal objective.  For a pair in different blocks, we retain the
inactivity product over the complete blocks strictly between the endpoint
blocks but omit the two partial products inside the endpoint blocks.  These
omissions give an upper bound, and the small block masses control the resulting
overestimate.  The
surrogate obtained from these two approximations is the \emph{block
objective}.

\sharat{A non-root vertex is said to be $\delta$-heavy if $p_v > \delta/2$, and \emph{light} otherwise. Do the same for the definition of admissible blocks and partitions. Call them $\delta$-admissible blocks and partitions.}

We first specify what blocks are allowed.  Fix $0<\delta<1/4$ and recall
that $q_v=1-p_v$.  A non-root vertex is \emph{heavy} if
$p_v>\delta/2$ and \emph{light} otherwise.  For a set of non-root
vertices~$B$, write
\[
    \mu(B):=\sum_{v\in B}p_v,
    \qquad
    q(B):=\prod_{v\in B}q_v.
\]
For the root block $\{r\}$, set $q(\{r\})=1$.

Heavy vertices must remain singletons.  Light vertices may be grouped as long
as their total mass is at most~$\delta$; this upper bound makes the two
approximations above accurate.  We also prevent a run of light blocks from
being split into arbitrarily many blocks of tiny mass.  Except for the first
block in such a run, every light block must have mass at least~$\delta/2$.
This lower bound will later ensure that the dynamic program needs to remember
only a bounded number of recent blocks.

\begin{samepage}
\begin{definition}[Admissible blocks and partitions]\label{def:admissible-partition}
A block is \emph{admissible} if it is either a singleton containing one heavy
vertex or a nonempty set of light vertices with total mass at most~$\delta$.
An ordered partition $(B_1,\dots,B_m)$ is admissible if all its blocks are
admissible and, whenever $B_{t-1}$ and $B_t$ are both light blocks,
\[
    \mu(B_t)\ge\delta/2.
\]
Thus, in every consecutive run of light blocks, only the first block may have
mass below~$\delta/2$.
\end{definition}
\end{samepage}

Fix an admissible ordered partition
$\mathcal B=(B_1,\dots,B_m)$ and introduce the two root blocks
\[
    B_0=B_{m+1}:=\{r\}.
\]
We now formalize the two parts of the block objective.

First consider pairs whose endpoints lie in the same block.
Let $B$ be a light block and let $\tau$ be an order of its vertices.  For
$x<_\tau y$, the inactivity factor contributed by the vertices of~$B$
strictly between $x$ and~$y$ satisfies
\[
    \prod_{x<_\tau z<_\tau y}q_z
    \ge
    1-\sum_{z\in B}p_z
    \ge 1-\delta.
\]
Thus this factor is uniformly close to~$1$, regardless of the precise
separation of $x$ and~$y$ inside the block.  We replace it by~$1$ but still
choose the internal order that gives the smallest resulting contribution.
This gives the definition
\begin{equation}\label{eq:internal-block-cost}
    I(B)
    :=
    \min_{\tau}
    \sum_{x<_\tau y}p_xp_y d(x,y),
\end{equation}
where the minimum is over all orders~$\tau$ of~$B$.  For heavy singletons and
root blocks, set $I(B)=0$.

The value $I(B)$ and a minimizing order can be computed by choosing the last
vertex of the order:
\begin{equation}\label{eq:internal-block-recurrence}
    I(B)
    =
    \min_{y\in B}
    \left(
        I(B\setminus\{y\})
        +\sum_{x\in B\setminus\{y\}}p_xp_y d(x,y)
    \right),
    \qquad I(\emptyset)=0.
\end{equation}
Indeed, the second term is exactly the contribution of the pairs whose last
vertex is~$y$.

Next fix endpoints $x\in B_i$ and $y\in B_j$ in two different blocks, where
$i<j$.  Their exact inactivity factor contains the product of three parts:
the inactive probabilities of the vertices in~$B_i$ after~$x$, the inactive
probabilities of all vertices in the complete blocks strictly between $B_i$
and~$B_j$, and the inactive probabilities of the vertices in~$B_j$ before~$y$.
The first and third parts depend on the internal endpoint positions.  We drop
these two partial-block factors and retain only the product over the complete
blocks strictly between the endpoint blocks.  To aggregate the remaining base
costs between two blocks, define
\[
    C(B_i,B_j)
    :=
    \sum_{x\in B_i}\sum_{y\in B_j}p_xp_y d(x,y).
\]
As in the pairwise truncation from \cref{sec:dp-pairwise-objective}, we keep a block pair only when the retained
inactivity product exceeds~$\eta$.  Combining the internal and cross-block
contributions gives the block objective
\begin{equation}\label{eq:block-objective}
    F_\eta(\mathcal B)
    :=
    \sum_{t=1}^m I(B_t)
    +
    \sum_{\substack{0\le i<j\le m+1\\
                     \prod_{i<t<j}q(B_t)>\eta}}
    \left(\prod_{i<t<j}q(B_t)\right)C(B_i,B_j).
\end{equation}
The first sum summarizes pairs within the same block.  The second sum
summarizes pairs in different blocks using only the product of~$q(B_t)$ over
the complete blocks strictly between the endpoint blocks.

Finally, order every light block according to a minimizer in
\cref{eq:internal-block-cost} and concatenate the block orders, with the root
at the beginning and end.  Denote the resulting a~priori tour by
$T_{\mathcal B}$.  To justify minimizing the block objective, we need two
comparisons.  First, $F_\eta(\mathcal B)$ should be close to the cost of its
corresponding tour $T_{\mathcal B}$; otherwise, a small block value could hide
a tour with large true cost.  Second, every a~priori tour~$T$ should admit an
admissible consecutive partition~$\mathcal B_T$ whose block value is not much
larger than $\cost(T)$.  Applied to an optimal tour, this ensures that the
minimum block value is not much larger than~$\optprior$.  The next lemma makes
both comparisons precise.

\sharat{Mention that $\delta$ is some fixed parameter in the below lemma. Part (i) should say for every $\delta$-admissible ordered partition... Same for the second part.}

\begin{lemma}\label{lem:block-objective-bounds}
The block objective satisfies the following inequalities.
\begin{enumerate}[label=(\roman*)]
    \item For every admissible ordered partition~$\mathcal B$,
    \begin{align}
        \cost^{\loc}_{\eta}(T_{\mathcal B})
        &\le F_\eta(\mathcal B),
        \label{eq:block-objective-dominates-local}\\
        F_\eta(\mathcal B)
        &\le(1-\delta)^{-2}\cost(T_{\mathcal B}).
        \label{eq:block-objective-faithful-upper}
    \end{align}
    Consequently, if $\rho\in(0,1]$ and
    $0<\eta\le\rho^2/16$, then
    \begin{equation}\label{eq:block-objective-two-sided}
        \frac{1}{1+\rho}\cost(T_{\mathcal B})
        \le F_\eta(\mathcal B)
        \le \frac{1}{(1-\delta)^2}\cost(T_{\mathcal B}).
    \end{equation}

    \item For every a~priori tour~$T$, there is an admissible ordered
    partition $\mathcal B_T$, consecutive in~$T$, such that
    \begin{equation}\label{eq:block-objective-from-order}
        F_\eta(\mathcal B_T)
        \le(1-\delta)^{-2}\cost(T).
    \end{equation}
\end{enumerate}
\end{lemma}

\begin{proof}
We first record the estimate that controls every omitted partial-block
factor.  If~$B$ is a light block and $B'\subseteq B$, then
\begin{equation}\label{eq:partial-block-factor}
    q(B')
    \ge 1-\sum_{v\in B'}p_v
    \ge 1-\delta.
\end{equation}
For a heavy singleton or a root block, the corresponding factor equals~$1$.

For the first inequality in~(i), fix an admissible ordered
partition~$\mathcal B$ and consider
endpoints $x\in B_i$ and $y\in B_j$ with $i<j$.  In the tour
$T_{\mathcal B}$, their exact coefficient is
\[
    p_xp_y\,
    q(B_i^{>x})\,
    \left(\prod_{i<t<j}q(B_t)\right)\,
    q(B_j^{<y}),
\]
where $B_i^{>x}$ denotes the vertices of $B_i$ after~$x$, and
$B_j^{<y}$ the vertices of $B_j$ before~$y$.  If $(x,y)$ is retained in
$\cost^{\loc}_{\eta}(T_{\mathcal B})$, then its exact inactivity probability
exceeds~$\eta$.  Removing the two endpoint-block factors can only increase
that probability, so the product over the complete blocks strictly between
the endpoint blocks also exceeds~$\eta$.  Hence the corresponding block pair occurs in
\cref{eq:block-objective}, with a coefficient at least as large as the exact
one.

If both endpoints lie in the same light block~$B$, their exact inactivity
factor is at most~$1$.  Their total retained contribution is therefore at
most the corresponding sum of $p_xp_y d(x,y)$ in the chosen internal order,
which is exactly~$I(B)$.  Summing these bounds over all retained pairs proves
\cref{eq:block-objective-dominates-local}.

For the second inequality in~(i), first consider a cross-block term of
$F_\eta(\mathcal B)$.  Relative to the exact coefficient of the same pair in
$T_{\mathcal B}$, its coefficient omits only the two partial-block factors
$q(B_i^{>x})$ and $q(B_j^{<y})$.  Each is either~$1$ or at least $1-\delta$
by~\cref{eq:partial-block-factor}.  Hence the block coefficient is at most
$(1-\delta)^{-2}$ times the exact coefficient.  Now let~$B$ be a light block
and let $W_B(T_{\mathcal B})$ denote the exact contribution of the pairs
internal to~$B$.  The internal order of $T_{\mathcal B}$ realizes~$I(B)$,
and every inactivity factor between two vertices of~$B$ is at least
$1-\delta$.  Therefore
\[
    I(B)\le \frac{1}{1-\delta}W_B(T_{\mathcal B}).
\]
The cross-block sum in $F_\eta$ contains only a subset of all cross-block
pairs, so summing these comparisons gives
\cref{eq:block-objective-faithful-upper}.
Combining~\cref{eq:block-objective-dominates-local} with
\cref{lem:long-gap-truncation} proves
\cref{eq:block-objective-two-sided}.

For~(ii), start from an arbitrary a~priori tour~$T$ and keep every heavy vertex
as a singleton.  In each maximal consecutive run of light vertices, scan from
right to left.  Close a block as soon as its mass reaches~$\delta/2$, and then
continue with the remaining prefix of the run.  Since every light vertex has
probability at most~$\delta/2$, each closed block has mass in
$[\delta/2,\delta]$.  Only the leftmost block of a run may have mass below
$\delta/2$.  The resulting consecutive partition~$\mathcal B_T$ is therefore
admissible.

Consider a cross-block term of $F_\eta(\mathcal B_T)$ with endpoints
$x\in B_i$ and $y\in B_j$, where $i<j$.  Relative to the exact coefficient of
this pair in~$T$, the corresponding term of $F_\eta(\mathcal B_T)$ omits
only the factors corresponding to the vertices after~$x$ in~$B_i$ and
before~$y$ in~$B_j$.  By~\cref{eq:partial-block-factor}, every cross-block
coefficient in $F_\eta(\mathcal B_T)$ is at most
$(1-\delta)^{-2}$ times the corresponding exact coefficient in~$T$.

It remains to compare the internal terms.  Let $B$ be a light block and retain
the order induced by~$T$ on~$B$.  If $W_B(T)$ denotes the exact contribution
of pairs with both endpoints in~$B$, then the definition of~$I(B)$ in
\cref{eq:internal-block-cost} and the order induced by~$T$ give
\[
    I(B)
    \le
    \sum_{\substack{x,y\in B\\x<_T y}}p_xp_y d(x,y)
    \le
    \frac{1}{1-\delta}W_B(T).
\]
The second inequality holds because every inactivity factor between two
vertices of~$B$ is at least $q(B)\ge1-\delta$.  Summing the cross-block and
internal comparisons proves~\cref{eq:block-objective-from-order}.
\end{proof}

Part~(i) says that the block objective is a two-sided approximation to the
cost of every tour that it represents.  Part~(ii) says that every a~priori
tour, and in particular an optimal one, has an admissible consecutive
partition whose block value is within a factor $(1-\delta)^{-2}$ of its cost.
Together, these comparisons explain why minimizing $F_\eta$ is a valid
surrogate for minimizing the exact tour cost.

\section{The dynamic program}\label{sec:dp-dynamic-program}

We now minimize $F_\eta$ exactly over all admissible ordered partitions.
Consider a partial partition
\[
    (B_0,B_1,\dots,B_t),
    \qquad B_0=\{r\}.
\]
If
\[
    \prod_{s=1}^t q(B_s)>\eta,
\]
its \emph{relevant suffix} is the entire prefix
$(B_0,B_1,\dots,B_t)$.  Otherwise, let $a$ be the largest index such that
\[
    \prod_{s=a}^t q(B_s)\le\eta.
\]
Equivalently,
\[
    \prod_{s=a}^t q(B_s)\le\eta
    \quad\text{and}\quad
    \prod_{s=a+1}^t q(B_s)>\eta.
\]
The relevant suffix is then $(B_a,B_{a+1},\dots,B_t)$.

A state consists of the set $U\subseteq V^\circ$ of already placed vertices and the
relevant suffix~$\Omega$ of their partial partition.  Its value is the minimum,
over all represented prefixes, of the internal terms of the placed blocks plus
all cross-block terms whose right endpoint has already been placed.  The
initial state is $(\emptyset,(B_0))$ with value~$0$.

A feasible next block~$B$ is either a singleton containing an unplaced heavy
vertex or a nonempty set of unplaced light vertices with
$\mu(B)\le\delta$.  If the last block of the current prefix and~$B$ are both
light, we additionally require $\mu(B)\ge\delta/2$.  Suppose
$\Omega=(B_a,\dots,B_t)$.  Appending~$B$ adds
\begin{equation}\label{eq:dp-increment}
    I(B)
    +
    \sum_{s=a}^t
    \left(\prod_{s<u\le t}q(B_u)\right)C(B_s,B).
\end{equation}
Every block in the relevant suffix occurs in this sum, because the complete
blocks between it and~$B$ have inactivity product greater than~$\eta$.
After the transition, the suffix is trimmed according to the same rule.  Once
every non-root vertex has been placed, the terminal root block~$\{r\}$ is appended
using the cross-block part of~\cref{eq:dp-increment} and internal cost zero.

\begin{lemma}[Correctness and running time of the DP]\label{lem:local-dp}
For fixed $\delta$ and~$\eta$, the dynamic program minimizes
$F_\eta$ over all admissible ordered partitions and runs in time
$c_{\delta,\eta}^n\poly(n)$.
\end{lemma}

\begin{proof}
When a block~$B$ is appended, the only new terms of~$F_\eta$ are its internal
term $I(B)$ and the cross-block terms whose right endpoint is~$B$.  These are
exactly the terms in~\cref{eq:dp-increment}.

If a previous block is not in the relevant suffix, then the complete blocks
between it and the current end have inactivity product at most~$\eta$.
Appending more blocks only multiplies this product by numbers in~$[0,1]$.
Hence a forgotten block can never occur in a future term of~$F_\eta$.  The
state therefore contains precisely the information needed to evaluate every
continuation, and the recurrence is an exact shortest-path computation over all
admissible partial partitions.

It remains to bound the number of states.  We use the elementary inequality
\[
    \prod_{v\in S}(1-p_v)
    \le \exp\left(-\sum_{v\in S}p_v\right).
\]
If the whole prefix is retained, then its inactivity product exceeds
$\eta$, so its total activation mass is less than
$\log(1/\eta)$.  Otherwise, deleting the first retained block leaves an
inactivity product greater than~$\eta$, so the remaining blocks have total
activation mass less than $\log(1/\eta)$.  Every admissible block has mass at
most~$1$: a heavy block is a singleton of mass $p_v\le1$, while a light block
has mass at most~$\delta<1$.  Thus the total activation mass of a relevant
suffix is less than
\[
    M_\eta:=\log(1/\eta)+1.
\]

Every retained heavy singleton has mass greater than~$\delta/2$, and every
retained light block except the first one in a light run has mass at least
$\delta/2$.  Let $\mathcal G$ be the collection of these blocks.  Since their
total mass is less than $M_\eta$,
\[
    |\mathcal G|<2\frac{M_\eta}{\delta}.
\]
Every remaining retained block is the first light block of a light run.  The
number of light runs is at most one more than the number of retained heavy
blocks, and hence at most $|\mathcal G|+1$.  Therefore the number of retained
non-root-vertex blocks is at most
\[
    2|\mathcal G|+1
    <4\frac{M_\eta}{\delta}+1.
\]
Including the root block, the relevant suffix has at most
\[
    K:=\left\lceil4\frac{M_\eta}{\delta}\right\rceil+2
    =O\!\left(\frac{\log(1/\eta)+1}{\delta}\right)
\]
blocks.  A state can be encoded by assigning each non-root vertex one of the labels
``unplaced'', ``placed and forgotten'', or ``position $1,\dots,K$ in the
relevant suffix''.  Thus there are at most $(K+2)^n\poly(n)$ states.  A
transition may enumerate a subset of the unplaced light vertices, contributing
at most another factor~$2^n$, and all numerical terms can be evaluated in
polynomial time.  This gives a running time
$c_{\delta,\eta}^n\poly(n)$.  The recurrence
\cref{eq:internal-block-recurrence} for the values $I(B)$ takes
$O(n2^n)$ time and is absorbed in the same bound.
\end{proof}

\section{Approximation guarantee}\label{sec:dp-approximation}

The dynamic program minimizes $F_\eta$ exactly over admissible partitions.
It remains to choose the parameters so that the tour represented by this
minimum has cost at most $(1+\varepsilon)\optprior$.  Since the parameters
depend only on~$\varepsilon$, this choice will keep the running-time base
dependent only on~$\varepsilon$.

\begin{theorem}\label{thm:apriori-ptas}
For every fixed $0<\varepsilon\le1$, a~priori TSP on symmetric or asymmetric
metrics has a deterministic algorithm running in time
$c_\varepsilon^n\poly(n)$ and returning an a~priori tour~$\widehat T$ such that
\[
    \cost(\widehat T)\le(1+\varepsilon)\optprior.
\]
\end{theorem}

\begin{proof}
Set
\[
    \gamma:=\frac{\varepsilon}{10},
    \qquad
    \delta:=\rho:=\gamma,
    \qquad
    \eta:=\frac{\gamma^2}{16}.
\]

Let $T^*$ be an optimal a~priori tour.  By
\cref{lem:block-objective-bounds}(ii), there is an admissible ordered partition
$\mathcal B^*$ into consecutive sets of~$T^*$ such that
\[
    F_\eta(\mathcal B^*)
    \le(1-\gamma)^{-2}\optprior.
\]
Let $\widehat{\mathcal B}$ be a partition minimizing $F_\eta$, as computed by
\cref{lem:local-dp}, and put
$\widehat T:=T_{\widehat{\mathcal B}}$.  Since $0<\gamma\le1/10$,
\cref{lem:block-objective-bounds}(i), the optimality of
$\widehat{\mathcal B}$, and the preceding bound,
\begin{equation}\label{eq:final-approximation-chain}
    \begin{aligned}
        \cost(\widehat T)
        &\le (1+\gamma)F_\eta(\widehat{\mathcal B})\\
        &\le (1+\gamma)F_\eta(\mathcal B^*)\\
        &\le (1+\gamma)(1-\gamma)^{-2}\optprior.
    \end{aligned}
\end{equation}

Using $(1-x)^{-1}\le1+2x$ for $0\le x\le1/2$, we obtain
\[
    (1+\gamma)(1-\gamma)^{-2}
    \le(1+\gamma)(1+2\gamma)^2
    \le1+6\gamma
    \le1+\varepsilon.
\]
Together with~\cref{eq:final-approximation-chain}, this proves the claimed
approximation guarantee.

Finally,
\[
    \log(1/\eta)
    =2\log\frac{4}{\gamma},
\]
which depends only on~$\varepsilon$.  Since $\delta$ also depends only on
$\varepsilon$, \cref{lem:local-dp} shows that
the running time remains $c_\varepsilon^n\poly(n)$.
\end{proof}

\chapter{Conclusion and Future Work}\label{ch:conclusion}

\section{Summary of results}\label{sec:conclusion-summary}

We summarize the main results of this thesis in the following table.

\begin{center}
\small
\begin{tabular}{lll}
\toprule
\textbf{Result} & \textbf{Metric} & \textbf{Reference} \\
\midrule
Adaptivity gap $\ge 2$ & Symmetric & \cref{thm:random-metric-gap} \\
Clairvoyance gap $\ge 10/9$ & Symmetric & \cref{thm:three-chain-cost} \\
Clairvoyance gap $\ge 3/2$ & Asymmetric & \cref{thm:layered-poset-gap} \\
Clairvoyance gap $O(n^{1/3})$ & Asymmetric & \cref{thm:asymmetric-clairvoyance-upper} \\
$(1{+}\varepsilon)$-approx.\ for a~priori TSP & Both & \cref{thm:apriori-ptas} \\
\bottomrule
\end{tabular}
\end{center}

\section{Directions for future research}\label{sec:future-research}

The results of this thesis suggest several possible directions for subsequent
work.

\begin{enumerate}
    \item \textbf{Symmetric clairvoyance gap.}
    What is the correct supremum of the clairvoyance gap for symmetric \StochTSP{}?
    The three-chain construction gives a lower bound
    of~$10/9$.  The upper bound in \cref{thm:general-upper} applies when
    there is exactly one semi-active vertex, but a matching upper bound for
    unrestricted symmetric instances remains open.

    \item \textbf{Symmetric adaptivity gap.}
    Our probabilistic construction gives a lower bound of~$2$, while the
    obliviousness-gap upper bound of~$3$ also bounds the adaptivity gap.
    Is the adaptivity gap equal to~$2$, or can a different instance produce
    a larger separation?

    \item \textbf{Asymmetric clairvoyance gap.}
    Our results place the clairvoyance gap between~$3/2$ and
    $O(n^{1/3})$.  The lower bound comes from the layered-poset construction,
    strengthening the recursive-triangle family sketched in
    \cref{rem:recursive-triangle}, whose ratio tends to~$4/3$;
    the upper bound comes from the adaptive backbone-circulation policy in
    \cref{ch:circulation}.  Is the gap bounded by a constant, or must it grow
    with~$n$?

    \item \textbf{Faster a~priori approximation algorithms.}
    Our $(1{+}\varepsilon)$-approximation for a~priori TSP runs in
    time $c_\varepsilon^n\poly(n)$ and is therefore exponential in~$n$.
    A polynomial-time $(1{+}\varepsilon)$-approximation scheme would also
    apply when every non-root vertex has activation probability~$1$, and would thus
    give a PTAS for the classical TSP.  Since that problem is
    APX-hard, this is impossible unless $\mathrm{P}=\mathrm{NP}$~\cite{Karpinski2015}.
    A more realistic direction is
    to seek a polynomial-time constant-factor approximation for a~priori
    TSP, or more generally to improve the tradeoff between running
    time and approximation ratio.
\end{enumerate}


% ============================================================
% ACKNOWLEDGEMENTS
% ============================================================
\chapter*{Acknowledgements}
\addcontentsline{toc}{chapter}{Acknowledgements}

I would like to thank Dr.\ Sharat Ibrahimpur for his guidance throughout this
project, Manuel Christalla for the continuous discussions and valuable
suggestions, and Prof.\ Dr.\ Vera Traub for advising this thesis.

I would also like to give special thanks to my parents and my girlfriend for
their constant care, understanding, and support in my life.


% ============================================================
% APPENDICES
% ============================================================
\appendix
\chapter{Technical Proof for Pairwise Truncation}
\label{app:pairwise-truncation}

\sharat{This section requires proper verification.}

This appendix proves the charging estimate used in
\cref{lem:long-gap-truncation}.  We use the notation for the a~priori tour,
the probabilities~$p_i$, the distances~$d_{ij}$, and the products~$Q(i,j)$
from \cref{sec:dp-pairwise-objective}.

\begin{lemma}[Discarded-pair charging]
\label{lem:discarded-pair-charging}
For every a~priori tour~$T$ and every $\eta\in(0,1)$,
\[
    \cost(T)-\cost^{\loc}_{\eta}(T)
    \le
    \frac{\eta\bigl(3+\ln(1/\eta)\bigr)}{1-\eta}\cost(T).
\]
\end{lemma}

\begin{proof}
Write
\[
    w_{ij}:=p_ip_jQ(i,j),
    \qquad
    D_\eta:=\cost(T)-\cost^{\loc}_{\eta}(T).
\]
Thus $D_\eta$ is the total contribution of the discarded pairs, namely those
pairs $(i,j)$ for which $Q(i,j)\le\eta$.

Fix a discarded pair $(i,j)$.  Draw the activation states of the vertices
strictly between $v_i$ and~$v_j$ according to their product distribution,
conditioned on at least one of them being active.  The conditioning event has
probability
\[
    1-Q(i,j)\ge1-\eta.
\]
Together with $v_i$ and~$v_j$, the sampled active vertices form a path in
their order in~$T$; denote its edge set by~$P_{ij}$.  The triangle inequality
gives
\[
    d_{ij}\le\sum_{(a,b)\in P_{ij}}d_{ab}.
\]
We charge $w_{ij}d_{ij}$ to the edges of this random path.

Fix a target pair $(a,b)$ with $w_{ab}>0$, and let $\Gamma_{ab}$ be the
expected coefficient charged to~$d_{ab}$.  A source pair $(i,j)$ can charge
$(a,b)$ only if $i\le a<b\le j$.  The source $(a,b)$ itself contributes
nothing because the conditioning ensures that its sampled path contains an
active vertex strictly between its endpoints.  Dividing the contributions of
the other sources by~$w_{ab}$ and separating the three possible strict
containment patterns gives
\begin{equation}\label{eq:appendix-truncation-target-congestion}
\frac{\Gamma_{ab}}{w_{ab}}
\le\frac{1}{1-\eta}
\left(
    \sum_{\substack{j>b\\Q(a,j)\le\eta}}p_jQ(a,j)
    +\sum_{\substack{i<a\\Q(i,b)\le\eta}}p_iQ(i,b)
    +\sum_{\substack{i<a,\,j>b\\Q(i,j)\le\eta}}p_ip_jQ(i,j)
\right).
\end{equation}
The first two sums are each at most~$\eta$.  Indeed, $Q(a,j)$ is
nonincreasing in~$j$, so the indices satisfying $Q(a,j)\le\eta$ form a
suffix, and telescoping gives
\[
    \sum_{\substack{j>b\\Q(a,j)\le\eta}}p_jQ(a,j)\le\eta.
\]
The second sum is bounded symmetrically.

It remains to bound the third sum in
\cref{eq:appendix-truncation-target-congestion}.  Put
\[
    H:=(1-p_a)Q(a,b)(1-p_b),
    \qquad
    x_i:=Q(i,a),
    \qquad
    y_j:=Q(b,j).
\]
For $i<a<b<j$, we have $Q(i,j)=Hx_iy_j$.  If either outside index set is
empty, the third sum vanishes.  Otherwise, telescoping gives
\[
    \sum_{i<a}p_ix_i=1,
    \qquad
    \sum_{j>b}p_jy_j=1.
\]
If $H\le\eta$, the third sum is therefore at most $H\le\eta$.  Suppose
instead that $H>\eta$, and set $t:=\eta/H\in(0,1)$.  For every $z>0$,
telescoping on the right gives
\[
    \sum_{\substack{j>b\\y_j\le z}}p_jy_j\le\min\{1,z\}.
\]
Consequently,
\begin{align*}
    \sum_{\substack{i<a,\,j>b\\Q(i,j)\le\eta}}p_ip_jQ(i,j)
    &=H\sum_{\substack{i<a\\x_i>0}}p_ix_i
        \sum_{\substack{j>b\\y_j\le t/x_i}}p_jy_j\\
    &\le H\sum_{i<a}p_i\min\{x_i,t\}.
\end{align*}

We use the following elementary survival-product estimate:
\begin{equation}\label{eq:appendix-survival-product-estimate}
    \sum_{i<a}p_i\min\{Q(i,a),t\}
    \le t\left(1+\ln\frac1t\right)
    \qquad(0<t<1).
\end{equation}
To verify it, recall that $x_i=Q(i,a)$ is nondecreasing in~$i$ and
$x_{i-1}=(1-p_i)x_i$.  If $x_0>t$, then
\[
    \sum_{i<a}p_i\min\{x_i,t\}
    =t\sum_{i<a}p_i
    \le t\left(1+\ln\frac1{x_0}\right)
    \le t\left(1+\ln\frac1t\right).
\]
Otherwise, let~$\ell$ be the largest index for which $x_\ell\le t$, and put
$x:=x_{\ell+1}$ and $z:=x_\ell$.  Since
$z=(1-p_{\ell+1})x$, telescoping and the inequality
$u\le-\ln(1-u)$ give
\begin{align*}
    \sum_{i<a}p_i\min\{x_i,t\}
    &\le z+t\left(1-\frac{z}{x}+\ln\frac1x\right)\\
    &\le t\left(2-\frac{t}{x}+\ln\frac1x\right)\\
    &\le t\left(1+\ln\frac1t\right).
\end{align*}
The last inequality is equivalent to
$1-t/x+\ln(t/x)\le0$, which follows from $\ln u\le u-1$.  This proves
\cref{eq:appendix-survival-product-estimate}.

Applying the estimate, the third sum in
\cref{eq:appendix-truncation-target-congestion} is at most
\[
    Ht\left(1+\ln\frac1t\right)
    =\eta\left(1+\ln\frac{H}{\eta}\right)
    \le\eta\left(1+\ln\frac1\eta\right).
\]
It follows that every target pair has congestion at most
\[
    \beta(\eta):=\frac{\eta\bigl(3+\ln(1/\eta)\bigr)}{1-\eta}.
\]
Pairs with $w_{ab}=0$ receive no charge, so summing over all target pairs
gives
\[
    D_\eta\le\beta(\eta)\cost(T),
\]
which is the claimed inequality.
\end{proof}

\chapter{Declaration of Originality}\label{app:declaration}

I hereby declare that I authored this thesis independently.  Suggestions from
the supervisor regarding language and content are excepted.  All sources and
aids used in preparing this thesis, including any authorized generative
artificial intelligence technologies, are documented truthfully and fully in
consultation with the supervisor.  I have adhered to the applicable citation
guidelines, and all passages taken verbatim or paraphrased from published or
unpublished works of others are indicated as such.

\bigskip

\noindent
\begin{tabular}{@{}l p{0.8\textwidth}@{}}
Title: & \thesistitle \\
Name: & \thesisauthor \\
Date: & \thesisdate \\
Signature: & \rule{5cm}{0.4pt}
\end{tabular}


% ============================================================
% BIBLIOGRAPHY
% ============================================================
\printbibliography[heading=bibintoc]

\end{document}
